\chapter{Functoriality and change of field}
\index{functoriality!BFJ}\index{base change!non-Archimedean}
\label{chap:bfj-functoriality}

Functoriality in non-Archimedean pluripotential theory is strongest on the
model locus.  There, pullback is pullback of line bundles and divisors,
while pushforward is governed by the projection formula.  Passing to
continuous metrics preserves many of these identities.  Passing to
singular psh functions or currents requires more care: pullback of a general
current is not defined, a birational pullback of an ample class is usually
only nef and big, and a Radon measure has no canonical pullback under an
arbitrary map.

This chapter tracks the normalization of pullback, pushforward and field
extension, from model intersections to the analytic statements their
hypotheses permit.  The reader should learn to distinguish a degree factor
from scalar-extension invariance, to check ampleness before using the
singular operator upstairs, and to separate topological descent from
positivity descent.  The worked curve and surface examples test these
interfaces numerically.
Throughout, morphisms of varieties induce continuous maps of Berkovich
spaces, denoted by the same letter with a superscript $\An$ only when
ambiguity is possible.

\section{Orientation}
\label{sec:ch30-orientation}

\subsection{Three ways of changing the base, and why they are different}
\label{sec:ch30-three-kinds}

The word \emph{functoriality} covers three operations in this book, and they
behave differently enough that conflating them produces false statements.

\emph{(A) A morphism over a fixed ground field.}  A morphism $f:Y\to X$ of
smooth projective $K$-varieties induces $f^{\An}:Y^{\An}\to X^{\An}$
(\Cref{sec:berk-functoriality}).  Here the ground valuation, its value group,
the residue field, and the meaning of a model are unchanged; the
variety moves.  The two cases that matter are $f$ finite surjective and
$f$ birational.  Their projection formulae differ by exactly one number,
the degree: a finite morphism of degree $d$ multiplies every top
intersection number, hence every Monge--Amp\`ere mass and every energy, by
$d$, while a birational morphism contributes the factor $1$.  Their
analytic scopes nevertheless differ, because a nontrivial birational
pullback destroys ampleness of the reference class.

\emph{(B) An extension of the ground field.}  A complete valued extension
$K'/K$ induces $\pi:X_{K'}^{\An}\to X^{\An}$
(\Cref{sec:berk-functoriality}).  Here the variety is unchanged --- its
function field merely acquires new scalars --- while the ground valuation,
its value group, and the residue field all move.  Top intersection numbers
are \emph{invariant} under scalar extension, so masses and energies are
unchanged; the analytic spaces, models and available potentials can change.  This
does not assert strict growth for every extension or every variety.  The bookkeeping that replaces the degree
$d$ of case (A) is a pair of numbers attached to each component of the
special fibre: a ramification index and a residue degree, whose product is
constrained by $[K':K]$ and which cancel against each other in the measure
formula.  This cancellation is the content of
\Cref{thm:base-change-ma}, and it is the reason no factor $[K':K]$ appears
there.

\emph{(C) The Gauss section of a special field extension.}  If $k$ is
trivially valued and $K=k(\!(t)\!)$, the inclusion $k\subset K$ is a
complete valued extension, so case (B) gives the usual base-change
projection
\[
                   r:X_K^{\An}\longrightarrow X^{\An}.
\]
What is special is an additional canonical continuous \emph{section}
$\sigma:X^{\An}\to X_K^{\An}$, built from Gauss valuations.  This section
is not induced by a morphism $X\to X_K$ and has no analogue for a general
valued extension; it is not surjective, and its image is a very thin subset
of $X_K^{\An}$.  What transfers along $\sigma$, as opposed to what follows
from the ordinary base-change projection $r$, is the subject of
\Cref{sec:ch30-gauss}; it is the item on which \Cref{part:bj} depends.

The three operations are compared once more, side by side, in the ledger of
\Cref{sec:ch30-ledger}.

\subsection{Four operations, and which of them exist}
\label{sec:ch30-four-operations}

Against each of (A), (B), (C) one may try to transport four kinds of object.
The following table is the summary of the chapter; each entry is proved or
disproved below.

\begin{center}
\small
\begin{tabular}{p{.20\textwidth}|p{.24\textwidth}|p{.22\textwidth}|p{.24\textwidth}}
\toprule
&\textbf{(A) morphism}&\textbf{(B) field extension}&\textbf{(C) Gauss
 section}\\
\midrule
numerical classes, pull back
 &yes, \eqref{eq:ch21-pullback-form}
 &yes, and $\NS$ may grow
 &yes through the underlying base change,
  $\theta\mapsto\theta_K$\\
potentials, pull back
 &model and continuous potentials always; singular potentials when the
  pulled-back class is ample and the polar condition of
  \Cref{thm:ch30-pullback-psh} holds
 &yes, \Cref{thm:ch30-base-change-psh}
 &yes, $\sigma^*$ on PL and model data\\
measures, push forward
 &yes, \eqref{eq:measure-pushforward}
 &yes, \Cref{thm:base-change-ma}
 &yes for matched model measures, along $\sigma$ and $r$\\
measures, lift
 &requires additional structure, such as finite local degrees
 &a unique invariant lift for a finite Galois extension
 &yes, by $\sigma_*$; $r_*$ inverts it on measures carried by the image\\
\bottomrule
\end{tabular}
\end{center}

Composition of continuous test functions defines pushforward of Radon
measures; an arbitrary continuous map does not specify a preferred
inverse.  \Cref{prop:ch30-pushforward-not-inverse} illustrates that
nonuniqueness.  Additional structure can select a lift: a finite Galois
quotient has the unique invariant lift proved in
\Cref{cor:ch30-invariant-functions}, and the Gauss section gives
$\sigma_*\mu$ for every measure downstairs.
The Monge--Amp\`ere identities below are projection formulae; these
special lifts must be distinguished from an inverse to unrestricted
pushforward.

\subsection{Standing hypotheses and notation}
\label{sec:ch30-hypotheses}

Except in \Cref{sec:ch30-gauss}, the original BFJ signature
\eqref{eq:ch20-standing} is in force on the base: $K$ is complete and
discretely valued with residue field $k$ of characteristic zero, $X$ is a
smooth connected projective $K$-variety of dimension $n$, models are
normal flat projective $R$-schemes with generic fibre $X$, and
$\theta$ is a closed model $(1,1)$-form whose class is ample.  We choose
a semipositive reference representative when using zero potentials and
constant truncations.  This is possible by
\Cref{prop:ch20-positive-representative}: replace $\theta$ by
$\theta+\ddc h$ and each potential $\varphi$ by $\varphi-h$.
The total curvature is unchanged, and
\Cref{prop:ch26-reference-change} records the energy adjustment.
The model projection identities apply to the original closed reference as
well, by real multilinearity and this reference change.  Write
\[
                      V=(\{\theta\}^n)>0 .
\]

Two departures from the global notation are made, and are used without
further comment.

\emph{The uniformizer is $\varpi$.}  The letter $\pi$ is reserved
throughout this chapter for the projection $X_{K'}^{\An}\to X^{\An}$
attached to a field extension.  We therefore write $\varpi$ for a
uniformizer of $R=K^\circ$, so that \eqref{eq:ch20-standing} reads
$v_K(\varpi)=1$, and $\varpi'$ for a uniformizer of $R'=K'^\circ$.
This is the convention already used in
\Cref{chap:trivial-analytification,chap:bj-test-configurations}.  The one
exception is \Cref{sec:ch30-gauss}, where the ground field is
$K=k(\!(t)\!)$ and its uniformizer is called $t$, in agreement with
\eqref{eq:gauss-extension}; its base-change projection is denoted by $r$, so no clash arises.

\emph{Geometric integrality.}  From \Cref{sec:ch30-base-change} on we assume
in addition that $X$ is geometrically integral over $K$, that is, that
$K$ is algebraically closed in $K(X)$.  This guarantees that $X_{K'}$
is again integral for every field extension $K'/K$, so that
\eqref{eq:ch20-standing} continues to hold upstairs and the degree
computations of \Cref{sec:ch30-base-change} are not obscured by a sum over
connected components.  \Cref{rem:ch30-not-geometrically-integral} records
what changes without this hypothesis.  Note that $\operatorname{char}K=0$,
since $\operatorname{char}k=0$, so all field extensions occurring below are
separable and no inseparability phenomena arise; this is used silently in
several ramification computations and is flagged again in
\Cref{sec:ch30-scope}.

For a morphism $f:Y\to X$ of $K$-varieties we write $f^*\varphi$ for
$\varphi\circ f^{\An}$ when $\varphi$ is a function on $X^{\An}$, and
$f_*\nu$ for the image measure of a Radon measure $\nu$ on $Y^{\An}$.
The two notations are adjoint in the elementary sense
\eqref{eq:measure-pushforward} and in no other.

\subsection{What is proved here and what is imported}
\label{sec:ch30-imported}

The model calculations reduce to intersection theory; their analytic
extensions use the preceding psh and operator results.  The quotient and
Gauss-section arguments also use the stated elementary valuation and
descent inputs.  The main dependencies are the following.

\begin{enumerate}
\item\label{itm:ch30-import:projection}
The algebraic projection formula for vertical intersection numbers,
\Cref{thm:model-projection-formula}, together with its two companions
\Cref{lem:vertical-intersection-degree} and \Cref{prop:nef-pullback}.  This
is what makes every statement in
\Cref{sec:ch30-finite,sec:ch30-birational} true, and it is the reason the
degree $d$ appears exactly once in each formula.

\item\label{itm:ch30-import:pullback}
Pullback of model and continuous psh functions,
\Cref{prop:ch21-pullback}, and the stability calculus
\Cref{thm:ch21-stability}, in particular
\Cref{thm:ch21-stability}\,\cref{itm:ch21-stab:max} and
\Cref{thm:ch21-stability}\,\cref{itm:ch21-stab:unif}.

\item\label{itm:ch30-import:regularization}
Decreasing regularization, \Cref{thm:bfj-decreasing-regularization}, and
stability of $\PSH(X,\theta)$ under decreasing nets,
\Cref{thm:ch21-decreasing-stability}.  These two together upgrade every
statement about model functions to a statement about arbitrary psh
functions whenever the source and target both satisfy the ample-class BFJ
signature.  Their availability includes the underlying model and resolution
hypotheses.  Residue characteristic zero also supplies the resolution
used in the model constructions and the defectless valued-field setting
of \eqref{eq:ch30-efn}; it is not merely a hypothesis of the last
approximation step.

\item\label{itm:ch30-import:ma}
The bounded Monge--Amp\`ere operator \Cref{thm:bounded-ma-extension},
with the weighted monotone convergence property
\eqref{eq:weighted-monotone-ma}, and its locality
\Cref{cor:ch24-model-locality-open}.

\item\label{itm:ch30-import:energy}
The energy calculus of \Cref{chap:bfj-energy}: the definition
\eqref{eq:model-energy}, monotone continuity
\Cref{prop:e1-monotone-continuity}, the truncation formula
\Cref{prop:ch26-truncation-limit}, and the mixed products
\Cref{thm:e1-mixed-products}.

\item\label{itm:ch30-import:analytic}
From \Cref{chap:berkovich}: functoriality and base change of
analytifications (\Cref{sec:berk-functoriality}), the description of a point
as a pair (\Cref{prop:berk-pair-description}), properness
(\Cref{thm:berk-proper-compact}), and the reduction map
(\Cref{sec:berk-reduction}).

\item\label{itm:ch30-import:descent}
Faithfully flat descent, in the form of effectivity for projective schemes
with a semilinear action of a finite group, \cite[Exp.~VIII]{SGA1}; this
occurs once, in \Cref{thm:ch30-model-descent}, and is the only external
input of \Cref{sec:ch30-descent}.
\end{enumerate}

Two things are \emph{not} imported, and their absence shapes the chapter.
There is no theory of psh functions in a nef and big class in
\Cref{part:bfj}; this is why \Cref{sec:ch30-birational} proves projection
formulae for pulled-back potentials and nothing else.  And there is no
pullback of currents in the sense of \Cref{part:cld}; this is why every
statement below is about functions and measures.

\subsection{Plan}
\label{sec:ch30-plan}

\Cref{sec:ch30-pullback} treats pullback along an arbitrary morphism, and
separates two requirements for pulling back a singular psh function: the
pulled-back class must remain in the ample BFJ framework, and the image must
not be contained in the polar set.
\Cref{sec:ch30-finite} specializes to finite morphisms and proves the
projection formula for mixed Monge--Amp\`ere measures and for energy, first
for model potentials, then for bounded ones, then on $\cE^1$; the degree
appears as a multiplicative factor throughout, and
\Cref{ex:ch30-power-map} carries the bookkeeping to numbers on
$\mathbf P^1$.  \Cref{sec:ch30-birational} does the same for birational
morphisms, where the factor is $1$, and where the essential structural fact
is that a birational morphism is a bijection on the locus of points of zero
support (\Cref{lem:ch30-birational-bijection}); \Cref{ex:ch30-surface-blowup}
is a surface computation.  \Cref{sec:ch30-base-change} treats extension of
the ground field, with the ramification bookkeeping displayed in full, and
\Cref{ex:ch30-conic-base-change} tracks a skeleton, a measure and an energy
through a ramified quadratic extension.  \Cref{sec:ch30-descent} proves what
descends and states precisely what is not known to.
\Cref{sec:ch30-pushforward} isolates the general projection formula and the
elementary reason pushforward is not invertible.  \Cref{sec:ch30-gauss}
constructs the Gauss extension and proves the four transfer statements ---
centers, supports, classes, total mass --- required by the comparison
programme, with \Cref{ex:ch30-gauss-normal-cone} and
\Cref{ex:ch30-gauss-two-rays} carrying them to numbers.
\Cref{sec:ch30-scope} lists where each hypothesis is used,
\Cref{sec:ch30-ledger} is the dictionary, and the source notes close the
chapter.

\section{Pullback along a morphism}
\label{sec:ch30-pullback}

\subsection{The pulled-back class}
\label{sec:ch30-pullback-class}

Let $f:Y\to X$ be a morphism of smooth connected projective
$K$-varieties, and let $\theta\in\mathcal Z^{1,1}(X)$ be a closed model
form determined on a model $\cX$.  The construction of the pulled-back
form was made in \Cref{sec:ch21-pullback}: choose any model $\cY_0$ of
$Y$, let $\cY$ be a normalization of the closure of the graph of $f$
inside $\cY_0\times_R\cX$ followed by a resolution, so that there is an
$R$-morphism $\tilde f:\cY\to\cX$ extending $f$, and set
$f^*\theta:=[\tilde f^*\theta_{\cX}]$ as in
\eqref{eq:ch21-pullback-form}.  Two such extensions agree on a model of
$Y$ dominating both, so the class is well defined.

Three features of this construction deserve to be stated before any analysis
begins, because each of them is a potential source of error.

\emph{The pulled-back class need not be ample.}  Its generic-fibre class is $f^*\{\theta\}$; it is semipositive when
$\theta$ is, by \Cref{prop:nef-pullback}.
If $f$ is finite, $f^*\{\theta\}$ is again ample; if $f$ contracts a
curve, $f^*\{\theta\}$ has degree zero on it and is not ample.
The pullback of an ample class is nef; it is also big for a dominant
generically finite map of equal-dimensional varieties, including the
birational case.
The standing hypothesis \eqref{eq:ch20-standing} therefore fails on $Y$
for a birational $f$, and no theorem of
\Crefrange{chap:bfj-psh}{chap:bfj-energy} may be applied there without
comment.  This is the single most important boundary in the chapter, and it
is discussed again in \Cref{sec:ch30-birational}.

\emph{The dimension may drop.}  If $\dim Y<\dim X$, then
$(\{f^*\theta\}^{\dim Y})$ is an intersection number on $Y$, unrelated to
$V$, and there is no projection formula for top-degree measures.  All
measure-theoretic statements below assume $\dim Y=\dim X$; the
function-theoretic statement \Cref{thm:ch30-pullback-psh} does not, but it
does require $f^*\{\theta\}$ to be ample.

\emph{Analytification is not a bijection.}  The map $f^{\An}$ is
surjective when $f$ is (properness, \Cref{thm:berk-proper-compact}), but
it need not be injective, including for a nontrivial birational $f$; see
\Cref{lem:ch30-birational-bijection} for the exact statement of what is
preserved.

\subsection{Model functions pull back to model functions}
\label{sec:ch30-pullback-model}

The elementary layer of the theory is \Cref{prop:ch21-pullback}, proved in
\Cref{chap:bfj-psh}, which we restate in the form used below and supplement
with two identities that are needed in the measure computations.

\begin{lemma}
\label{lem:ch30-pullback-model-calculus}
Let $f:Y\to X$ be a morphism of smooth connected projective
$K$-varieties and let $\tilde f:\cY\to\cX$ be an extension to models as
above.
\begin{enumerate}
\item\label{itm:ch30-pbm:divisor}
If $D$ is a vertical $\mathbf R$-Cartier divisor on $\cX$, then
$\tilde f^*D$ is a vertical $\mathbf R$-Cartier divisor on $\cY$ and
\begin{equation}
\label{eq:ch30-pullback-model-function}
              f^*\varphi_D=\varphi_D\circ f^{\An}=\varphi_{\tilde f^*D}.
\end{equation}
\item\label{itm:ch30-pbm:linear}
The map $f^*$ is linear on real model functions and preserves constants.
As identities of functions it satisfies
$f^*\max\{\varphi,\psi\}=\max\{f^*\varphi,f^*\psi\}$ and
$f^*(\varphi+c)=f^*\varphi+c$ for $c\in\mathbf R$.
A maximum remains a model function when both inputs are rational model
functions; this is not asserted for arbitrary real model functions.
\item\label{itm:ch30-pbm:ddc}
On model functions, $\ddc$ commutes with pullback:
\begin{equation}
\label{eq:ch30-pullback-ddc}
     f^*\theta+\ddc f^*\varphi_D=\bigl[\tilde f^*(\theta_{\cX}+D)\bigr].
\end{equation}
In particular, if $\varphi_D$ is $\theta$-psh, then $f^*\varphi_D$ is
$f^*\theta$-psh.
\item\label{itm:ch30-pbm:sup}
$\sup_{Y^{\An}}f^*\varphi\leq\sup_{X^{\An}}\varphi$, with equality if
$f$ is surjective.
\end{enumerate}
\end{lemma}

\begin{proof}
Assertions \cref{itm:ch30-pbm:divisor}, \cref{itm:ch30-pbm:ddc} and
\cref{itm:ch30-pbm:sup} are
\Cref{prop:ch21-pullback}\,\cref{itm:ch21-pb:model} and
\cref{itm:ch21-pb:sup}, together with the sign convention
\eqref{eq:ch20-sign}: the local equation of $\tilde f^*D$ is the pullback
of a local equation of $D$, so
$\varphi_{\tilde f^*D}(y)=v_y(\tilde f^\#f_D)=v_{f^{\An}(y)}(f_D)
=\varphi_D(f^{\An}(y))$.  Verticality holds because
$\tilde f^{-1}(\cX_s)=\cY_s$, both being the fibres over the closed point
of $\Spec R$.

For \cref{itm:ch30-pbm:linear}, all four identities are pointwise identities
between compositions of functions with the fixed map $f^{\An}$, and
composition on the right commutes with linear combinations, maxima and the
addition of constants.  Linear combinations and constants are represented
by the corresponding real vertical divisors on a common model.
For rational inputs, \Cref{prop:ch20-model-max} represents the maximum
on a dominating model.  For real inputs the pointwise maximum identity
still holds, but \Cref{ex:ch20-real-max-irrational-breakpoint} shows why
its result need not be a real model function.
\end{proof}

The reader should notice which item of
\Cref{lem:ch30-pullback-model-calculus} carries content.  Items
\cref{itm:ch30-pbm:divisor} and \cref{itm:ch30-pbm:linear} are bookkeeping
with local equations.  Item \cref{itm:ch30-pbm:ddc} is where positivity
enters, and it rests on exactly one geometric fact, namely that the pullback
of a relatively nef class along a morphism of models is relatively nef
(\Cref{prop:nef-pullback}), which in turn is the projection formula for a
curve in a special fibre.  Everything else in this chapter is downstream of
that one line.

\subsection{Pullback of a singular potential}
\label{sec:ch30-pullback-singular}

\Cref{chap:bfj-psh} deliberately stopped at continuous potentials and
deferred the general case (\Cref{prop:ch21-pullback} and the warning
following it).  The obstruction identified there was twofold: the
approximation argument used for continuous potentials is a uniform one and
is unavailable in general, and $f^*\varphi$ may be identically $-\infty$,
in which case it is not a psh function at all by
\Cref{def:ch21-bfj-psh}.  Decreasing regularization removes the first
obstruction completely.  Once the pulled-back class remains in the ample
BFJ framework, avoiding the identically $-\infty$ limit is the remaining
analytic condition.

\begin{theorem}
\label{thm:ch30-pullback-psh}
\index{pullback!of psh functions}
Assume \eqref{eq:ch20-standing} and let $f:Y\to X$ be a morphism of smooth
connected projective $K$-varieties such that the generic-fibre class
$f^*\{\theta\}$ is ample on $Y$.  Let
$\varphi\in\PSH(X,\theta)$ and suppose that
\begin{equation}
\label{eq:ch30-pullback-hypothesis}
 f^{\An}(Y^{\An})\not\subset P(\varphi):=\{\varphi=-\infty\}.
\end{equation}
Then $f^*\varphi=\varphi\circ f^{\An}\in\PSH(Y,f^*\theta)$.  Moreover:
\begin{enumerate}
\item\label{itm:ch30-pbpsh:surj}
Hypothesis \eqref{eq:ch30-pullback-hypothesis} holds automatically if
$f$ is surjective, and then
$\sup_{Y^{\An}}f^*\varphi=\sup_{X^{\An}}\varphi$.
\item\label{itm:ch30-pbpsh:order}
On the psh functions satisfying
\eqref{eq:ch30-pullback-hypothesis}, $f^*$ is order preserving and
commutes with finite maxima, addition of constants and convex combinations.
\item\label{itm:ch30-pbpsh:decreasing}
If $\varphi_i\searrow\varphi$ is a decreasing net in $\PSH(X,\theta)$
and \eqref{eq:ch30-pullback-hypothesis} holds for $\varphi$, then
$f^*\varphi_i\searrow f^*\varphi$ in $\PSH(Y,f^*\theta)$.
\item\label{itm:ch30-pbpsh:polar}
$P(f^*\varphi)=(f^{\An})^{-1}\bigl(P(\varphi)\bigr)$, and if $f$ is
surjective then $\varphi\preceq\psi$ implies
$f^*\varphi\preceq f^*\psi$ for the singularity order of
\Cref{def:ch21-polar-singularity}.
\end{enumerate}
\end{theorem}

\begin{proof}
\emph{Step 1: regularization upstairs.}  By
\Cref{thm:bfj-decreasing-regularization} there is a decreasing net
$(\varphi_i)$ of $\theta$-psh model functions with
$\varphi_i\searrow\varphi$ pointwise on $X^{\An}$.  By
\Cref{lem:ch30-pullback-model-calculus}\,\cref{itm:ch30-pbm:ddc} each
$f^*\varphi_i$ is an $f^*\theta$-psh model function on $Y$, and the
net $(f^*\varphi_i)$ is decreasing because composition with a fixed map
preserves pointwise inequalities.  Its pointwise infimum is
\[
 \inf_i\,(\varphi_i\circ f^{\An})
 =\bigl(\inf_i\varphi_i\bigr)\circ f^{\An}
 =f^*\varphi,
\]
the middle equality because the infimum of a net of extended-real numbers is
computed pointwise and $f^{\An}$ is a fixed map.

\emph{Step 2: the limit is not identically $-\infty$.}  Hypothesis
\eqref{eq:ch30-pullback-hypothesis} provides $y\in Y^{\An}$ with
$\varphi(f^{\An}(y))>-\infty$, that is, $f^*\varphi(y)>-\infty$.

\emph{Step 3: conclusion.}  The ampleness of $f^*\{\theta\}$ places
$(Y,f^*\theta)$ in the standing framework of \Cref{chap:bfj-psh}.
By \Cref{thm:ch21-decreasing-stability} applied on $Y$, the decreasing net
$(f^*\varphi_i)$ has infimum either identically $-\infty$ or in
$\PSH(Y,f^*\theta)$; Step~2 excludes the first alternative.  For a finite
morphism the required ampleness is automatic.

\cref{itm:ch30-pbpsh:surj}  If $f$ is surjective then so is $f^{\An}$, by
properness, so $f^{\An}(Y^{\An})=X^{\An}\not\subset P(\varphi)$ because
$\varphi\not\equiv-\infty$.  The equality of suprema is
\Cref{lem:ch30-pullback-model-calculus}\,\cref{itm:ch30-pbm:sup}, which is a
statement about the range of a function and does not use psh-ness.

\cref{itm:ch30-pbpsh:order}  Each identity holds pointwise before any
positivity is discussed, exactly as in
\Cref{lem:ch30-pullback-model-calculus}\,\cref{itm:ch30-pbm:linear}.  The
pullbacks of the input functions are psh by the first part of the theorem;
their maxima and convex combinations are psh by
\Cref{thm:ch21-stability}, now applied in the ample class
$f^*\{\theta\}$.  Addition of constants and preservation of order are
immediate.

\cref{itm:ch30-pbpsh:decreasing}  The net $(f^*\varphi_i)$ is decreasing
with pointwise infimum $f^*\varphi$ by the computation of Step~1, which
did not use that the $\varphi_i$ were model functions.  Each term is psh
by the first part of the theorem, whose hypothesis holds for
$\varphi_i\geq\varphi$ once it holds for $\varphi$.

\cref{itm:ch30-pbpsh:polar}  The description of the polar set is immediate
from $f^*\varphi(y)=\varphi(f^{\An}(y))$.  For the singularity order,
recall from \Cref{def:ch21-polar-singularity} that $\varphi\preceq\psi$
means $\varphi\leq\psi+C$ for some constant $C$; composing with
$f^{\An}$ preserves this, and surjectivity guarantees that the pulled-back
functions are again not identically $-\infty$.
\end{proof}

Even when the pulled-back class is ample, hypothesis
\eqref{eq:ch30-pullback-hypothesis} is not a technical convenience.  It
fails, and the conclusion fails with it, in the simplest situation in
which a subvariety is contained in a polar set.

\begin{example}[A pullback that is identically $-\infty$]
\label{ex:ch30-polar-pullback-fails}
Let $X=\mathbf P^2_K$ with homogeneous coordinates $z_0,z_1,z_2$, let
$L=\cO(1)$ with the standard model metric coming from
$\mathbf P^2_R$, and let $\theta$ be its curvature form, so that
$V=(\{\theta\}^2)=1$.  The section $z_1$ has a well-defined weight
relative to the reference metric,
\[
        \varphi(x):=\log\lVert z_1(x)\rVert_0
        =\log\frac{|z_1(x)|}
        {\max\{|z_0(x)|,|z_1(x)|,|z_2(x)|\}},
\]
and $\varphi\in\PSH(X,\theta)$ with
$P(\varphi)=\{z_1=0\}^{\An}=H^{\An}$, where $H\subset\mathbf P^2_K$ is
the line $z_1=0$.  Indeed $\varphi$ is a decreasing limit of the
$\theta$-psh model functions
$\varphi_m:=\max\{\varphi,-m\}$, and its polar set is exactly the locus
where the section vanishes; this is the two-dimensional analogue of
\Cref{ex:ch26-infinite-energy}; the truncations are the model weights of
the ideals $(z_1,\varpi^m z_0,\varpi^m z_1,\varpi^m z_2)$ on the
standard homogeneous charts.

Take $Y=H\simeq\mathbf P^1_K$ and let $f:Y\to X$ be the inclusion.  Then
$f^{\An}(Y^{\An})=H^{\An}=P(\varphi)$, hypothesis
\eqref{eq:ch30-pullback-hypothesis} fails, and
$f^*\varphi\equiv-\infty$.  This is not a psh function in the sense of
\Cref{def:ch21-bfj-psh}, and it cannot be repaired: the class
$\PSH(Y,f^*\theta)$ is nonempty (it contains the constants, since
$f^*\theta$ is semipositive with ample class $\cO_{\mathbf P^1}(1)$),
so the conclusion of \Cref{thm:ch30-pullback-psh} is genuinely false here
and not merely unproved.

The example also shows that no analogue of the complex \emph{restriction}
theorem --- pull back, subtract the generic Lelong number along $Y$, and
obtain a psh function --- is available in \Cref{part:bfj}: such a statement
requires a theory of psh functions with prescribed singularities along a
divisor, and that is not developed here.
\end{example}

\begin{remark}[Which images are admissible]
\label{rem:ch30-admissible-images}
Since $P(\varphi)$ contains no quasimonomial point
(\Cref{prop:ch21-polar-structure}\,\cref{itm:ch21-pol:skeleta}) and has empty
interior (\cref{itm:ch21-pol:interior}), hypothesis
\eqref{eq:ch30-pullback-hypothesis} holds whenever $f^{\An}(Y^{\An})$
meets some dual complex $\Delta_{\cX}\subset X^{\An}$, in particular
whenever $f$ is dominant.  It can fail only when $f(Y)$ is a proper
subvariety of $X$, and then only for those $\varphi$ whose polar set
happens to contain $f(Y)^{\An}$.  The failure is therefore confined to a
situation which is invisible on skeleta and which
\Cref{part:bfj} never needs: all applications of
\Cref{thm:ch30-pullback-psh} below have $f$ surjective.
\end{remark}

\begin{remark}[Pullback is injective but far from surjective]
\label{rem:ch30-pullback-not-onto}
If $f$ is surjective, then $f^*$ is injective on functions, since
$f^{\An}$ is surjective.  Surjectivity onto $\PSH(Y,f^*\theta)$ is not guaranteed: a potential
upstairs can separate points in a fibre and then cannot be a pullback.
For example, under $f(z)=z^2$ on $\mathbf P^1$, the rational points
$z=1$ and $z=-1$ have the same image.  With the degree-two pulled-back
reference, the model potential
$-\min\{1,\max\{0,v(z-1)\}\}$ takes different values at them and is
psh by the curve slope formula.  The identity map, in contrast, is
surjective on potentials.
For a fibrewise constant function, positivity descent requires its own
argument; lack of such an argument here is not a counterexample.
\Cref{sec:ch30-descent} supplies precise sufficient conditions for field
extensions, and \Cref{rem:ch30-psh-grows} gives an actual example of
potentials that do not descend.
\end{remark}

\section{Finite morphisms}
\label{sec:ch30-finite}

\subsection{The setting and the one number that matters}
\label{sec:ch30-finite-setting}

Let $f:Y\to X$ be a finite surjective morphism of degree $d$ between
smooth projective $K$-varieties of the same dimension $n$.  Recall that
$d=[K(Y):K(X)]$ and that, $f$ being finite and flat over the smooth
locus, every fibre of $f$ has $d$ points counted with multiplicities and
residue degrees.  If $\theta$ is represented by a model class on $\mathcal X$, choose
a model of $Y$ and use the normalized graph closure in its product with
$\mathcal X$, as in \Cref{sec:ch30-pullback-class}; then take a
dominating SNC model.  This supplies the model morphism needed to
represent $f^*\theta$.

Everything in this section is a consequence of the single algebraic identity
\begin{equation}
\label{eq:ch30-finite-projection-input}
 \bigl(\tilde f^*D\cdot\tilde f^*\mathcal L_1\cdots
       \tilde f^*\mathcal L_n\bigr)_{\cY}
 =d\,\bigl(D\cdot\mathcal L_1\cdots\mathcal L_n\bigr)_{\cX},
\end{equation}
valid for every vertical Cartier divisor $D$ on $\cX$ and every
$R$-morphism $\tilde f:\cY\to\cX$ of models whose generic fibre is
generically finite of degree $d$; this is
\Cref{thm:model-projection-formula}.  The reader who remembers only
\eqref{eq:ch30-finite-projection-input} can reconstruct
\Cref{thm:finite-ma-projection}, \Cref{cor:finite-pullback-energy} and their
mixed and finite-energy refinements without further thought, because the
Chambert--Loir measure and the energy are defined by pairing model functions
against vertical intersection numbers
(\eqref{eq:metrics-cl-pairing}, \eqref{eq:model-energy}).

\begin{proposition}\label{prop:finite-pullback-psh}
For a model function $\varphi_D$,
\begin{equation}
\label{eq:finite-pullback-model-function}
                  f^*\varphi_D=\varphi_D\circ f^{\An}
                  =\varphi_{f^*D}.
\end{equation}
Pullback preserves semipositive model forms and model psh functions.  If
$\theta$ has ample class, then $f^*\theta$ also has ample class and
\begin{equation}
\label{eq:finite-pullback-psh}
 \varphi\in\PSH(X,\theta)
 \quad\Longrightarrow\quad
 \varphi\circ f^{\An}\in\PSH(Y,f^*\theta).
\end{equation}
The same assertion holds for continuous semipositive metric weights.
\end{proposition}

\begin{proof}
The equality in \eqref{eq:finite-pullback-model-function} is checked from a
local equation of $D$; it is
\Cref{lem:ch30-pullback-model-calculus}\,\cref{itm:ch30-pbm:divisor}.
Relative nefness is preserved by pullback
(\Cref{prop:nef-pullback}), which gives the assertion about semipositive
model forms and hence about $\theta$-psh model functions through
\eqref{eq:ch30-pullback-ddc}.

For ampleness of the pulled-back class, recall that the pullback of an
ample line bundle along a finite morphism is ample: if $A$ is ample on
$X$ then $f^*A$ is ample on $Y$, because $f$ is finite, hence
quasi-affine, and the pullback of an ample sheaf along a quasi-affine
morphism of noetherian schemes is ample \cite[II, 5.1.12]{EGAII}.  This is
not \Cref{sec:ch20-imported}\,\cref{itm:ch20-import:ega}, which is the
fibrewise openness of \emph{relative} ampleness and says nothing about
pullback.  Applying this to
each rational ample class in an ample-cone decomposition of
$\{\theta\}$, as in \eqref{eq:ch20-ample-decomposition}, shows that
$f^*\{\theta\}$ is ample, and therefore that \eqref{eq:ch20-standing} holds on $Y$ with
reference form $f^*\theta$.

The implication \eqref{eq:finite-pullback-psh} is now the case of
\Cref{thm:ch30-pullback-psh} in which $f$ is surjective; hypothesis
\eqref{eq:ch30-pullback-hypothesis} is automatic there.  For continuous
potentials one may argue more directly and without decreasing
regularization: approximate $\varphi$ uniformly by $\theta$-psh model
functions (\Cref{thm:ch21-continuous-zhang}), compose with $f^{\An}$, and
use that composition with a fixed map does not increase the sup-norm, so
that the pulled-back sequence converges uniformly; then apply
\Cref{thm:ch21-stability}\,\cref{itm:ch21-stab:unif}.  This is
\Cref{prop:ch21-pullback}\,\cref{itm:ch21-pb:cont}.
\end{proof}

\begin{remark}[Where surjectivity is used, and where it is not]
\label{rem:ch30-finite-surjectivity}
Under the standing irreducibility and equal-dimension hypotheses, a finite
morphism has closed image of the same dimension and is therefore
surjective.  Surjectivity is used in
\eqref{eq:finite-pullback-psh} only to guarantee the polar condition and to
identify suprema.  For a non-dominant morphism of smaller-dimensional
source, the pullback theorem still applies when the restricted class is
ample, but one must check the polar condition explicitly, as
\Cref{ex:ch30-polar-pullback-fails} shows.
\end{remark}

\subsection{Pushforward of measures}
\label{sec:ch30-finite-pushforward}

For a Radon measure $\nu$ on $Y^{\An}$, pushforward is unambiguous:
\begin{equation}
\label{eq:measure-pushforward}
       f_*\nu(B)=\nu((f^{\An})^{-1}B),\qquad
       \int_X h\,\rd(f_*\nu)=\int_Y h\circ f^{\An}\,\rd \nu.
\end{equation}
The elementary properties used below are collected once.

\begin{lemma}
\label{lem:ch30-pushforward-basic}
Let $g:Z\to W$ be a continuous map of compact Hausdorff spaces and let
$g_*$ be defined by \eqref{eq:measure-pushforward}.
\begin{enumerate}
\item\label{itm:ch30-pf:radon}
$g_*\nu$ is a positive Radon measure on $W$ with
$g_*\nu(W)=\nu(Z)$; in particular $g_*$ preserves total mass.
\item\label{itm:ch30-pf:linear}
$g_*$ is linear and order preserving, and $(g'\circ g)_*=g'_*\circ g_*$.
\item\label{itm:ch30-pf:weak}
$g_*$ is continuous for the weak topologies: if $\nu_\alpha\to\nu$
weakly, then $g_*\nu_\alpha\to g_*\nu$ weakly.
\item\label{itm:ch30-pf:support}
$\Supp(g_*\nu)=g(\Supp\nu)$, and $g_*\delta_z=\delta_{g(z)}$.
\end{enumerate}
\end{lemma}

\begin{proof}
\cref{itm:ch30-pf:radon}  The set function $B\mapsto\nu(g^{-1}B)$ is a
Borel measure because $g^{-1}$ preserves countable disjoint unions, and it
is finite with total mass $\nu(Z)$.  On a compact Hausdorff space every
finite Borel measure that is inner regular on open sets is Radon, and
regularity is inherited because $g^{-1}$ of an open set is open and
$g^{-1}$ of a compact set is closed, hence compact.  Alternatively, and
this is the form used below, define $g_*\nu$ by the Riesz representation
theorem from the positive linear functional
$h\mapsto\int_Z h\circ g\,\rd\nu$ on $C^0(W)$
(\Cref{thm:riesz-interface}); the two definitions agree by a monotone class
argument, and the second makes the second identity in
\eqref{eq:measure-pushforward} a definition.

\cref{itm:ch30-pf:linear} and \cref{itm:ch30-pf:weak} are immediate from the
functional description: $h\circ g\in C^0(Z)$ whenever $h\in C^0(W)$, so
$\int h\,\rd(g_*\nu_\alpha)=\int(h\circ g)\,\rd\nu_\alpha$ converges to
$\int(h\circ g)\,\rd\nu=\int h\,\rd(g_*\nu)$.

\cref{itm:ch30-pf:support}  If $U\subset W$ is open and disjoint from
$g(\Supp\nu)$, then $g^{-1}U$ is open and disjoint from $\Supp\nu$, so
$g_*\nu(U)=0$; hence $\Supp(g_*\nu)\subseteq\overline{g(\Supp\nu)}
=g(\Supp\nu)$, the last equality because $\Supp\nu$ is compact and $g$
continuous.  Conversely if $w=g(z)$ with $z\in\Supp\nu$ and $U\ni w$
is open, then $g^{-1}U$ is an open neighbourhood of $z$, so
$g_*\nu(U)=\nu(g^{-1}U)>0$.
\end{proof}

The identity $g_*\delta_z=\delta_{g(z)}$ is the whole of the projection
formula on the model locus, because model Monge--Amp\`ere measures are finite
atomic measures carried by divisorial points
(\Cref{prop:ch24-model-support}\,\cref{itm:ch24-ms:atomic}).  Everything
reduces to identifying which divisorial point upstairs maps to which
divisorial point downstairs, and with what weight.

\subsection{The projection formula for Monge--Amp\`ere measures}
\label{sec:ch30-finite-projection}

The measure normalization is the one fixed in \eqref{eq:appJ-cl-measure};
in particular no probability normalization is in force and every measure
below has the total mass recorded in \eqref{eq:appJ-cl-total-mass}.

\begin{theorem}\label{thm:finite-ma-projection}
Let $\varphi_1,\ldots,\varphi_n$ be $\theta$-psh model potentials or
continuous $\theta$-psh potentials.  Then
\begin{equation}
\label{eq:finite-ma-projection}
\begin{split}
 f_*\bigl(& (f^*\theta+\ddc f^*\varphi_1)\wedge\cdots\wedge
                 (f^*\theta+\ddc f^*\varphi_n)\bigr)\\
 & =d\,(\theta+\ddc\varphi_1)\wedge\cdots\wedge
                 (\theta+\ddc\varphi_n).
\end{split}
\end{equation}
In particular,
\begin{equation}
\label{eq:finite-diagonal-projection}
 f_*\MA_{f^*\theta}(f^*\varphi)=d\MA_\theta(\varphi),
 \qquad
 (\{f^*\theta\}^n)=d(\{\theta\}^n).
\end{equation}
\end{theorem}

\begin{proof}
\emph{Step 1: reduction to a test against model functions.}  Both sides of
\eqref{eq:finite-ma-projection} are positive Radon measures on $X^{\An}$.
Rational model functions are dense in $C^0(X^{\An})$ for the uniform norm
(\Cref{prop:metrics-model-function-density}), so two Radon measures agreeing
against every model function agree.  It therefore suffices to prove
\begin{equation}
\label{eq:ch30-finite-test}
 \int_{Y^{\An}}(h\circ f^{\An})\,
  \bigwedge_{i=1}^n(f^*\theta+\ddc f^*\varphi_i)
 =d\int_{X^{\An}}h\,\bigwedge_{i=1}^n(\theta+\ddc\varphi_i)
\end{equation}
for every rational model function $h$ on $X$, where the left-hand side has been
rewritten using \eqref{eq:measure-pushforward}.

\emph{Step 2: the model case.}  Assume first that all $\varphi_i$ are
$\theta$-psh model functions and that $h=\varphi_{D_h}$.  Choose a
single model $\cX$ on which $\theta$, all the $\varphi_i$ and $h$
are determined, so that
$\theta+\ddc\varphi_i=[\mathcal L_i]$ for real model classes
$\mathcal L_i$ and $h=\varphi_{D_h}$ with $D_h$ vertical
$\mathbf Q$-Cartier; all intersections extend by real multilinearity.
Then choose a model $\cY$ of $Y$ with an $R$-morphism
$\tilde f:\cY\to\cX$ extending $f$, as in
\Cref{sec:ch30-pullback-class}.  By
\Cref{lem:ch30-pullback-model-calculus} the pulled-back data are
$f^*\theta+\ddc f^*\varphi_i=[\tilde f^*\mathcal L_i]$ and
$h\circ f^{\An}=\varphi_{\tilde f^*D_h}$.  The Chambert--Loir pairing
\eqref{eq:metrics-cl-pairing} converts each side of
\eqref{eq:ch30-finite-test} into a vertical intersection number:
\[
 \int_{Y^{\An}}\varphi_{\tilde f^*D_h}\,
   \bigwedge_i[\tilde f^*\mathcal L_i]
 =\bigl(\tilde f^*D_h\cdot\tilde f^*\mathcal L_1\cdots
        \tilde f^*\mathcal L_n\bigr)_{\cY},
\]
and similarly downstairs.  Now $\tilde f$ is proper and dominant with
generically finite generic fibre of degree $d$, so
\eqref{eq:ch30-finite-projection-input} applies and gives exactly
\eqref{eq:ch30-finite-test}.

\emph{Step 3: the continuous case.}  Let now the $\varphi_i$ be
continuous $\theta$-psh.  By \Cref{thm:ch21-continuous-zhang} choose
$\theta$-psh model functions $\varphi_i^{(m)}\to\varphi_i$ uniformly on
$X^{\An}$.  Then $f^*\varphi_i^{(m)}\to f^*\varphi_i$ uniformly on
$Y^{\An}$, since composition with $f^{\An}$ does not increase the
sup-norm.  By \Cref{prop:ch24-continuous-uniform}\,\cref{itm:ch24-cu:unif}
the mixed measures converge weakly on each side, and by
\Cref{lem:ch30-pushforward-basic}\,\cref{itm:ch30-pf:weak} pushforward is
weakly continuous.  Passing to the limit in the model identity proved in
Step~2 gives \eqref{eq:finite-ma-projection}.

\emph{Step 4: the two consequences.}  Taking all $\varphi_i$ equal gives
the first identity of \eqref{eq:finite-diagonal-projection}.  Taking
$\varphi_i=0$ and evaluating total masses gives, by
\Cref{lem:ch30-pushforward-basic}\,\cref{itm:ch30-pf:radon} and
\eqref{eq:appJ-cl-total-mass},
\[
 \begin{aligned}
 (\{f^*\theta\}^n)
 &=\MA_{f^*\theta}(0)(Y^{\An})\\
 &=\bigl(f_*\MA_{f^*\theta}(0)\bigr)(X^{\An})\\
 &=d\,\MA_\theta(0)(X^{\An})\\
 &=d(\{\theta\}^n),
 \end{aligned}
\]
which is the second.  Of course the last identity is also the algebraic
projection formula on the generic fibre; the point of deriving it this way
is that it is the mass test of \Cref{sec:appJ-unit-tests} applied to
\eqref{eq:finite-ma-projection}, and a sign or normalization error anywhere
in the chapter would be detected by it.
\end{proof}

The restriction to model and continuous potentials in
\Cref{thm:finite-ma-projection} is not necessary.  The extension to bounded
potentials is the first place where the monotone machinery of
\Cref{chap:bfj-ma} is used, and it is worth writing out because the same
three-line argument recurs verbatim in
\Cref{sec:ch30-birational,sec:ch30-base-change}.

\begin{theorem}
\label{thm:ch30-finite-bounded-projection}
Let $f:Y\to X$ be finite surjective of degree $d$ between smooth
projective $K$-varieties of dimension $n$, and let
$\varphi_1,\ldots,\varphi_n$ be bounded $\theta$-psh functions on
$X$.  Then $f^*\varphi_1,\ldots,f^*\varphi_n$ are bounded
$f^*\theta$-psh functions on $Y$ and
\begin{equation}
\label{eq:ch30-bounded-projection}
 f_*\MA_{f^*\theta}(f^*\varphi_1,\ldots,f^*\varphi_n)
 =d\,\MA_\theta(\varphi_1,\ldots,\varphi_n).
\end{equation}
\end{theorem}

\begin{proof}
Boundedness of $f^*\varphi_i$ is clear, and psh-ness is
\Cref{prop:finite-pullback-psh}.  By
\Cref{thm:bfj-decreasing-regularization} choose, for each $i$, a
decreasing net $\varphi_i^\alpha\searrow\varphi_i$ of $\theta$-psh model
functions.  Each net is already bounded below by the bounded limit.
Restrict to a cofinal tail below one fixed model majorant, and use the
finite product of the directed index sets for the $n$ potentials.
The resulting nets are uniformly bounded above and below.  Then
$f^*\varphi_i^\alpha\searrow f^*\varphi_i$ is a decreasing net of bounded
$f^*\theta$-psh model functions.

Fix $h\in C^0(X^{\An})$; we must show that the two sides of
\eqref{eq:ch30-bounded-projection} agree against $h$.  Since model
functions are uniformly dense in $C^0(X^{\An})$ and both sides are Radon
measures of the same finite mass by
\Cref{thm:ch24-mass} and
\Cref{lem:ch30-pushforward-basic}\,\cref{itm:ch30-pf:radon}, we may assume
$h$ is a rational model function.  By \Cref{lem:ch24-psh-difference},
write $h=(h_+-h_-)/(2t)$ with $\theta$-psh model functions $h_\pm$
and $t>0$.  Their pullbacks are $f^*\theta$-psh.
Apply \eqref{eq:weighted-monotone-ma} with the constant nets $h_\pm$
downstairs and $h_\pm\circ f^{\An}$ upstairs; subtract and divide by
$2t$.  Thus it applies to the displayed model test $h$ by linearity.  Set
\[
 \begin{aligned}
 \mu_\alpha^Y
   &:=\MA_{f^*\theta}(f^*\varphi_1^\alpha,\ldots,
      f^*\varphi_n^\alpha),\\
 \mu^Y
   &:=\MA_{f^*\theta}(f^*\varphi_1,\ldots,f^*\varphi_n),\\
 \mu_\alpha^X
   &:=\MA_\theta(\varphi_1^\alpha,\ldots,\varphi_n^\alpha),\\
 \mu^X
   &:=\MA_\theta(\varphi_1,\ldots,\varphi_n).
 \end{aligned}
\]
The two weighted convergences are
\[
 \int_{Y^{\An}}(h\circ f^{\An})\,\rd\mu_\alpha^Y
 \longrightarrow
 \int_{Y^{\An}}(h\circ f^{\An})\,\rd\mu^Y,
 \qquad
 \int_{X^{\An}}h\,\rd\mu_\alpha^X
 \longrightarrow
 \int_{X^{\An}}h\,\rd\mu^X.
\]
For every $\alpha$ the two left-hand sides are related by
\eqref{eq:ch30-finite-test}, that is, the first equals $d$ times the
second.  Passing to the limit gives the identity against $h$, and hence
\eqref{eq:ch30-bounded-projection}.
\end{proof}

\begin{remark}[Why the test function must be pulled back]
\label{rem:ch30-test-function}
Pushforward is tested against a continuous function $h$ downstairs,
so its test upstairs is $h\circ f^{\An}$.  The proof first uses rational
model tests, decomposed into psh model differences by
\Cref{lem:ch24-psh-difference}; it then reaches arbitrary continuous
tests by uniform density and the common mass bound.
An arbitrary continuous function is not asserted to be a difference of
two continuous psh functions.  Nor is there a corresponding pushforward
operation on model functions.  These are the reasons for choosing the
test space downstairs.
\end{remark}

\subsection{Energy, and the finite-energy class}
\label{sec:ch30-finite-energy}

\begin{corollary}\label{cor:finite-pullback-energy}
For a model or continuous semipositive potential $\varphi$,
\begin{equation}
\label{eq:finite-pullback-energy}
               E_{f^*\theta}(f^*\varphi)
               =dE_\theta(\varphi).
\end{equation}
\end{corollary}

\begin{proof}
Apply \eqref{eq:finite-ma-projection} to every mixed term in the energy
formula \eqref{eq:model-energy}, first for model potentials and then by
uniform approximation.  In detail, for a $\theta$-psh model function
$\varphi$,
\[
\begin{aligned}
 E_{f^*\theta}(f^*\varphi)
 &=\frac1{n+1}\sum_{j=0}^n\int_{Y^{\An}}f^*\varphi\,
    \MA_{f^*\theta}\bigl((f^*\varphi)^{(j)},0^{(n-j)}\bigr)\\
 &=\frac1{n+1}\sum_{j=0}^n\int_{X^{\An}}\varphi\,
    \Bigl(f_*\MA_{f^*\theta}\bigl((f^*\varphi)^{(j)},0^{(n-j)}\bigr)\Bigr),
\end{aligned}
\]
by \eqref{eq:measure-pushforward} applied to the function $\varphi$, and
each pushforward equals $d\,\MA_\theta(\varphi^{(j)},0^{(n-j)})$ by
\eqref{eq:finite-ma-projection}.  Summing gives
\eqref{eq:finite-pullback-energy}.  Continuous potentials follow by uniform
approximation and \Cref{prop:ch26-energy-lipschitz}.
\end{proof}

The energy identity extends to the whole finite-energy class, and the
extension is an equivalence: $f^*$ neither creates nor destroys finite
energy.  This is the statement actually used when a variational solution is
transported along a finite cover.

\begin{theorem}
\label{thm:ch30-finite-energy-transfer}
Let $f:Y\to X$ be finite surjective of degree $d$ as above and let
$\varphi\in\PSH(X,\theta)$.  Then
\begin{equation}
\label{eq:ch30-finite-energy-general}
 E_{f^*\theta}(f^*\varphi)=d\,E_\theta(\varphi)
 \qquad\text{in }\mathbf R\cup\{-\infty\},
\end{equation}
and consequently
\[
 \varphi\in\cE^1(X,\theta)
 \iff
 f^*\varphi\in\cE^1(Y,f^*\theta).
\]
Moreover \eqref{eq:ch30-bounded-projection} holds for
$\varphi_1,\ldots,\varphi_n\in\cE^1(X,\theta)$, with mixed
Monge--Amp\`ere measures in the sense of \Cref{thm:e1-mixed-products}.
\end{theorem}

\begin{proof}
\emph{Step 1: energy.}  Both sides of
\eqref{eq:ch30-finite-energy-general} are continuous along decreasing nets:
downstairs this is \Cref{prop:e1-monotone-continuity}, upstairs the same
statement applied on $Y$ with reference form $f^*\theta$, whose class is
ample by \Cref{prop:finite-pullback-psh}, so that
\eqref{eq:ch20-standing} holds on $Y$.  Take a decreasing net
$\varphi^\alpha\searrow\varphi$ of $\theta$-psh model functions; then
$f^*\varphi^\alpha\searrow f^*\varphi$, and
\eqref{eq:ch30-finite-energy-general} for the $\varphi^\alpha$ is
\eqref{eq:finite-pullback-energy}.  Passing to the limit on both sides,
which is legitimate in $\mathbf R\cup\{-\infty\}$ because both nets are
decreasing, gives the identity for $\varphi$.  The equivalence of finite
energy follows since $d>0$.

\emph{Step 2: mixed measures.}  Let
$\varphi_1,\ldots,\varphi_n\in\cE^1(X,\theta)$ and let
$\varphi_i^{\langle t\rangle}=\max\{\varphi_i,-t\}$ be their truncations,
which are bounded $\theta$-psh functions decreasing to $\varphi_i$ as
$t\to\infty$.  Pullback commutes with truncation,
\[
 f^*\bigl(\varphi_i^{\langle t\rangle}\bigr)
 =\max\{f^*\varphi_i,-t\}
 =(f^*\varphi_i)^{\langle t\rangle},
\]
since composition with a fixed map commutes with maxima.  By Step~1 each
$f^*\varphi_i$ lies in $\cE^1(Y,f^*\theta)$.  By
\Cref{thm:e1-mixed-products} the mixed measures of the truncations converge
weakly to those of the limits, on $X$ and on $Y$ alike, and pushforward
is weakly continuous
(\Cref{lem:ch30-pushforward-basic}\,\cref{itm:ch30-pf:weak}).  Since
\eqref{eq:ch30-bounded-projection} holds for the truncations by
\Cref{thm:ch30-finite-bounded-projection}, it passes to the limit.
\end{proof}

\begin{corollary}
\label{cor:ch30-finite-IJ}
Write $I$ and $J$ for the functionals of \Cref{def:ch26-ij}, formed on
$X$ with reference form $\theta$ and on $Y$ with reference form
$f^*\theta$.  For $\varphi,\psi\in\cE^1(X,\theta)$,
\begin{equation}
\label{eq:ch30-finite-IJ}
 I(f^*\varphi,f^*\psi)=d\,I(\varphi,\psi),
 \qquad
 J_{f^*\psi}(f^*\varphi)=d\,J_\psi(\varphi).
\end{equation}
\end{corollary}

\begin{proof}
By \eqref{eq:ch26-I-definition} and \eqref{eq:ch26-J-definition} both
functionals are built from integrals of the difference $\varphi-\psi$
against mixed Monge--Amp\`ere measures, and from energies.  For each such
integral, \eqref{eq:measure-pushforward} applied to the function
$\varphi-\psi$ --- legitimate on $\cE^1$, where each potential is
integrable against every relevant mixed measure by
\Cref{thm:e1-mixed-products}; assign the difference value zero on the
common polar set, which is null for these measures ---
turns the integral upstairs into an integral downstairs against the
pushforward, and the pushforward carries the factor $d$ by
\Cref{thm:ch30-finite-energy-transfer}.  The energy terms carry the same
factor by \eqref{eq:ch30-finite-energy-general}.
\end{proof}

\subsection{A finite pullback on a curve, carried to numbers}
\label{sec:ch30-finite-example}

The next example verifies every clause of
\Cref{thm:finite-ma-projection,cor:finite-pullback-energy} on an explicit
curve.  It is worth doing once with the models written down, because it is
the smallest situation in which the factor $d$ can be seen migrating from
the degree of a line bundle to the mass of a measure and from there to an
energy.

\begin{example}[A power map on the projective line]
\label{ex:ch30-power-map}
\index{projection formula!worked example}
Let $K$ be complete and discretely valued with $v_K(\varpi)=1$ and
residue characteristic zero, let $X=\mathbf P^1_K$ with affine coordinate
$z$, and let
\[
 f:\mathbf P^1_K\longrightarrow\mathbf P^1_K,
 \qquad
 [Z_0:Z_1]\longmapsto[Z_0^d:Z_1^d],
\]
be the $d$-th power map, $d\geq1$, which is finite surjective of degree
$d$ and extends to a finite $R$-morphism
$\mathbf P^1_R\to\mathbf P^1_R$.  Equip $L=\cO_{\mathbf P^1}(1)$ with the
model metric from $\mathbf P^1_{R}$ and let $\theta$ be its curvature,
so that $V=(\{\theta\})=\deg L=1$.  Then $f^*L\simeq\cO(d)$ with the
pulled-back model metric, and $f^*\theta$ is its curvature, of degree
$d$.

\emph{The map on the tree.}  Write $s(x):=v_x(z)\in[-\infty,+\infty]$ for
the tropical coordinate of \Cref{sec:ch08-p1-tree}, and let $x_s$ denote
the monomial point with $s(x_s)=s$; thus $x_0=x_G$ is the Gauss point,
$x_{+\infty}$ is the rigid point $z=0$ and $x_{-\infty}$ the rigid
point $z=\infty$.  Since $f^\#z=z^d$,
\begin{equation}
\label{eq:ch30-power-map-tropical}
 s\bigl(f^{\An}(x)\bigr)=v_x(z^d)=d\,s(x),
 \qquad\text{so}\qquad f^{\An}(x_s)=x_{ds}.
\end{equation}
The power map stretches the tropical coordinate by $d$ and in particular
fixes the Gauss point.

\emph{A model potential downstairs.}  Fix $a\in\mathbf Z_{>0}$ and put
\begin{equation}
\label{eq:ch30-power-map-potential}
 \varphi(x):=-\min\bigl\{a,\max\{0,s(x)\}\bigr\}.
\end{equation}
Thus $\varphi$ vanishes off the branch towards $z=0$, decreases with
slope $-1$ along the segment $0\leq s\leq a$, and is constant equal to
$-a$ beyond $s=a$.  It is a model function: let $\cX_a$ be the model
obtained from $\mathbf P^1_R$ by $a$ successive blowups of the closed
point $z=0$ of the special fibre, so that
$(\cX_a)_s=E_0+E_1+\cdots+E_a$ is a chain of projective lines, all of
multiplicity one, with $x_{E_i}=x_i$ and $E_0$ the strict transform of
the original special fibre.  Let $\cL$ be the pullback of
$\cO_{\mathbf P^1_R}(1)$; then $\varphi=\varphi_D$ with
$D=-\sum_{i=1}^{a}iE_i$, because $\varphi_D(x_{E_i})=a_i/b_i=-i$ by
\eqref{eq:ch20-sign}.

\emph{Its Monge--Amp\`ere measure.}  Apply the weighted graph Laplacian
\eqref{eq:ch20-curve-laplacian}: with $u_i=-i$ for $0\leq i\leq a$ and
adjacency $(E_i\cdot E_{i+1})=1$,
\[
\begin{aligned}
 (D\cdot E_0)&=u_1-u_0=-1,\\
 (D\cdot E_i)&=(u_{i-1}-u_i)+(u_{i+1}-u_i)=0\qquad(0<i<a),\\
 (D\cdot E_a)&=u_{a-1}-u_a=1 .
\end{aligned}
\]
Since $(\cL\cdot E_0)=1$ and $(\cL\cdot E_i)=0$ for $i\geq1$, the
model measure \eqref{eq:appJ-cl-measure} is
\begin{equation}
\label{eq:ch30-power-map-ma-down}
 \MA_\theta(\varphi)
 =\sum_{i=0}^a\bigl((\cL+D)\cdot E_i\bigr)\delta_{x_i}
 =\delta_{x_a},
 \qquad
 \MA_\theta(\varphi)(X^{\An})=1=V .
\end{equation}
All coefficients are nonnegative, so $\cL+D$ is relatively nef and
$\varphi$ is $\theta$-psh; the effect of $\varphi$ has been to move
the unit of mass from the Gauss point to $x_a$.

\emph{Upstairs.}  Assume $d\mid a$ and write $a=dm$.  Let $\cY_m$ be
the analogous chain model of $\mathbf P^1_K$ with components
$F_0,\ldots,F_m$ and $x_{F_j}=x_j$, and let $\cL'$ be the pullback of
$\cO_{\mathbf P^1_R}(d)$, so that $(\cL'\cdot F_0)=d$ and
$(\cL'\cdot F_j)=0$ for $j\geq1$; this represents $f^*\theta$, and
\begin{equation}
\label{eq:ch30-power-map-ma-ref}
 \MA_{f^*\theta}(0)=d\,\delta_{x_G},
 \qquad
 \MA_{f^*\theta}(0)(Y^{\An})=d=(\{f^*\theta\}) .
\end{equation}
By \eqref{eq:ch30-power-map-tropical} and
\eqref{eq:ch30-power-map-potential},
\[
 (f^*\varphi)(x_j)=\varphi(x_{dj})=-\min\{a,dj\}=-dj
 \qquad (0\leq j\leq m),
\]
so $f^*\varphi=\varphi_{D'}$ on $\cY_m$ with
$D'=-\sum_{j=1}^m dj\,F_j$, a model function with values in
$\mathbf Z$.  The same Laplacian computation, now with $u'_j=-dj$, gives
$(D'\cdot F_0)=-d$, $(D'\cdot F_j)=0$ for $0<j<m$, and
$(D'\cdot F_m)=d$, whence
\begin{equation}
\label{eq:ch30-power-map-ma-up}
 \MA_{f^*\theta}(f^*\varphi)=d\,\delta_{x_m},
 \qquad
 \MA_{f^*\theta}(f^*\varphi)(Y^{\An})=d=(\{f^*\theta\}) .
\end{equation}

\emph{The three checks.}  First, $f^{\An}(x_m)=x_{dm}=x_a$ by
\eqref{eq:ch30-power-map-tropical}, so by
\Cref{lem:ch30-pushforward-basic}\,\cref{itm:ch30-pf:support},
\[
 f_*\MA_{f^*\theta}(f^*\varphi)
 =d\,f_*\delta_{x_m}
 =d\,\delta_{x_a}
 =d\,\MA_\theta(\varphi),
\]
which is \eqref{eq:finite-diagonal-projection}.  Second, the masses are
$1$ downstairs and $d$ upstairs, in agreement with
$(\{f^*\theta\})=d(\{\theta\})$.  Third, the energies: by
\eqref{eq:model-energy} with $n=1$,
\[
 \begin{aligned}
 E_\theta(\varphi)
 &=\tfrac12\Bigl(\int\varphi\,\MA_\theta(\varphi)
    +\int\varphi\,\MA_\theta(0)\Bigr)\\
 &=\tfrac12\bigl(\varphi(x_a)+\varphi(x_G)\bigr)
 =\tfrac12(-a+0)=-\frac a2 ,
 \end{aligned}
\]
\[
 E_{f^*\theta}(f^*\varphi)
 =\tfrac12\Bigl(d\,f^*\varphi(x_m)+d\,f^*\varphi(x_G)\Bigr)
 =\tfrac12\bigl(d(-a)+0\bigr)=-\frac{da}2
 =d\,E_\theta(\varphi),
\]
which is \eqref{eq:finite-pullback-energy}.  With $d=3$ and $m=2$, so
that $a=6$, the numbers are
\[
 \MA_\theta(\varphi)=\delta_{x_6},\quad
 \MA_{f^*\theta}(f^*\varphi)=3\delta_{x_2},\quad
 E_\theta(\varphi)=-3,\quad
 E_{f^*\theta}(f^*\varphi)=-9 .
\]

\emph{Two warnings read off from the same computation.}  The measure
upstairs depends on the map $f$: its atom is $x_2$ when $d=3$,
while the downstairs atom is $x_6$.  The projection formula records
their relation; the support downstairs alone, without the map and its
degree data, does not specify the measure upstairs.  And the constant potentials show that the factor $d$ in
\eqref{eq:finite-pullback-energy} is a genuine feature of the normalization
and not an artefact: since $E_\theta(c)=cV=c$ and
$E_{f^*\theta}(c)=c\,d\,V=cd$, one has
$E_{f^*\theta}(f^*c)=dE_\theta(c)$ exactly, with the common normalization
$E_\theta(0)=E_{f^*\theta}(0)=0$.  Had energies been normalized by the
volume, as in \eqref{eq:bj-energy-definition}, the factor $d$ would have
disappeared from \eqref{eq:finite-pullback-energy} and reappeared in the
mass identity; the two conventions are compared once more in
\Cref{sec:ch30-gauss-mass}.
\end{example}

Equation \eqref{eq:finite-diagonal-projection} is often the correct way to
transport an equation.  It does \emph{not} define $f^*\mu$ for an
arbitrary Radon measure $\mu$.  Even for a point mass, a pullback requires
local degrees and residue-field weights; for a general continuous map there
is no canonical choice at all.  The precise statement, and the elementary
reason behind it, are in \Cref{prop:ch30-pushforward-not-inverse}.

\section{Birational morphisms}
\label{sec:ch30-birational}

\subsection{What changes and what does not}
\label{sec:ch30-birational-setting}

Let $f:Y\to X$ be projective birational between smooth projective
varieties.  The identity
\begin{equation}
\label{eq:birational-top-intersection}
                         (f^*L)^n=(L^n)
\end{equation}
is harmless, but the class $f^*L$ is generally not ample: it has degree
zero on curves contracted by $f$.  Hence the original ample-class BFJ
theory on $Y$ cannot simply be invoked with reference form $f^*\theta$.

To be precise about what remains available, fix once and for all a
semipositive representative of the class, which exists by
\Cref{prop:ch20-positive-representative}: after replacing $\theta$ by
$\theta':=\theta+\ddc\psi$ for a suitable model function $\psi$, we
simultaneously replace each potential $\varphi$ by
$\varphi':=\varphi-\psi$.  The total form
$\theta'+\ddc\varphi'=\theta+\ddc\varphi$, and hence its measure, is
unchanged; the corresponding energy is shifted by the explicit reference
term $E_\theta(\psi)$, by
\Cref{cor:ch20-reference-change,prop:ch26-reference-change}.  After this
reparameterization we may and do assume that $\theta$ itself is
semipositive with ample class.  Then
$f^*\theta$ is a semipositive closed model form
(\Cref{prop:nef-pullback}) whose class is nef and big, and:

\begin{itemize}
\item the constants and the pullbacks of $\theta$-psh model functions are
$f^*\theta$-psh model functions, by
\Cref{lem:ch30-pullback-model-calculus};
\item if $\varphi$ is a continuous $\theta$-psh potential, then
$f^*\varphi$ is the uniform limit of the pulled-back model approximants,
so continuous pullback remains available by
\Cref{prop:ch21-pullback}\,\cref{itm:ch21-pb:cont}; no assertion about an
arbitrary singular potential is made in this nef and big class;
\item the \emph{model} Monge--Amp\`ere operator is available upstairs,
because a mixed model measure is computed by vertical intersection numbers
on a model and never uses ampleness of the generic class; the same applies
to the intersection-theoretic energy \eqref{eq:ch26-energy-intersection} and
hence, by uniform approximation, to continuous potentials;
\item but the \emph{bounded} Monge--Amp\`ere operator
\Cref{thm:bounded-ma-extension}, whose construction quantifies over the
whole of \eqref{eq:ch24-standing} and in particular over an ample class, is
\emph{not} available on $(Y,f^*\theta)$; nor are the envelope theory of
\Cref{chap:bfj-envelopes}, the compactness theorem
\Cref{thm:bfj-psh-compactness}, the capacity theory of
\Cref{chap:bfj-capacity}, the class $\cE^1$ of \Cref{chap:bfj-energy}, or
the variational method of
\Cref{chap:bfj-differentiability,chap:bfj-finite-energy-measures}.
\end{itemize}

This list is the reason the present section proves projection formulae and
nothing else.  There are extensions of global pluripotential theory to nef
and big classes, but they require their own definitions of full mass and
singularity type.  \Cref{prop:birational-model-projection} is not a
substitute for that theory.  In particular, we do not claim that an
arbitrary BFJ current pulls back through a birational map.

\subsection{The valuative bijection}
\label{sec:ch30-birational-bijection}

A birational morphism induces a bijection between the zero-support loci,
which contain the supports of all \emph{model} Monge--Amp\`ere measures.
This is the structural fact that upgrades the model projection formula from
a pushforward identity to an identification on its finite atomic supports.
Continuous measures are obtained as weak limits and need not remain
supported on this locus.

\begin{lemma}
\label{lem:ch30-birational-bijection}
\index{birational morphism!valuative bijection}
Let $f:Y\to X$ be a proper birational morphism of integral proper
$K$-varieties.  Write
\[
 X^{\An}_0:=\{x\in X^{\An}\mid \Supp(x)=0\}
\]
for the set of points of zero support, that is, of absolute values on the
function field $K(X)$ extending $|\cdot|_K$, and similarly for $Y$.
Then $f^{\An}$ restricts to a bijection
$Y^{\An}_0\to X^{\An}_0$ which
\begin{enumerate}
\item\label{itm:ch30-bb:val}
is the identity under the identification $K(Y)=K(X)$;
\item\label{itm:ch30-bb:divisorial}
carries divisorial points to divisorial points and quasimonomial points to
quasimonomial points, preserving the normalization $x_E=b_E^{-1}\ord_E$;
\item\label{itm:ch30-bb:centers}
is compatible with centers: for $y\in Y^{\An}_0$ with center
$\xi\in\cY$ on a model $\cY$ of $Y$ and a model morphism
$\tilde f:\cY\to\cX$, the center of $f^{\An}(y)$ on $\cX$ is
$\tilde f(\xi)$.
\end{enumerate}
\end{lemma}

\begin{proof}
By \Cref{prop:berk-pair-description} a point of $X^{\An}$ is a pair
consisting of a schematic point $\xi\in X$ and an absolute value on
$\kappa(\xi)$ extending $|\cdot|_K$, taken up to the usual
identification; the support of the point is $\xi$.  Thus $X^{\An}_0$ is
the set of absolute values on $K(X)$ extending $|\cdot|_K$, and
likewise for $Y$.  Since $f$ is birational it induces an isomorphism
$f^\#:K(X)\to K(Y)$, and by the description of $f^{\An}$ in
\Cref{sec:berk-functoriality} the map $f^{\An}$ sends the absolute value
$|\cdot|_y$ on $K(Y)$ to $|\cdot|_y\circ f^\#$ on $K(X)$.  This is a
bijection with inverse induced by $(f^\#)^{-1}$, which proves the first
assertion and \cref{itm:ch30-bb:val}.

\cref{itm:ch30-bb:divisorial}  A divisorial point of $X^{\An}$ is by
definition $b_E^{-1}\ord_E$ for a prime divisor $E$ on some model of
$X$; this is a condition on the valuation of $K(X)$ alone, namely that
its value group be a discrete subgroup commensurable with
$v_K(K^\times)=\mathbf Z$ and its residue field have transcendence degree
$n$ over $k$; this is the relative equality case of
\Cref{lem:berk-abhyankar-valued}.  Indeed
\[
 v_{x_E}(K(X)^\times)=b_E^{-1}\mathbf Z,
 \qquad
 \operatorname{trdeg}_k\widetilde{K(X)}_{x_E}=n.
\]
The normalization is intrinsic because $v_{x_E}(\varpi)=1$, while
\[
 b_E=
 \bigl[v_{x_E}(K(X)^\times):v_K(K^\times)\bigr]
\]
is the finite index of the base value group in the full value group.  Thus
$b_E$ is \emph{not} the reciprocal of $v_{x_E}(\varpi)$.  These
conditions are preserved by \cref{itm:ch30-bb:val}.  The same applies to
quasimonomial points, whose defining relative Abhyankar equality is
intrinsic.

\cref{itm:ch30-bb:centers}  Immediate from the definition of a center: if
$\cO_{\cY,\xi}\subset R_y$ with $\mathfrak m_\xi=\mathfrak m_y\cap
\cO_{\cY,\xi}$, then $\cO_{\cX,\tilde f(\xi)}\subset\cO_{\cY,\xi}\subset
R_y$, and the maximal ideals correspond; uniqueness of the center is
\Cref{prop:centers-proper-separated}.
\end{proof}

\begin{warning}{The bijection is not a homeomorphism of the analytifications}
\Cref{lem:ch30-birational-bijection} does \emph{not} say that $f^{\An}$ is
injective.  Points supported on proper subvarieties behave differently:
if an exceptional prime divisor $E\subset Y$ of the generic fibre is
contracted to a point $p\in X$, all of $E^{\An}$ maps into
$\{x\in X^{\An}:\Supp(x)=p\}$.  The points of $E^{\An}$ can
have support at its generic point or at smaller subvarieties; they do
not all have support exactly $E$.  What the
lemma says is that this collapsing is invisible to Monge--Amp\`ere theory,
because a model Monge--Amp\`ere measure is carried by divisorial points
(\Cref{prop:ch24-model-support}\,\cref{itm:ch24-ms:atomic}), all of which lie
in $X^{\An}_0$.  It says nothing about polar sets or about points of type
$1$.
\end{warning}

\subsection{The projection formula in the birational case}
\label{sec:ch30-birational-projection}

\begin{proposition}\label{prop:birational-model-projection}
For semipositive model metrics, and by uniform approximation for continuous
semipositive metrics, the Chambert--Loir measure satisfies
\begin{equation}
\label{eq:birational-cl-projection}
              f_*c_1(f^*L,f^*\|\cdot\|)^n
              =c_1(L,\|\cdot\|)^n.
\end{equation}
The same equality holds for mixed measures of pulled-back metrized line
bundles.
\end{proposition}

\begin{proof}
Resolve the rational map between chosen models: given models $\cX$ of
$X$ carrying the metrized data and $\cY_0$ of $Y$, let $\cY$ be a
normalization of the graph closure of $f$ inside $\cY_0\times_R\cX$,
so that there is an $R$-morphism $\tilde f:\cY\to\cX$ whose generic
fibre is $f$, hence generically finite of degree $d=1$.  Testing both
sides of \eqref{eq:birational-cl-projection} against a model function
$h=\varphi_{D_h}$ on $X$, and using
\eqref{eq:measure-pushforward} together with
\eqref{eq:ch30-pullback-model-function}, reduces the identity to
\[
 \bigl(\tilde f^*D_h\cdot(\tilde f^*\mathcal L)^n\bigr)_{\cY}
 =\bigl(D_h\cdot\mathcal L^n\bigr)_{\cX},
\]
which is \Cref{thm:model-projection-formula} with $d=1$.  Model functions
are uniformly dense in $C^0(X^{\An})$
(\Cref{prop:metrics-model-function-density}), so the two Radon measures
coincide.  Uniform approximation and weak continuity give the continuous
statement, exactly as in Step~3 of the proof of
\Cref{thm:finite-ma-projection}; the mixed case is identical, with
$\mathcal L^n$ replaced by $\mathcal L_1\cdots\mathcal L_n$.
\end{proof}

For model metrics, \Cref{lem:ch30-birational-bijection} upgrades this
pushforward identity to an identification on the supports.

\begin{theorem}
\label{thm:ch30-birational-identification}
Let $f:Y\to X$ be projective birational between smooth projective
$K$-varieties of dimension $n$, and let $\varphi_1,\ldots,\varphi_n$
be $\theta$-psh model functions on $X$, or continuous
$\theta$-psh functions.  Then
\begin{equation}
\label{eq:ch30-birational-bounded}
 f_*\MA_{f^*\theta}(f^*\varphi_1,\ldots,f^*\varphi_n)
 =\MA_\theta(\varphi_1,\ldots,\varphi_n),
\end{equation}
both measures have total mass $V$, and $f^{\An}$ restricts to a
measure-preserving bijection between the supports when all the potentials
are model functions; in that case the supports are contained in
$Y^{\An}_0$ and $X^{\An}_0$, respectively.  No support-containment
assertion is made for continuous potentials.  Moreover
\begin{equation}
\label{eq:ch30-birational-energy}
 E_{f^*\theta}(f^*\varphi)=E_\theta(\varphi)
\end{equation}
for such $\varphi$, both sides being computed by the
intersection-theoretic formula \eqref{eq:ch26-energy-intersection}.
\end{theorem}

\begin{proof}
Identity \eqref{eq:ch30-birational-bounded} for model potentials is
\Cref{prop:birational-model-projection}, whose proof is
\eqref{eq:ch30-finite-projection-input} with $d=1$; the passage to
continuous potentials is Step~3 of the proof of
\Cref{thm:finite-ma-projection}, which uses only uniform approximation and
the weak continuity of $f_*$, both of which are available for a
semipositive closed model form.  The notation upstairs denotes these
continuous Chambert--Loir products of pulled-back metrics, defined by
uniform model approximation; it does not invoke the ample-class singular
operator on $Y$.  Real model classes are handled by multilinearity.
Equality of masses is
\eqref{eq:birational-top-intersection}.  The statement about supports
follows, for model potentials, from
\Cref{lem:ch30-pushforward-basic}\,\cref{itm:ch30-pf:support} and
\Cref{lem:ch30-birational-bijection}: both measures are carried by
divisorial points, which lie in $Y^{\An}_0$ and $X^{\An}_0$, and on
those loci $f^{\An}$ is injective, so
\eqref{eq:ch30-birational-bounded} determines each measure from the other.
Weak limits of divisorial atomic measures can acquire support outside the
zero-support locus, so this argument does not extend the support assertion
to continuous metrics.  Finally
\eqref{eq:ch30-birational-energy} is the computation in the proof of
\Cref{cor:finite-pullback-energy} with $d=1$: every mixed term
transforms by \eqref{eq:ch30-birational-bounded}, and the
mixed vertical $(n+1)$-fold terms occurring in
\eqref{eq:ch26-energy-intersection} are preserved by
$\tilde f$ because $\tilde f$ has degree one.
\end{proof}

\begin{warning}{No bounded or finite-energy statement upstairs}
\Cref{thm:ch30-birational-identification} is deliberately not stated for
bounded or finite-energy potentials.  The obstruction is not the projection
formula, which is an identity between vertical intersection numbers and is
insensitive to positivity of the generic class; it is that the objects
$\MA_{f^*\theta}(\psi_1,\ldots,\psi_n)$ and $\cE^1(Y,f^*\theta)$ have
not been constructed.  Their construction in
\Cref{chap:bfj-ma,chap:bfj-energy} quantifies over
\eqref{eq:ch24-standing}, and ampleness of the generic class is used there
--- through \Cref{prop:ch20-positive-representative} and
\Cref{lem:ch24-psh-difference} --- to know that $\theta$-psh model
functions span a dense subspace of $C^0$.  For a nontrivial birational
$f$ that hypothesis fails on $Y$.  A reader who has at hand a
pluripotential theory in a nef and big class may of course combine it with
\eqref{eq:ch30-birational-bounded}; nothing in this book supplies one.
\end{warning}

\begin{remark}[Exceptional mass]
\label{rem:ch30-exceptional-mass}
Formula \eqref{eq:birational-cl-projection} controls the pushforward.  It
does not say that an unrelated semipositive metric on $f^*L$ has no mass
near the exceptional locus, nor that every solution upstairs descends.
Such conclusions require the metric itself to be pulled back or a separate
descent argument.

The distinction is worth restating in the language of
\Cref{thm:ch30-birational-identification}.  That theorem concerns measures
of the form $\MA_{f^*\theta}(f^*\varphi_1,\ldots)$, whose potentials are
themselves pullbacks.  For model potentials the supports correspond bijectively under
$f^{\An}$.  This support assertion was not extended to continuous
potentials.  A divisorial point has zero algebraic support, so its
image cannot have support at a proper generic-fibre center of $f$.
Its reduction center on a chosen model can nevertheless lie above
the specialization of that center.  Model reduction and algebraic
support are different data.  The projection formula for pulled-back
metrics does not by itself classify arbitrary semipositive metrics
upstairs.
\end{remark}

\subsection{A birational pullback on a surface, carried to numbers}
\label{sec:ch30-birational-example}

\begin{example}[Blowing up a rational point of the projective plane]
\label{ex:ch30-surface-blowup}
\index{projection formula!birational example}
Let $K$ be as in \eqref{eq:ch20-standing} with $v_K(\varpi)=1$, let
$X=\mathbf P^2_K$, let $L=\cO(1)$ with the model metric coming from
$\cX=\mathbf P^2_R$, and let $\theta$ be its curvature form, so that
\[
                V=(\{\theta\}^2)=(L^2)=1 .
\]

\emph{A model potential downstairs.}  Let $p_0\in\mathbf P^2(k)$ be a
$k$-rational point of the special fibre and let
$\rho:\cX'=\mathrm{Bl}_{p_0}\mathbf P^2_R\to\cX$ be its blowup, with
exceptional divisor $E\simeq\mathbf P^2_k$ and $\widetilde X$ the strict
transform of $\mathbf P^2_k$.  Since $\cX_s=\mathbf P^2_k$ is a smooth
Cartier divisor on the threefold $\cX$, its multiplicity at $p_0$ is
one, so
\[
                \cX'_s=\widetilde X+E,
                \qquad b_{\widetilde X}=b_E=1,
\]
and the associated divisorial points are $x_{\widetilde X}=x_G$, the
Gauss point of $\mathbf P^{2,\An}_K$, and $x_E=\ord_E$.  For
$c\in\mathbf Q$ put
\[
 \cL_c:=\rho^*\cO(1)+cE,
 \qquad
 \varphi_c:=\varphi_{cE},
\]
so that $\varphi_c(x_G)=0$ and $\varphi_c(x_E)=c$ by
\eqref{eq:ch20-sign}, and $\theta+\ddc\varphi_c=[\cL_c]$.

\emph{The semipositivity range.}  We have $\cO(E)|_E=\cO_E(-1)$, so
$\cL_c|_E=\cO_E(-c)$, which is nef exactly for $c\leq0$.  On
$\widetilde X=\mathrm{Bl}_{p_0}\mathbf P^2_k$, writing $H$ for the
pullback of a line and $\ell=E\cap\widetilde X$ for the exceptional curve,
one has $\cL_c|_{\widetilde X}=H+c\ell$ with $(H\cdot\ell)=0$ and
$(\ell^2)=-1$; testing against $\ell$ gives $-c\geq0$ and against the
strict transform $H-\ell$ of a line through $p_0$ gives $1+c\geq0$.
Hence
\begin{equation}
\label{eq:ch30-surface-nef-range}
 \varphi_c\ \text{is}\ \theta\text{-psh}
 \iff -1\leq c\leq 0 .
\end{equation}

\emph{The Monge--Amp\`ere measure.}  For $-1\leq c\leq0$,
\[
 (\cL_c|_E)^2=(\cO_E(-c))^2=c^2,
 \qquad
 (\cL_c|_{\widetilde X})^2=(H+c\ell)^2=1-c^2 ,
\]
so \eqref{eq:appJ-cl-measure} gives
\begin{equation}
\label{eq:ch30-surface-ma-down}
 \MA_\theta(\varphi_c)=(1-c^2)\,\delta_{x_G}+c^2\,\delta_{x_E},
 \qquad
 \MA_\theta(\varphi_c)(X^{\An})=1=V .
\end{equation}
For $c=-\tfrac12$ the split is $\bigl(\tfrac34,\tfrac14\bigr)$; for
$c=-1$ the whole unit of mass has migrated to $x_E$.

\emph{Blowing up the generic fibre.}  Let $p\in\mathbf P^2(K)$ be a
rational point whose reduction is $p_0$, let
$f:Y=\mathrm{Bl}_p\mathbf P^2_K\to X$ be its blowup, and let
$E_Y\subset Y$ be the exceptional curve of $f$.  Then
$(\{f^*\theta\}^2)=1$ while $(\{f^*\theta\}\cdot E_Y)=0$, so
$\{f^*\theta\}$ is nef and big but not ample: this is
\eqref{eq:birational-top-intersection} and the failure of ampleness, in the
smallest possible example.

\emph{A model upstairs.}  Let $\overline P\subset\cX$ be the closure of
$p$, a section of $\cX\to\Spec R$ meeting $\cX_s$ transversally at
$p_0$, and let $\overline P'\subset\cX'$ be its strict transform.  Since
$\overline P$ is transverse to $\cX_s$, the point
$q:=\overline P'\cap E$ does not lie on
$\ell=E\cap\widetilde X=\mathbf P(T_{p_0}\cX_s)$.  Put
\[
 \cY':=\mathrm{Bl}_{\overline P'}\cX'
 \xrightarrow{\ \pi'\ }\cX'\xrightarrow{\ \rho\ }\cX .
\]
Its generic fibre is $\mathrm{Bl}_p\mathbf P^2_K=Y$, so $\cY'$ is a
model of $Y$ with an $R$-morphism $\tilde f=\rho\circ\pi'$ to $\cX$
extending $f$.  The center $\overline P'$ is not contained in
$\cX'_s$, so
\[
 \cY'_s=(\pi')^*\cX'_s=\widetilde X'+E',
 \qquad b_{\widetilde X'}=b_{E'}=1,
\]
where $\widetilde X'\simeq\widetilde X$ (the blowup center misses it) and
$E'=\mathrm{Bl}_q E$.  The new exceptional divisor of $\pi'$ is a
$\mathbf P^1$-bundle over $\overline P'\simeq\Spec R$ and is therefore
horizontal, not a component of the special fibre.

\emph{The Monge--Amp\`ere measure upstairs.}  Let
$\cL'_c:=(\pi')^*\cL_c$, which represents
$f^*\theta+\ddc f^*\varphi_c$ by
\eqref{eq:ch30-pullback-ddc}.  Self-intersection numbers are preserved by
pullback along a birational morphism of surfaces, so
\[
 (\cL'_c|_{\widetilde X'})^2=1-c^2,
 \qquad
 (\cL'_c|_{E'})^2=(\cL_c|_E)^2=c^2 ,
\]
and therefore
\begin{equation}
\label{eq:ch30-surface-ma-up}
 \MA_{f^*\theta}(f^*\varphi_c)
 =(1-c^2)\,\delta_{x_{\widetilde X'}}+c^2\,\delta_{x_{E'}},
 \qquad
 \text{total mass}\ 1 .
\end{equation}

\emph{The checks.}  The divisorial points $x_{\widetilde X'}$ and
$x_{E'}$ are, as valuations of $K(Y)=K(X)$, equal to $x_G$ and
$x_E$; this is
\Cref{lem:ch30-birational-bijection}\,\cref{itm:ch30-bb:val} in the concrete
case, and one can see it directly, since $\pi'$ is an isomorphism near the
generic points of $\widetilde X$ and $E$.  Consequently
$f^{\An}(x_{\widetilde X'})=x_G$ and $f^{\An}(x_{E'})=x_E$, and
\eqref{eq:ch30-surface-ma-up} pushes forward to
\eqref{eq:ch30-surface-ma-down}: no factor appears, in contrast with
\Cref{ex:ch30-power-map}.

\emph{The energy.}  With $n=2$, the intersection form
\eqref{eq:ch26-energy-intersection} of the energy reads
\[
\begin{aligned}
 E_\theta(\varphi_c)
 & =\tfrac13\sum_{j=0}^2
   \bigl(cE\cdot\cL_c^{\,j}\cdot(\rho^*\cO(1))^{2-j}\bigr)\\
 & =\tfrac13\Bigl(
   3c\,(\rho^*\cO(1))^2\!\cdot\!E
   +3c^2(\rho^*\cO(1))\!\cdot\!E^2
   +c^3(E^3)\Bigr).
\end{aligned}
\]
where every displayed summand is legitimate because its first factor is the
vertical divisor $cE$, in the sense of
\Cref{lem:ch24-vertical-numbers}.  By the projection formula for $\rho$, the first two summands
vanish, since $\rho_*E=0$ and $\rho_*(E^2)=0$; and $(E^3)=(\cO_E(-1)^2)
=1$ for the exceptional divisor of the blowup of a smooth point of a
threefold.  Hence
\begin{equation}
\label{eq:ch30-surface-energy}
 E_\theta(\varphi_c)=\frac{c^3}{3},
 \qquad
 E_\theta(\varphi_{-1})=-\frac13,
 \qquad
 E_\theta(\varphi_{-1/2})=-\frac1{24} .
\end{equation}
The same value is obtained from the mixed formula
\eqref{eq:model-energy}: the only nonvanishing term is $j=2$, which
contributes $\varphi_c(x_E)\,(\cL_c|_E)^2=c\cdot c^2$.  Upstairs every one
of these vertical intersection numbers is computed on $\cY'$ and is
preserved by $\pi'$, so
$E_{f^*\theta}(f^*\varphi_c)=E_\theta(\varphi_c)$, as
\eqref{eq:ch30-birational-energy} asserts --- with the factor $1$, not
$d$.

\emph{The boundary of the example.}  Further vertical blowups of $\cY'$
produce new model functions on $Y^{\An}$, but twisting by the negative of
a new exceptional component does not automatically preserve relative
nefness.  One must test every curve in the new special fibre before calling
such a function $f^*\theta$-psh or attaching a positive
Monge--Amp\`ere measure to it.  No such additional metric is used here; the
numerical computation above concerns only the genuine pullbacks
$f^*\varphi_c$.
\end{example}

\section{Extension of the ground field}
\label{sec:ch30-base-change}

Let $K'/K$ be a finite extension of complete non-Archimedean fields, with
the absolute value on $K'$ extending the one on $K$, and let
\begin{equation}
\label{eq:ch30-base-change-map}
               \pi:X_{K'}^{\An}\longrightarrow X^{\An}
\end{equation}
be the canonical projection.  We keep the same absolute-value
normalization; we do not renormalize a uniformizer of $K'$ to have
valuation one.

\subsection{The base change map}
\label{sec:ch30-bc-map}

We first collect what is true for an arbitrary complete valued extension,
before specializing to the finite case in which models exist.

\begin{lemma}
\label{lem:ch30-base-change-map}
Let $K'/K$ be any complete valued extension and $X$ a proper
$K$-scheme.
\begin{enumerate}
\item\label{itm:ch30-bcm:cont}
$\pi$ is continuous, proper and surjective, and both spaces are compact
Hausdorff.
\item\label{itm:ch30-bcm:fibre}
For $x\in X^{\An}$,
$\pi^{-1}(x)\simeq\mathcal M\bigl(\mathcal H(x)\widehat\otimes_KK'\bigr)$.
\item\label{itm:ch30-bcm:support}
$\Supp\bigl(\pi(x')\bigr)$ is the image of $\Supp(x')$ under
$X_{K'}\to X$; in particular $\pi$ maps points of zero support to points
of zero support when $X$ is geometrically integral.
\item\label{itm:ch30-bcm:sup}
For any function $u$ on $X^{\An}$,
$\sup_{X_{K'}^{\An}}u\circ\pi=\sup_{X^{\An}}u$, and $u\circ\pi$ is
continuous, respectively upper semicontinuous, if and only if $u$ is.
\end{enumerate}
\end{lemma}

\begin{proof}
\cref{itm:ch30-bcm:cont} and \cref{itm:ch30-bcm:fibre} are recalled in
\Cref{sec:berk-functoriality}, the fibre description being Berkovich's
\cite[\S1.4 and \S3.4]{Berk12}; compactness is
\Cref{thm:berk-proper-compact}.  Properness of $\pi$ follows from
compactness of both spaces and continuity.

\cref{itm:ch30-bcm:support}  A point of $X_{K'}^{\An}$ is a pair
$(\xi',|\cdot|)$ with $\xi'\in X_{K'}$ and $|\cdot|$ an absolute value
on $\kappa(\xi')$ (\Cref{prop:berk-pair-description}); the map $\pi$
sends it to $(\xi,|\cdot|_{|\kappa(\xi)})$ where $\xi$ is the image of
$\xi'$.  If $X$ is geometrically integral, $X_{K'}$ is integral and
its generic point maps to the generic point of $X$.

\cref{itm:ch30-bcm:sup}  The first assertion is surjectivity of $\pi$;
the second follows because $\pi$ is a continuous surjection between
compact Hausdorff spaces, hence a topological quotient map, so a function
downstairs is continuous (upper semicontinuous) exactly when its pullback
is.
\end{proof}

\begin{remark}[The fibres are not points]
\label{rem:ch30-fibres-not-points}
Item \cref{itm:ch30-bcm:fibre} is the source of most of the difficulties in
this section, and it should be read as a warning.  For $X=\mathbf P^1$ and $K'=K(\sqrt{\varpi})$, the downstairs
closed point defined by $z^2-\varpi$ is a type-$1$ point with completed
residue field $K'$.  Its fibre consists of the two $K'$-rational
points $z=\sqrt{\varpi}$ and $z=-\sqrt{\varpi}$, since
$K'\otimes_KK'\simeq K'\times K'$.
A $K$-rational point, on the other hand, has one point in its fibre,
with completed residue field extended from $K$ to $K'$.  A statement about
$X^{\An}$ therefore has, in general, no single ``lift'' to
$X_{K'}^{\An}$.  What survives is the pair of operations
$u\mapsto u\circ\pi$ on functions and $\nu\mapsto\pi_*\nu$ on measures,
and \Cref{lem:ch30-base-change-map}\,\cref{itm:ch30-bcm:sup} says exactly
that the first of these loses nothing.
\end{remark}

\subsection{The normalization of the extended valuation}
\label{sec:ch30-bc-normalization}

Let $K'/K$ be finite, $e=e(K'/K)$ the ramification index and
$f=f(K'/K)=[k':k]$ the residue degree, where $k'$ is the residue field
of $R'=K'^\circ$.  A complete discretely valued field of residue characteristic zero is
defectless; its finite extensions are separable and their valuation rings
are finite over $R$.  Thus the fundamental equality
\begin{equation}
\label{eq:ch30-efn}
                    e\,f=[K':K]
\end{equation}
holds; see \cite[Ch.~VI]{ZariskiSamuelII}.  With the normalization fixed
above --- namely $v_{K'}$ restricts to $v_K$ on $K$ --- one has
\begin{equation}
\label{eq:ch30-uniformizer-value}
 v_{K'}(\varpi)=1,
 \qquad
 v_{K'}(\varpi')=\frac1e,
 \qquad
 v_{K'}(K'^\times)=\frac1e\mathbf Z .
\end{equation}
The value group therefore grows by the index $e$, while the residue field
grows by the degree $f$.  These are the two numbers that replace the
degree $d$ of \Cref{sec:ch30-finite}, and \eqref{eq:ch30-efn} is the
constraint that makes them cancel.

\begin{center}
\fbox{\begin{minipage}{0.91\textwidth}
\textbf{Normalization checkpoint.}
If one instead rescales the extended valuation so that a uniformizer of
$K'$ has value one, its restriction to $K$ is multiplied by the
ramification index.  Model functions, edge lengths, and energies then scale.
Formula \eqref{eq:base-change-ma} presupposes an extension of the original
absolute value without this renormalization.

Explicitly, write $v''=e\,v_{K'}$, so that $v''(\varpi')=1$ and
$v''|_K=e\,v_K$.  The identity map of sets
$X_{K'}^{\An}\to X_{K'}^{\An}$ intertwining the two normalizations is the
scaling $x\mapsto e\cdot x$ of \eqref{eq:appJ-scaling-action}, and under
it
\begin{equation}
\label{eq:ch30-ramified-renormalization}
 \varphi''=e\,\varphi',
 \qquad
 \ell''(\gamma)=e\,\ell'(\gamma),
 \qquad
 \MA''=\MA',
 \qquad
 E''=e\,E',
\end{equation}
where $\varphi'$ is a model weight, $\ell'$ an edge length of a dual
complex, and $E'$ an energy.  Monge--Amp\`ere \emph{measures} are
unchanged as measures --- only the labelling of their atoms moves --- while
weights, lengths and energies are multiplied by $e$.  A comparison
theorem that reports an unexplained factor $e$, or $[K':K]$, has almost
always mixed the two conventions.
\end{minipage}}
\end{center}

\subsection{Models, ramification, and residue degrees}
\label{sec:ch30-bc-models}

\begin{lemma}
\label{lem:ch30-base-change-model}
\index{base change!of models}
Assume \eqref{eq:ch20-standing}, let $X$ be geometrically integral, let
$K'/K$ be finite, and let $\cX$ be a normal model of $X$.  Let
$\cX'$ be the normalization of $\cX\times_RR'$ and
$\rho:\cX'\to\cX$ the induced morphism.
\begin{enumerate}
\item\label{itm:ch30-bcmod:finite}
$\cX'$ is a normal model of $X_{K'}$ and $\rho$ is finite and
surjective, with generic fibre the base-change morphism
$X_{K'}\to X$, of degree $[K':K]$.
\item\label{itm:ch30-bcmod:mult}
For a prime component $E$ of $\cX_s$ with multiplicity
$b_E=\ord_E(\varpi)$, and for a prime component $E'$ of $\cX'_s$
lying over $E$, put
\[
 m_{E'}:=\ord_{E'}(f_E),
 \qquad
 d_{E'}:=[\kappa(E'):\kappa(E)],
\]
where $f_E$ is a local equation of $E$ at its generic point.  Then
\begin{equation}
\label{eq:ch30-ramification-sum}
 \sum_{E'\mapsto E}m_{E'}d_{E'}=[K':K]=ef,
\end{equation}
and, writing $b_{E'}=\ord_{E'}(\varpi')$,
\begin{equation}
\label{eq:ch30-multiplicity-relation}
 \ord_{E'}(\varpi)=b_Em_{E'}=e\,b_{E'} .
\end{equation}
\item\label{itm:ch30-bcmod:points}
With the normalization \eqref{eq:ch30-uniformizer-value}, the divisorial
point of $X_{K'}^{\An}$ attached to $E'$ is
\[
 x_{E'}=\bigl(\ord_{E'}(\varpi)\bigr)^{-1}\ord_{E'}
       =(e\,b_{E'})^{-1}\ord_{E'},
\]
and $\pi(x_{E'})=x_E$.
\item\label{itm:ch30-bcmod:degrees}
Let $\cL_1,\ldots,\cL_n$ be model line bundles on $\cX$, with pullbacks
$\cL'_i=\rho^*\cL_i$.  Then, with degrees taken over $k'$ upstairs and
over $k$ downstairs,
\begin{equation}
\label{eq:ch30-degree-transfer}
 \bigl(\cL'_1|_{E'}\cdots\cL'_n|_{E'}\bigr)_{k'}
 =\frac{d_{E'}}{f}\,\bigl(\cL_1|_E\cdots\cL_n|_E\bigr)_{k} .
\end{equation}
\item\label{itm:ch30-bcmod:cancel}
Consequently
\begin{equation}
\label{eq:ch30-coefficient-cancellation}
 \sum_{E'\mapsto E}
   b_{E'}\bigl(\cL'_1|_{E'}\cdots\cL'_n|_{E'}\bigr)_{k'}
 =b_E\bigl(\cL_1|_E\cdots\cL_n|_E\bigr)_{k}.
\end{equation}
\end{enumerate}
\end{lemma}

\begin{proof}
\cref{itm:ch30-bcmod:finite}  The ring $R'$ is a finite $R$-module
because $R$ is a complete discrete valuation ring and $K'/K$ is finite,
so $\cX\times_RR'\to\cX$ is finite; it is also flat and surjective.  Its
generic fibre is $X\times_KK'=X_{K'}$, which is integral because $X$ is
geometrically integral, so $\cX\times_RR'$ is integral and its
normalization $\cX'$ is well defined; normalization is finite because
finite-type schemes over a complete discrete valuation ring are excellent,
hence Nagata (\cite[Tag 035S]{stacks-project}).  Thus $\rho$ is finite and
surjective, and $\cX'$ is a normal flat projective $R'$-scheme with
generic fibre $X_{K'}$, that is, a model.  Its degree is
$[K(X_{K'}):K(X)]=[K':K]$, the extension $K'\otimes_KK(X)$ being a field
by geometric integrality.

\cref{itm:ch30-bcmod:mult}  Localize at the generic point $\eta_E$ of
$E$.  Since $\cX$ is normal, $\cO_{\cX,\eta_E}$ is a discrete
valuation ring with fraction field $K(X)$ and residue field
$\kappa(E)$, and $\cX'$ is the integral closure of $\cX$ in
$K(X_{K'})$, so the semilocal ring of $\cX'$ over $\eta_E$ is the
integral closure of $\cO_{\cX,\eta_E}$ in the finite separable extension
$K(X_{K'})/K(X)$ of degree $[K':K]$.  Its maximal ideals correspond
bijectively to the components $E'$ lying over $E$, with ramification
indices $m_{E'}=\ord_{E'}(f_E)$ and residue degrees $d_{E'}$; finiteness of the integral closure makes this semilocal ring a finite
torsion-free, hence free, module over the downstairs discrete valuation
ring.  Its rank is $[K':K]$; reducing a local uniformizer gives the
length identity
\cite[Ch.~VI]{ZariskiSamuelII}, namely
\eqref{eq:ch30-ramification-sum}.  For
\eqref{eq:ch30-multiplicity-relation}, apply $\ord_{E'}$ to
$\varpi=u\,f_E^{\,b_E}$ at the generic point of $E$, where $u$ is a
unit, to get $\ord_{E'}(\varpi)=b_Em_{E'}$; and to
$\varpi=u'\varpi'^{\,e}$ to get $\ord_{E'}(\varpi)=e\,b_{E'}$.

\cref{itm:ch30-bcmod:points}  The normalization of a divisorial point is
determined by the requirement that it restrict to $v_{K'}$ on $K'$,
equivalently to $v_K$ on $K$; since $v_K(\varpi)=1$, the correct
scaling factor is $\ord_{E'}(\varpi)^{-1}$.  For the second assertion,
let $g\in K(X)^\times$.  At the generic point of $E$ write
$g=u\,f_E^{\,r}$ with $u$ a unit and $r=\ord_E(g)$; applying
$\ord_{E'}$ gives $\ord_{E'}(g)=m_{E'}\ord_E(g)$, whence
\[
 v_{x_{E'}}(g)=\frac{\ord_{E'}(g)}{b_Em_{E'}}
              =\frac{\ord_E(g)}{b_E}=v_{x_E}(g).
\]
Thus $x_{E'}$ restricts to $x_E$ on $K(X)$, which is exactly
$\pi(x_{E'})=x_E$.

\cref{itm:ch30-bcmod:degrees}  The morphism $g:=\rho|_{E'}:E'\to E$ is
finite and surjective of degree $d_{E'}$, and
$\cL'_i|_{E'}=g^*(\cL_i|_E)$.  The projection formula for $g$ gives
$\deg_k\bigl(g^*\alpha\bigr)=d_{E'}\deg_k(\alpha)$ for a top-degree class
$\alpha$ on $E$.  Since $E'$ is a scheme over $k'$, every closed
point of $E'$ has residue field containing $k'$, so
$\deg_k=f\cdot\deg_{k'}$ on $0$-cycles of $E'$.  Combining gives
\eqref{eq:ch30-degree-transfer}.

\cref{itm:ch30-bcmod:cancel}  Insert
$b_{E'}=b_Em_{E'}/e$ from \eqref{eq:ch30-multiplicity-relation} and
\eqref{eq:ch30-degree-transfer} into the left-hand side:
\[
 \sum_{E'\mapsto E}\frac{b_Em_{E'}}{e}\cdot\frac{d_{E'}}f
   \bigl(\cL_1|_E\cdots\cL_n|_E\bigr)_k
 =\frac{b_E}{ef}\Bigl(\sum_{E'\mapsto E}m_{E'}d_{E'}\Bigr)
   \bigl(\cL_1|_E\cdots\cL_n|_E\bigr)_k,
\]
and the bracket equals $ef$ by \eqref{eq:ch30-ramification-sum}.
\end{proof}

\begin{remark}[Reading the cancellation]
\label{rem:ch30-cancellation}
Identity \eqref{eq:ch30-coefficient-cancellation} is the whole content of
\Cref{thm:base-change-ma}, and it is worth isolating why it is not a
tautology.  Three different quantities move when the ground field is
extended: the multiplicity $b$ of a component in the special fibre
divisor, which is divided by $e$ up to the factor $m_{E'}$; the degree
of a line bundle on that component, which is multiplied by $d_{E'}$ and
divided by $f$; and the number of components over a fixed one, which is
constrained by \eqref{eq:ch30-ramification-sum}.  The product of the three
effects is $(1/e)(1/f)(ef)=1$.  Drop the convention that degrees on the
special fibre are taken over the residue field of the \emph{base of the
model}, and the identity acquires a spurious factor $f$; drop the
convention \eqref{eq:ch30-uniformizer-value} and it acquires a spurious
factor $e$.  Both mistakes are easy to make and both are detected by the
unit test of \Cref{sec:appJ-unit-tests}, since the total mass must remain
$(L_1\cdots L_n)$.
\end{remark}

\subsection{Model and continuous metrics under base change}
\label{sec:ch30-bc-metrics}

\begin{proposition}\label{prop:base-change-model-metrics}
Pullback takes model metrics to model metrics after normalization of the
base-changed model, and takes semipositive model metrics to semipositive
model metrics.  For a model weight $\varphi$,
\begin{equation}
\label{eq:base-change-weight}
                    \varphi_{K'}=\varphi\circ\pi.
\end{equation}
The analogous statement holds for continuous semipositive metrics by
uniform approximation.
\end{proposition}
\begin{proof}
Choose a normal projective model $\mathcal X$ of $X$ and a model line bundle
$\mathcal L$ representing the metric.  Form
$\mathcal X\times_{K^\circ}K'{}^\circ$ and let $\mathcal X'$ be its
normalization.  The generic fibre of $\mathcal X'$ is $X_{K'}$, and the
pullback $\mathcal L'$ of $\mathcal L$ is a model line bundle there.

To compare the metrics, work on a model open set on which a power
$\mathcal L^{\otimes q}$ is generated by a frame $e$.  If a rational section
$s$ is written $s=fe$ on the generic fibre, its model norm is determined by
$|f|$.  After scalar extension the same identity holds with $f$ viewed in
the base-changed affinoid algebra.  Evaluation at $x'\in X_{K'}^{\An}$ is
evaluation of $f$ at $\pi(x')$; hence the pulled-back weight is
$\varphi(\pi(x'))$.  These local identities glue and prove
\eqref{eq:base-change-weight}.

Suppose that $\mathcal L$ is relatively nef.  The morphism
$\rho:\mathcal X'\to\mathcal X$ is finite after passage through the finite
base change.  For every integral curve $C'$ in the special fibre, either
$\rho(C')$ is a curve and the projection formula gives
\[
 f(\mathcal L'\cdot C')_{k'}
   =[k(C'):k(\rho(C'))](\mathcal L\cdot\rho(C'))_k\geq0,
\]
or $\rho(C')$ is a point and $(\mathcal L'\cdot C')=0$.  Thus
$\mathcal L'$ is relatively nef, proving preservation of semipositivity.
Finally, if $\varphi_j\to\varphi$ uniformly with each $\varphi_j$ a
semipositive model weight, then
\[
 \lVert\varphi_j\circ\pi-\varphi\circ\pi\rVert_\infty
 \leq\lVert\varphi_j-\varphi\rVert_\infty\longrightarrow0.
\]
The continuous statement follows from the definition by uniform closure.
\end{proof}

Two supplements to \Cref{prop:base-change-model-metrics} are used below.
The first records that the semipositivity test may be run in either
direction, which is what makes descent possible in
\Cref{sec:ch30-descent}; the second identifies the pulled-back model
function in the notation of \Cref{lem:ch30-base-change-model}.

\begin{lemma}
\label{lem:ch30-semipositivity-descends}
Let $\rho:\cX'\to\cX$ be a finite surjective morphism of normal models
over $R'$ and $R$ respectively, and let $\cM$ be a line bundle on
$\cX$.  Then $\cM$ is relatively nef if and only if $\rho^*\cM$ is.
The same curve test applies to real model classes.
\end{lemma}

\begin{proof}
Necessity is \Cref{prop:nef-pullback}, or the computation in the proof of
\Cref{prop:base-change-model-metrics}.  For sufficiency, let
$C\subset\cX_s$ be an integral curve.  Since $\rho$ is finite and
surjective, $\rho^{-1}(C)$ has pure dimension one and dominates $C$;
choose an integral curve $C'\subseteq\rho^{-1}(C)$ with
$\rho(C')=C$, and let $\delta=[\kappa(C'):\kappa(C)]\geq1$.  With degrees taken over $k'$ upstairs and $k$ downstairs, the
projection formula gives
$f(\rho^*\cM\cdot C')_{k'}=\delta(\cM\cdot C)_k$,
where $f=[k':k]$.  Hence
$(\cM\cdot C)_k=(f/\delta)(\rho^*\cM\cdot C')_{k'}\geq0$.  As $C$ was
arbitrary, $\cM$ is relatively nef.
\end{proof}

\begin{lemma}
\label{lem:ch30-base-change-model-function}
In the notation of \Cref{lem:ch30-base-change-model}, let
$D=\sum_Ea_EE$ be a vertical $\mathbf R$-Cartier divisor on $\cX$.
Then $\rho^*D$ is vertical on $\cX'$,
$\varphi_{\rho^*D}=\varphi_D\circ\pi$, and the vertex values are related
by
\begin{equation}
\label{eq:ch30-model-function-vertices}
 \varphi_{\rho^*D}(x_{E'})=\frac{m_{E'}a_E}{b_Em_{E'}}
                          =\frac{a_E}{b_E}=\varphi_D(x_E).
\end{equation}
In particular the map $D\mapsto\rho^*D$ induces an injection of the group
of model functions on $X$ into that on $X_{K'}$, compatible with
$\ddc$ and with \eqref{eq:ch30-pullback-ddc}.
\end{lemma}

\begin{proof}
Verticality holds because $\rho^{-1}(\cX_s)=\cX'_s$.  The identity
$\varphi_{\rho^*D}=\varphi_D\circ\pi$ is the same local-equation
computation as in
\Cref{lem:ch30-pullback-model-calculus}\,\cref{itm:ch30-pbm:divisor}, and
\eqref{eq:ch30-model-function-vertices} is its numerical form: the
coefficient of $E'$ in $\rho^*D$ is $m_{E'}a_E$ and the multiplicity
of $E'$ in the divisor $\ord_{E'}(\varpi)$ used to normalize
$x_{E'}$ is $b_Em_{E'}$, by
\eqref{eq:ch30-multiplicity-relation}.  Injectivity follows from
\eqref{eq:base-change-weight} and surjectivity of $\pi$.
\end{proof}

\subsection{Measures and energies under base change}
\label{sec:ch30-bc-measures}

The measure normalization under field extension differs from that for a
degree-$d$ morphism of varieties.

\begin{theorem}\label{thm:base-change-ma}
For pulled-back model or continuous semipositive metrics,
\begin{equation}
\label{eq:base-change-ma}
 \pi_*\MA_{\theta_{K'}}(\varphi\circ\pi)=\MA_\theta(\varphi),
 \qquad
 E_{\theta_{K'}}(\varphi\circ\pi)=E_\theta(\varphi).
\end{equation}
The corresponding mixed identities also hold.
\end{theorem}

\begin{proof}
For a model metric, components of the normalized base-changed special fibre
lying over a component $E$ contribute residue degrees and ramification
multiplicities.  Their weighted sum is the original coefficient
\[
 b_E\bigl(c_1(\mathcal L)|_E\bigr)^n
\]
in the Chambert--Loir measure.  This proves the pushforward identity.  There
is no factor $[K':K]$, because the algebraic top intersection number of
the generic-fibre line bundle is invariant under scalar extension.  Apply
the mixed energy formula and then uniform approximation.

In detail, let $\varphi=\varphi_D$ be a model function determined on
$\cX$, and let $\cL_1,\ldots,\cL_n$ be model line bundles on $\cX$
representing the mixed data, with pullbacks $\cL'_i$ to the normalized
base change $\cX'$ of \Cref{lem:ch30-base-change-model}.  By
\eqref{eq:appJ-cl-measure} applied on $\cX'$, and using the normalization
\eqref{eq:ch30-uniformizer-value} for the divisorial points,
\[
 \bigwedge_{i=1}^n\bigl(\theta_{K'}+\ddc\varphi_i\circ\pi\bigr)
 =\sum_{E'}b_{E'}
   \bigl(\cL'_1|_{E'}\cdots\cL'_n|_{E'}\bigr)_{k'}\,\delta_{x_{E'}} .
\]
By \Cref{lem:ch30-base-change-model}\,\cref{itm:ch30-bcmod:points} we have
$\pi(x_{E'})=x_E$ whenever $E'$ lies over $E$, so by
\Cref{lem:ch30-pushforward-basic}\,\cref{itm:ch30-pf:support} the
pushforward groups the atoms by their images:
\[
\begin{aligned}
 \pi_*\Bigl(\bigwedge_i(\theta_{K'}+\ddc\varphi_i\circ\pi)\Bigr)
 &=\sum_E\Bigl(\sum_{E'\mapsto E}b_{E'}
    \bigl(\cL'_1|_{E'}\cdots\cL'_n|_{E'}\bigr)_{k'}\Bigr)\delta_{x_E}\\
 &=\sum_Eb_E\bigl(\cL_1|_E\cdots\cL_n|_E\bigr)_k\,\delta_{x_E},
\end{aligned}
\]
the last step being exactly
\eqref{eq:ch30-coefficient-cancellation}.  The right-hand side is
$\bigwedge_i(\theta+\ddc\varphi_i)$, which proves the mixed pushforward
identity for model data.  Taking all $\varphi_i$ equal gives the first
half of \eqref{eq:base-change-ma}.

For the energy, apply \eqref{eq:measure-pushforward} to the function
$\varphi$ in each of the $n+1$ terms of \eqref{eq:model-energy}:
\[
\begin{aligned}
 E_{\theta_{K'}}(\varphi\circ\pi)
 &=\frac1{n+1}\sum_{j=0}^n\int_{X^{\An}}\varphi\,
   \pi_*\MA_{\theta_{K'}}\bigl((\varphi\circ\pi)^{(j)},0^{(n-j)}\bigr)\\
 &=\frac1{n+1}\sum_{j=0}^n\int_{X^{\An}}\varphi\,
   \MA_\theta\bigl(\varphi^{(j)},0^{(n-j)}\bigr),
\end{aligned}
\]
which is $E_\theta(\varphi)$.  Finally, continuous semipositive metrics
are uniform limits of semipositive model metrics
(\Cref{thm:ch21-continuous-zhang}), their pullbacks converge uniformly by
\eqref{eq:base-change-weight}, mixed measures converge weakly
(\Cref{prop:ch24-continuous-uniform}\,\cref{itm:ch24-cu:unif}), pushforward
is weakly continuous
(\Cref{lem:ch30-pushforward-basic}\,\cref{itm:ch30-pf:weak}), and energy is
Lipschitz for the uniform norm (\Cref{prop:ch26-energy-lipschitz}).
\end{proof}

\begin{corollary}
\label{cor:ch30-base-change-mass}
With the above normalizations,
\begin{equation}
\label{eq:ch30-base-change-mass}
 \MA_{\theta_{K'}}(\varphi\circ\pi)\bigl(X_{K'}^{\An}\bigr)
 =\MA_\theta(\varphi)\bigl(X^{\An}\bigr)=V
 =(\{\theta\}^n)=(\{\theta_{K'}\}^n).
\end{equation}
No factor $[K':K]$ occurs anywhere.
\end{corollary}

\begin{proof}
Pushforward preserves total mass
(\Cref{lem:ch30-pushforward-basic}\,\cref{itm:ch30-pf:radon}), so the first
equality follows from \eqref{eq:base-change-ma}; the remaining ones are
\eqref{eq:appJ-cl-total-mass} together with the invariance of intersection
numbers under flat base change of the ground field.
\end{proof}

\begin{theorem}
\label{thm:ch30-base-change-psh}
Let $K'/K$ be a finite extension and let $X$ be geometrically integral.
The BFJ hypotheses hold upstairs with the extended absolute value.
Its value group is $(1/e)\mathbf Z$, rather than $\mathbf Z$;
\eqref{eq:ch30-ramified-renormalization} transports the theory to and
from the uniformizer-value-one convention without altering that
absolute value in the displayed formulae.  Then
\begin{equation}
\label{eq:ch30-base-change-psh}
 \varphi\in\PSH(X,\theta)
 \implies
 \varphi\circ\pi\in\PSH(X_{K'},\theta_{K'}),
\end{equation}
the map $\varphi\mapsto\varphi\circ\pi$ is injective, order preserving,
commutes with maxima, constants, convex combinations and decreasing nets,
and preserves boundedness, continuity, the finite-energy class and
singularity types.  Moreover
\begin{equation}
\label{eq:ch30-base-change-e1}
 E_{\theta_{K'}}(\varphi\circ\pi)=E_\theta(\varphi)
 \quad\text{in }\mathbf R\cup\{-\infty\},
 \qquad
 \varphi\in\cE^1(X,\theta)\iff
 \varphi\circ\pi\in\cE^1(X_{K'},\theta_{K'}),
\end{equation}
and \eqref{eq:base-change-ma} holds for bounded and for finite-energy
potentials, with mixed measures.
\end{theorem}

\begin{proof}
The proof is the one given for \Cref{thm:ch30-pullback-psh} and
\Cref{thm:ch30-finite-energy-transfer}, with $f^{\An}$ replaced by
$\pi$ and the degree replaced by $1$; we indicate only the two points
where the substitution has to be checked.  First, the model layer: a
$\theta$-psh model function on $X$ pulls back to a
$\theta_{K'}$-psh model function on $X_{K'}$, by
\Cref{prop:base-change-model-metrics} and
\Cref{lem:ch30-base-change-model-function}.  Second, the surjectivity used
to exclude the value $-\infty$ and to compare suprema is
\Cref{lem:ch30-base-change-map}\,\cref{itm:ch30-bcm:cont} and
\cref{itm:ch30-bcm:sup}.  Everything else --- decreasing regularization
downstairs, decreasing stability upstairs, truncation, and the monotone
convergence \eqref{eq:weighted-monotone-ma} --- is applied verbatim,
\eqref{eq:ch20-standing} being available on $X_{K'}$ because $K'$ is
again complete and discretely valued with residue characteristic zero and
$X_{K'}$ is again smooth, connected and projective.  Injectivity and the
statement about suprema and singularity types follow from surjectivity of
$\pi$; the energy identity is
\Cref{thm:base-change-ma} extended to decreasing limits by
\Cref{prop:e1-monotone-continuity}.
\end{proof}

\subsection{Intermediate fields: algebraic, non-complete, non-discrete}
\label{sec:ch30-bc-beyond}

The pullback map on psh functions is generally not onto.  Scalar extension
creates analytic points, models, and psh functions that need not descend to
$K$.  Likewise, $\pi_*\nu'$ forgets how a measure $\nu'$ is
distributed in the fibres of $\pi$.

Beyond finite extensions, three distinct phenomena have to be separated,
and it is a recurring source of confusion that all three are called
``base change''.

\emph{(i) Non-complete extensions.}  The construction $X\mapsto X_{K'}^{\An}$
requires $K'$ to be complete: the affinoid algebras out of which
$X_{K'}^{\An}$ is glued are complete by definition.  For a non-complete
valued extension $L/K$ one must first pass to the completion
$\widehat L$ and set $X_L^{\An}:=X_{\widehat L}^{\An}$.  This fixes the meaning of analytification over $L$; it does not assert
a new algebraic morphism $X_L\to X_{\widehat L}$.  The completion has
the same value group and residue field as $L$, but algebraic geometry
over the completion can change, because
$X_{\widehat L}$ has more subvarieties than $X_L$ and more models.  The
most important instance is $L=K^{a}$, an algebraic closure, which is never
complete when $K$ is discretely valued and nontrivially valued.  Its
completion $\widehat{K^{a}}$ is algebraically closed, by Krasner's lemma,
and is the standard target for arguments that wish to assume the ground
field algebraically closed.

\emph{(ii) Non-discrete value groups.}  Already for $K'=K(\varpi^{1/m})$
the value group is $(1/m)\mathbf Z$, still discrete; but the union over
all $m$ has value group $\mathbf Q$, and $\widehat{K^{a}}$ has value
group $\mathbf Q$ as well.  A rank-one valuation ring with non-discrete
value group is not Noetherian, and the notion of a model used throughout
\Cref{part:bfj} --- a normal flat \emph{projective} scheme over $R$ ---
needs a general-valuation-ring treatment or an admissible formal
model framework, for example in the sense of Raynaud
(\cite{BoschLutkebohmert1993}, \cite{RaynaudGruson1971}), or by the
algebraic models over a valuation ring studied by Gubler.  The dictionary
between the two is not automatic, and \eqref{eq:ch20-standing} explicitly
excludes it.  Consequently
\Cref{thm:ch30-base-change-psh} is stated for finite extensions only.  The
statements that do survive over an arbitrary complete extension are those
formulated for continuous semipositive metrics in Zhang's sense, which make
no reference to a discrete value group: preservation of formal/model semipositivity follows from
\cite[Proposition~2.18(a) and Lemma~3.3]{GM19}, and uniform approximation
then preserves continuous semipositivity.  The general-field measure
projection formula needs the corresponding Chambert--Loir/formal-model
measure theory in addition to these metric results; it is not proved by
the finite-discrete ramification calculation alone.

\emph{(iii) Non-algebraically-closed fields.}  Three separate effects
appear, and it is worth naming them.  The residue field $k$ may be non-algebraically closed, so components
can split after extension of their constant fields.  This is distinct
from the function-field residue degree $d_{E'}$, which can exceed one
even when $k$ is algebraically closed, as the conic example shows.
The factor $1/f$ in \eqref{eq:ch30-degree-transfer} changes the field
over which numerical degrees are taken; it does not simply remove all
component residue degrees.  The
variety $X$ may fail to be geometrically integral, which is
\Cref{rem:ch30-not-geometrically-integral}.  And the fibres of $\pi$ may
have several points even over a divisorial point, which is why
\Cref{lem:ch30-base-change-model}\,\cref{itm:ch30-bcmod:points} asserts only
that $\pi(x_{E'})=x_E$ and not that $E'$ is unique over $E$.

\begin{remark}[If $X$ is not geometrically integral]
\label{rem:ch30-not-geometrically-integral}
Suppose $X$ is smooth, connected and projective over $K$, but that the
algebraic closure $K_1$ of $K$ in $K(X)$ is a proper extension of
$K$, necessarily finite and separable.  Then for $K'\supseteq K_1$ the
finite $K'$-algebra
\[
                         K_1\otimes_KK'=\prod_\nu K'_\nu
\]
is a product of finite separable field extensions, and $X_{K'}$ decomposes
accordingly into integral components after base change to the factors
$K'_\nu$.  If $K'$ contains a normal closure of $K_1/K$, every factor
is $K'$ and there are exactly $[K_1:K]$ components; the weaker
assumption $K'\supseteq K_1$ alone does not imply this complete splitting.
The analytification is the corresponding disjoint union.  All statements
above remain true component by component,
with \eqref{eq:ch30-ramification-sum} replaced by the corresponding sum
over the components of $X_{K'}$, and with the total mass
\eqref{eq:ch30-base-change-mass} distributed among them.  The one statement
that genuinely fails is connectedness, hence the assertion in
\eqref{eq:ch20-standing} that $X_{K'}$ is again admissible; this is why
geometric integrality was imposed in \Cref{sec:ch30-hypotheses} rather than
carried as a case distinction.
\end{remark}

\begin{remark}[The psh class can grow]
\label{rem:ch30-psh-grows}
Take $X=\mathbf P^1_K$, the standard reference on $\cO(1)$, and
$K'=K(a)$ with $a^2=\varpi$.  The two $K'$-rational points $z=a$
and $z=-a$ have the same image downstairs.  The model potential
\[
 \psi=-\min\{1,\max\{0,v(z-a)\}\}
\]
is $\theta_{K'}$-psh: its curve curvature moves one unit of mass
from the Gauss point to the divisorial point of the disc
$|z-a|\leq\exp(-1)$.  Its values at $z=a$ and $z=-a$ are $-1$
and $-1/2$, since $v_{K'}(2a)=1/2$.
It is not constant on the fibre and therefore cannot be a pullback.
Thus strict growth is possible, without a universal claim for all
extensions.  The conic example below illustrates changed multiplicities
of an existing point; that fact alone does not prove growth of the
psh class.
\end{remark}

\subsection{A ramified base change, carried to numbers}
\label{sec:ch30-bc-example}

\begin{example}[The standard conic over a ramified quadratic extension]
\label{ex:ch30-conic-base-change}
\index{base change!worked example}
Let $K$ satisfy \eqref{eq:ch20-standing} with $v_K(\varpi)=1$, and take
the model
\[
 \mathcal C=\{XY=\varpi Z^2\}\subset\mathbf P^2_R
\]
of \eqref{eq:running-conic}, with generic fibre a smooth conic
$C\simeq\mathbf P^1_K$ and $L=\cO_{\mathbf P^2}(1)|_C$ of degree
$V=\deg L=2$.  In the affine chart $Z=1$ write $x=X$, $y=Y$, so
that $xy=\varpi$; the model $\mathcal C$ is regular, its special fibre
$C_0\cup C_1$ is a pair of lines meeting at a node, both of multiplicity
one, and its skeleton is the segment
\[
 \Gamma=\{v_s:0\leq s\leq1\},
 \qquad v_s(x)=s,\quad v_s(y)=1-s,
\]
of length $v_K(\varpi)=1$.  Its Chambert--Loir measure is
$\mu_0=\delta_{v_0}+\delta_{v_1}$ of mass $2$, by
\Cref{prop:running-conic-measure} and \eqref{eq:running-conic-measure}.

\emph{A potential downstairs.}  Blow up the node of $\mathcal C$ --- a regular
point of the regular surface $\mathcal C$ --- to obtain $\widetilde{\mathcal C}$ with
special fibre $C_0+2E+C_1$, where $\ord_E(x)=\ord_E(y)=1$, hence
$b_E=\ord_E(\varpi)=2$ and
\[
 x_E=b_E^{-1}\ord_E=v_{1/2}.
\]
Since $\widetilde{\mathcal C}_s$ is principal, $(\widetilde{\mathcal C}_s\cdot E)=0$
gives $(E^2)=-1$.  Let $\varphi=\varphi_D$ with $D=-E$; its vertex
values are
\[
 \varphi(v_0)=\varphi(v_1)=0,\qquad
 \varphi(v_{1/2})=\frac{-1}{b_E}=-\frac12,
\]
so $\varphi$ is the tent function on $\Gamma$ with depth $-\tfrac12$
at the midpoint.  Writing $\cL$ for the pullback of
$\cO_{\mathbf P^2_R}(1)$, the degrees are $(\cL\cdot C_i)=1$,
$(\cL\cdot E)=0$, $(D\cdot C_i)=-1$, $(D\cdot E)=-(E^2)=1$, whence
\begin{equation}
\label{eq:ch30-conic-ma-down}
 \MA_\theta(\varphi)
 =\sum_Fb_F\bigl((\cL+D)\cdot F\bigr)\delta_{x_F}
 =0\cdot\delta_{v_0}+2\cdot1\cdot\delta_{v_{1/2}}+0\cdot\delta_{v_1}
 =2\,\delta_{v_{1/2}},
\end{equation}
of mass $2=V$.  All three coefficients are nonnegative, so
$\varphi$ is $\theta$-psh.  The same answer comes from the graph
Laplacian \eqref{eq:curve-ma-laplacian}: the outgoing slopes of the tent
are $2\varphi(v_{1/2})=-1$ at each endpoint and $+2$ at the midpoint,
so $\MA_\theta(\varphi)=\mu_0+\Delta_\Gamma\varphi
=(1-1)\delta_{v_0}+(1-1)\delta_{v_1}+2\delta_{v_{1/2}}$.  Finally
\begin{equation}
\label{eq:ch30-conic-energy-down}
 E_\theta(\varphi)
 =\tfrac12\Bigl(\int\varphi\,\MA_\theta(\varphi)
   +\int\varphi\,\MA_\theta(0)\Bigr)
 =\tfrac12\Bigl(2\cdot\bigl(-\tfrac12\bigr)+0\Bigr)=-\frac12 .
\end{equation}

\emph{The extension.}  Let $K'=K(\varpi')$ with $\varpi'^2=\varpi$, so
that $e=2$, $k'=k$, $f=1$, $[K':K]=2$, and, in the normalization
\eqref{eq:ch30-uniformizer-value}, $v_{K'}(\varpi')=\tfrac12$.  The base
change $\mathcal C\times_RR'=\{XY=\varpi'^2Z^2\}$ has an $A_1$ singularity at
the node; it is already normal, and its special fibre is the pair of lines
$L_0\cup L_1$ with $b'_{L_i}=\ord_{L_i}(\varpi')=1$.  Blowing up the
singular point yields the semistable model $\mathcal C''$ over $R'$ with
\[
 \mathcal C''_s=L_0+E''+L_1,
 \qquad b'_{L_0}=b'_{E''}=b'_{L_1}=1,
\]
where $E''$ is a $(-2)$-curve, as $(\mathcal C''_s\cdot E'')=0$ shows.  The
associated divisorial points, normalized as in
\Cref{lem:ch30-base-change-model}\,\cref{itm:ch30-bcmod:points}, are
\[
 x_{L_0}=v'_0,\qquad
 x_{E''}=\bigl(\ord_{E''}(\varpi)\bigr)^{-1}\ord_{E''}
        =\tfrac12\ord_{E''}=v'_{1/2},\qquad
 x_{L_1}=v'_1 ,
\]
whose images under $\pi$ are $v_0,v_{1/2},v_1$, respectively.
The prime distinguishes a valuation on the scalar-extended function
field from its restriction to the original field.

\emph{The bookkeeping numbers.}  The model $\mathcal C''$ is the normalized base change of
$\widetilde{\mathcal C}$, with a finite morphism
$\rho:\mathcal C''\to\widetilde{\mathcal C}$.
Indeed, the node blowup downstairs has a chart
$y=ux$, $\varpi=ux^2$; after adjoining $\varpi'$ its normalization
has $w=\varpi'/x$, $u=w^2$, and $xw=\varpi'$.
The symmetric chart uses $w'=\varpi'/y$, with $ww'=1$ on the
exceptional overlap.  These are exactly the two regular charts of
the resolution of $xy=\varpi'^2$.
Over $E$ there is the single component $E''$, and
\eqref{eq:ch30-multiplicity-relation} gives
$\ord_{E''}(\varpi)=b_Em_{E''}=e\,b'_{E''}$, that is,
$2m_{E''}=2$, so $m_{E''}=1$; then
\eqref{eq:ch30-ramification-sum} forces
\[
 d_{E''}=[\kappa(E''):\kappa(E)]=\frac{ef}{m_{E''}}=2 .
\]
This is visible: $\varpi/x^2$ is a unit on $\widetilde{\mathcal C}$ whose
restriction to $E$ is a coordinate $u$ of $E\simeq\mathbf P^1_k$, and
adjoining $\varpi'$ adjoins $\sqrt u$; thus $E''\to E$ is the
degree-two cover $u=(u')^2$, again a projective line, branched at the two
nodes.  Equivalently $\mathcal H(v_{1/2})\widehat\otimes_KK'$ is a field
of degree two over $\mathcal H(v_{1/2})$, so
$\pi^{-1}(v_{1/2})$ is a \emph{single} point, with residue degree two.
Over $C_0$ the roles are exchanged: there is one component $L_0$, with
$\ord_{L_0}(\varpi)=b_{C_0}m_{L_0}=e\,b'_{L_0}=2$, so $m_{L_0}=2$ and
$d_{L_0}=1$.  The cancellation
\eqref{eq:ch30-coefficient-cancellation} reads, over $C_0$,
\[
 b'_{L_0}\bigl(\cL'|_{L_0}\bigr)_{k'}
 =1\cdot\frac{d_{L_0}}{f}\bigl(\cL|_{C_0}\bigr)_k
 =1\cdot1\cdot1=1=b_{C_0}\bigl(\cL|_{C_0}\bigr)_k ,
\]
and over $E$ both sides vanish because $(\cL|_E)=0$.

\emph{The skeleton.}  The skeleton $\Gamma''\subset C_{K'}^{\An}$ of
$\mathcal C''$ is again the segment of monomial points $v'_s$ with
$v'_s(x)=s$, $0\leq s\leq1$, and $\pi$ maps it bijectively and
isometrically onto $\Gamma$; its length is $v_{K'}(\varpi)=1$,
\emph{not} $2$.  What has changed is the multiplicity data: the point
$v_{1/2}$, divisorial of multiplicity $b=2$ over $K$, is divisorial of
multiplicity $b'=1$ over $K'$, and $\mathcal C''$ is semistable over $R'$.
This is the sense in which base change simplifies models without moving the
skeleton.

\emph{The measure and the energy.}  On $\mathcal C''$, with
$D''=\rho^*D=-E''$ and $\cL''$ the pullback of $\cL$, one has
$(\cL''\cdot L_i)=1$, $(\cL''\cdot E'')=0$, $(D''\cdot L_i)=-1$, and
$(D''\cdot E'')=-\bigl((E'')^2\bigr)=2$.  Hence
\begin{equation}
\label{eq:ch30-conic-ma-up}
 \MA_{\theta_{K'}}(\varphi\circ\pi)
 =0\cdot\delta_{v'_0}+1\cdot 2\cdot\delta_{v'_{1/2}}+0\cdot\delta_{v'_1}
 =2\,\delta_{v'_{1/2}},
\end{equation}
again of mass $2$, and $\pi_*$ of it is $2\delta_{v_{1/2}}$, which is
\eqref{eq:ch30-conic-ma-down}: this is \eqref{eq:base-change-ma}, with no
factor $[K':K]=2$.  Note that the coefficient $2$ has changed its
origin: downstairs it was $b_E\cdot1$, upstairs it is
$b'_{E''}\cdot2$.  The vertex value is unchanged,
$\varphi_{D''}(v'_{1/2})=(\ord_{E''}(\varpi))^{-1}\cdot(-1)=-\tfrac12$,
by \eqref{eq:ch30-model-function-vertices}, and therefore
\begin{equation}
\label{eq:ch30-conic-energy-up}
 E_{\theta_{K'}}(\varphi\circ\pi)
 =\tfrac12\Bigl(2\cdot\bigl(-\tfrac12\bigr)+0\Bigr)=-\frac12
 =E_\theta(\varphi).
\end{equation}

\emph{The same computation in the wrong normalization.}  Suppose one
insists that a uniformizer of $K'$ have valuation one, that is, replaces
$v_{K'}$ by $v''=2v_{K'}$.  Then the skeleton becomes a segment of
length $2$, the divisorial point formerly called $v'_{1/2}$ becomes
$v''_1$, the model function becomes the tent of depth $-1$, and
\[
 \MA''=2\,\delta_{v''_1},
 \qquad
 E''=\tfrac12\bigl(2\cdot(-1)+0\bigr)=-1=2\,E_\theta(\varphi).
\]
The measure is the same measure with relabelled atoms; the weight and the
energy have been multiplied by $e=2$.  This is exactly
\eqref{eq:ch30-ramified-renormalization}, and it is the reason the
checkpoint above is stated as an equation and not as a caution.
\end{example}

\section{Descent and uniqueness}
\label{sec:ch30-descent}

Suppose $K'/K$ is finite Galois with group $G$.  The group acts on
$X_{K'}^{\An}$ and on pulled-back metrics.

\subsection{Three different descent questions}
\label{sec:ch30-descent-questions}

The word \emph{descent} is used for three questions which have different
answers, and the difference matters.

\begin{description}[style=nextline,leftmargin=0pt,labelindent=0pt,
  labelsep=0pt,itemsep=.5\baselineskip,font=\normalfont\itshape]
\item[(Q1) Functions.]
Is every $G$-invariant function on $X_{K'}^{\An}$ of the form
$u\circ\pi$?  This is a question about the topological quotient, and it is
completely settled by \Cref{thm:ch30-galois-quotient}: the answer is yes,
with $u$ unique and inheriting continuity or upper semicontinuity.

\item[(Q2) Positivity.]
If $u$ is a function on $X^{\An}$ with $u\circ\pi$ in
$\PSH(X_{K'},\theta_{K'})$, is $u\in\PSH(X,\theta)$?  This is a genuine
question, because $\PSH$ is defined by an approximation property relative
to the models available over the base field, and there are more models
upstairs.  We prove it for model functions already defined downstairs
(\Cref{thm:ch30-model-psh-descent}) and, conditionally on the model
descent property, for arbitrary psh functions
(\Cref{thm:ch30-continuous-descent}).  The model descent property is
proved here for unramified extensions (\Cref{thm:ch30-model-descent})
and remains an explicit input otherwise.

\item[(Q3) Solutions.]
If $\psi$ solves a Monge--Amp\`ere equation upstairs, does it descend to a
solution downstairs?  Uniqueness first forces $G$-invariance; descending
that invariant potential still uses the positivity answer to (Q2).
\Cref{prop:galois-descent-potentials} keeps these two steps separate.
\end{description}

\subsection{The analytification is the quotient}
\label{sec:ch30-galois-quotient}

\begin{theorem}
\label{thm:ch30-galois-quotient}
\index{Galois action!on Berkovich spaces}\index{descent!Galois}
Let $K'/K$ be a finite Galois extension with group $G$ and let $X$ be
a proper $K$-scheme.  Then $G$ acts continuously on $X_{K'}^{\An}$,
the map $\pi$ is $G$-invariant, and the induced map
\begin{equation}
\label{eq:ch30-galois-quotient}
 X_{K'}^{\An}/G\longrightarrow X^{\An}
\end{equation}
is a homeomorphism.  In particular the fibres of $\pi$ are exactly the
$G$-orbits, and they are finite.
\end{theorem}

\begin{proof}
\emph{Step 1: the action.}  Each $g\in G$ is a $K$-automorphism of
$K'$, hence induces a $K$-automorphism of the scheme
$X_{K'}=X\times_KK'$ over $X$, and therefore a homeomorphism of
$X_{K'}^{\An}$ commuting with $\pi$
(\Cref{sec:berk-functoriality}).  Continuity of the action of the finite
group $G$ is automatic.

\emph{Step 2: the fibres are finite and are single orbits.}  Fix
$x\in X^{\An}$.  By
\Cref{lem:ch30-base-change-map}\,\cref{itm:ch30-bcm:fibre},
\[
 \pi^{-1}(x)\simeq
 \mathcal M\bigl(\mathcal H(x)\widehat\otimes_KK'\bigr).
\]
Since $[K':K]<\infty$, the algebra $A:=\mathcal H(x)\otimes_KK'$ is a
finite-dimensional $\mathcal H(x)$-algebra, hence already complete, so no
completion is needed.  As $\operatorname{char}K=0$, the extension
$K'/K$ is separable and $A$ is a finite étale
$\mathcal H(x)$-algebra; write
$A\simeq L_1\times\cdots\times L_r$ with $L_i$ finite separable field
extensions of $\mathcal H(x)$.  Each $L_i$ is a finite extension of the
complete field $\mathcal H(x)$, so it carries a unique absolute value
extending that of $\mathcal H(x)$, and its Berkovich spectrum is a single
point.  A multiplicative seminorm on a finite product of fields is supported
on one factor, so
\[
 \pi^{-1}(x)=\{y_1,\ldots,y_r\},
 \qquad \mathcal H(y_i)=L_i .
\]

It remains to see that $G$ permutes the $L_i$ transitively.  Fix an
algebraic closure $\Omega$ of $\mathcal H(x)$.  The factors $L_i$ of
$A=\mathcal H(x)\otimes_KK'$ correspond bijectively to the orbits of
$\Gamma:=\Aut(\Omega/\mathcal H(x))$ acting on the set
$S:=\operatorname{Hom}_K(K',\Omega)$, by the standard description of a
finite étale algebra.  Now $\#S=[K':K]=\#G$, because $K'/K$ is
separable, and $G$ acts on $S$ by $g\cdot\iota=\iota\circ g^{-1}$,
freely (as $\iota$ is injective) and hence simply transitively.  The two
actions on $S$ commute, so the $G$-action descends to the orbit set
$S/\Gamma$ and is transitive there, being transitive on $S$.  Hence
$G$ permutes $\{L_1,\ldots,L_r\}$ transitively, that is,
$\pi^{-1}(x)$ is one $G$-orbit.  (The schematic counterpart of this
transitivity is \cite[Tag 0BRI]{stacks-project}, recorded in
\Cref{app:stacks}.)

\emph{Step 3: conclusion.}  The map \eqref{eq:ch30-galois-quotient} is a
continuous bijection, by Step~2 and surjectivity of $\pi$
(\Cref{lem:ch30-base-change-map}\,\cref{itm:ch30-bcm:cont}).  Its source is a
quotient of a compact space, hence compact; its target is Hausdorff.  A
continuous bijection from a compact space to a Hausdorff space is a
homeomorphism.
\end{proof}

\begin{corollary}
\label{cor:ch30-invariant-functions}
Let $K'/K$ be finite Galois with group $G$.  The map
$u\mapsto u\circ\pi$ is a bijection from the set of functions
$X^{\An}\to[-\infty,+\infty]$ onto the set of $G$-invariant such
functions on $X_{K'}^{\An}$.  It preserves and reflects continuity, upper
semicontinuity, boundedness, suprema, infima, maxima, pointwise order,
and sums whenever the extended-real sums are defined.

Moreover, pushforward is a bijection from $G$-invariant positive Radon
measures on $X_{K'}^{\An}$ to positive Radon measures on $X^{\An}$.
Its inverse is the mass-preserving invariant lift $\mu\mapsto\mu^G$,
characterized by
\begin{equation}
\label{eq:ch30-invariant-measure-lift}
 \int_{X_{K'}^{\An}}h\,\rd\mu^G
 =\int_{X^{\An}}\overline h\,\rd\mu,\qquad
 \overline h\circ\pi=\frac1{\#G}\sum_{g\in G}h\circ g,
 \quad h\in C^0(X_{K'}^{\An}).
\end{equation}
\end{corollary}

\begin{proof}
Injectivity is surjectivity of $\pi$.  Given a $G$-invariant $w$ on
$X_{K'}^{\An}$, it factors through the set-theoretic quotient, which is
$X^{\An}$ by \Cref{thm:ch30-galois-quotient}; the factorization
$u$ is unique.  Because \eqref{eq:ch30-galois-quotient} is a
homeomorphism, $\pi$ is a topological quotient map, so $u$ is continuous
or upper semicontinuous exactly when $w$ is.  The remaining assertions
about functions are pointwise identities.

For measures, the average on the right of
\eqref{eq:ch30-invariant-measure-lift} is invariant and continuous, hence
descends to a unique continuous $\overline h$ by the quotient theorem.
The functional $h\mapsto\int\overline h\,\rd\mu$ is positive and
bounded, so Riesz representation gives a Radon measure $\mu^G$.
Averaging is unchanged by the group action; hence this measure is
invariant.  For $h=u\circ\pi$ the average is $u\circ\pi$, so
$\pi_*\mu^G=\mu$, including equality of total masses.
Conversely, if $\nu$ is invariant and $\pi_*\nu=\mu$, averaging the
integral of $h$ gives
$\int h\,\rd\nu=\int\overline h\circ\pi\,\rd\nu
=\int\overline h\,\rd\mu$.  Thus $\nu=\mu^G$.
No choice of individual points in a fibre is involved.
\end{proof}

\subsection{Descent of positivity}
\label{sec:ch30-descent-positivity}

\begin{theorem}
\label{thm:ch30-model-psh-descent}
Let $K'/K$ be a finite extension --- Galois or not --- and let
$\varphi\in\mathcal D(X)_{\mathbf R}$ be a model function on $X$.  Then
\[
 \varphi\ \text{is}\ \theta\text{-psh}
 \iff
 \varphi\circ\pi\ \text{is}\ \theta_{K'}\text{-psh}.
\]
More generally, for a model line bundle $\cM$ on a model $\cX$ of
$X$, $\cM$ is relatively nef if and only if its pullback to the
normalized base change $\cX'$ is.
\end{theorem}

\begin{proof}
Write $\varphi=\varphi_D$ on a model $\cX$, so that
$\theta+\ddc\varphi=[\cM]$ with $\cM=\theta_{\cX}+D$ a real model class, and let
$\rho:\cX'\to\cX$ be as in \Cref{lem:ch30-base-change-model}, a finite
surjective morphism.  By
\Cref{lem:ch30-base-change-model-function}, $\varphi\circ\pi$ is the
model function of $\rho^*D$ on $\cX'$ and
$\theta_{K'}+\ddc(\varphi\circ\pi)=[\rho^*\cM]$.  Now apply
\Cref{lem:ch30-semipositivity-descends}, which is precisely the assertion
that relative nefness of $\cM$ and of $\rho^*\cM$ are equivalent for a
finite surjective $\rho$.  Semipositivity of a model form is independent
of the model on which it is computed
(\Cref{prop:ch20-semipositive-well-defined}), so the conclusion does not
depend on the choices.
\end{proof}

The corresponding statement for arbitrary potentials is not formal, because
the class $\PSH$ is defined by approximation and the supply of
approximants grows with the ground field
(\Cref{rem:ch30-psh-grows}).  What is needed is the descent of the
approximants themselves.

\begin{definition}[The model descent property]
\label{def:ch30-model-descent-property}
Let $K'/K$ be finite Galois with group $G$.  Say that $K'/K$ has the
\emph{model descent property} if
\begin{equation}
\label{eq:ch30-model-descent-property}
 \pi^*:\mathcal D(X)_{\mathbf Q}
 \longrightarrow
 \mathcal D(X_{K'})_{\mathbf Q}^{\,G}
\end{equation}
is surjective, that is, if every $G$-invariant model function on
$X_{K'}$ is the pullback of a model function on $X$.
\end{definition}

The map \eqref{eq:ch30-model-descent-property} is always injective, by
\Cref{cor:ch30-invariant-functions}, and its image consists of model
functions, so the only issue is surjectivity.

\begin{theorem}
\label{thm:ch30-model-descent}
If $K'/K$ is finite Galois and unramified, then it has the model descent
property.
\end{theorem}

\begin{proof}
Since $K'/K$ is unramified, $R'$ is a finite étale $R$-algebra and
$\Spec R'\to\Spec R$ is a Galois covering with group $G$; in
particular $R'\otimes_RR'\simeq\prod_{g\in G}R'$, so a semilinear
$G$-action on an $R'$-scheme or on a quasi-coherent sheaf is the same
thing as a descent datum for the faithfully flat morphism
$\Spec R'\to\Spec R$ \cite[Exp.~VIII]{SGA1}.

Let $\psi\in\mathcal D(X_{K'})_{\mathbf Q}^G$.  Choose a normal model
$\cX_1$ of $X_{K'}$ and a vertical $\mathbf Q$-Cartier divisor
$D_1$ with $\psi=\varphi_{D_1}$.  For $g\in G$ the translate
$g(\cX_1)$ is again a normal model, so the normalization $\cX'$ of the
closure of the diagonal image of $X_{K'}$ in
$\prod_{g\in G}g(\cX_1)$ is a normal model of $X_{K'}$ dominating all
of them, and it is $G$-stable by construction; thus $G$ acts on
$\cX'$, semilinearly over $R'$.  On $\cX'$, $\psi=\varphi_{D'}$ for
a unique vertical $\mathbf Q$-Cartier divisor $D'$ --- uniqueness
because $D\mapsto\varphi_D$ is injective on a fixed model
(\Cref{lem:metrics-model-lattice}) --- and $D'$ is $G$-invariant,
because $\psi$ is.

Choose a relatively ample line bundle $A$ on $\cX'$ and tensor its
translates: $\bigotimes_{g\in G}g^*A$ is relatively ample and has the
canonical $G$-linearization obtained by permuting the factors.
Effective descent with this ample datum
\cite[Exp.~VIII, Proposition~7.8 and Corollary~5.8]{SGA1}, together with
$R'\otimes_RR'\simeq\prod_GR'$, gives a projective $R$-scheme $\cX$
with $\cX\times_RR'\simeq\cX'$.
Its generic fibre is $X$; flatness and normality descend under this
finite faithfully flat étale cover, so it is a normal model.

Clear a denominator $q$ of $D'$.  The invariant Cartier divisor $qD'$
defines an invariant invertible subsheaf
$\cO_{\cX'}(-qD')$ of the sheaf of rational functions.  Its inclusion,
as well as the invertible sheaf, is a descent datum; faithfully flat
descent \cite[Exp.~VIII, Theorem~1.1]{SGA1} therefore gives a Cartier divisor $D_0$ on $\cX$ whose pullback
is $qD'$.  It is vertical since its pullback has no generic-fibre
support.  Put $D=D_0/q$.  Finally $\varphi_D\circ\pi=\varphi_{\rho^*D}=\varphi_{D'}=\psi$
by \Cref{lem:ch30-base-change-model-function}, so $\psi$ is in the image
of \eqref{eq:ch30-model-descent-property}.
\end{proof}

\begin{theorem}
\label{thm:ch30-continuous-descent}
Let $K'/K$ be finite Galois with group $G$ and assume the model descent
property.  Let
$u:X^{\An}\to[-\infty,+\infty)$ be not identically $-\infty$.  If
$u\circ\pi\in\PSH(X_{K'},\theta_{K'})$, then
$u\in\PSH(X,\theta)$.  In particular the assertion holds for continuous
potentials.
\end{theorem}

\begin{proof}
\emph{Step 1: invariant approximants upstairs.}  The function
$u\circ\pi$ is $\theta_{K'}$-psh, so
\Cref{thm:bfj-decreasing-regularization} gives a decreasing net of
$\theta_{K'}$-psh model functions
$\psi_\alpha\downarrow u\circ\pi$.  Put
\[
 \overline\psi_\alpha:=\frac1{\#G}\sum_{g\in G}g^*\psi_\alpha .
\]
Each $g^*\psi_\alpha$ is again a $\theta_{K'}$-psh model function: $g$
carries models over $R'$ to models over $R'$ and vertical divisors to
vertical divisors, and $g^*\theta_{K'}=\theta_{K'}$ because
$\theta_{K'}$ is pulled back from $X$.  Hence
$\overline\psi_\alpha$ is a
model function, and it is $\theta_{K'}$-psh as a convex combination
(\Cref{thm:ch21-stability}\,\cref{itm:ch21-stab:sum}).  It is
$G$-invariant.  The net is decreasing, and its pointwise limit is
$\#G^{-1}\sum_g g^*(u\circ\pi)=u\circ\pi$.

\emph{Step 2: descent of the approximants.}  By the model descent property
there are rational model functions
$\chi_\alpha$ with $\chi_\alpha\circ\pi=\overline\psi_\alpha$.
Here the regularization theorem supplies rational model functions
$\psi_\alpha\in\mathcal D(X)$, and the finite average stays rational.
More generally the descent property extends to real model functions:
write an invariant real model function as a finite real linear
combination of rational ones and average each summand over $G$;
each invariant rational summand descends.  By
\Cref{thm:ch30-model-psh-descent}, each $\chi_\alpha$ is $\theta$-psh.
Surjectivity of $\pi$ reflects pointwise order and pointwise convergence,
so $\chi_\alpha$ is decreasing and $\chi_\alpha\downarrow u$.

\emph{Step 3: conclusion.}  Decreasing stability
\Cref{thm:ch21-decreasing-stability} gives
$u\in\PSH(X,\theta)$.
\end{proof}

\subsection{Descent of solutions}
\label{sec:ch30-descent-solutions}

\begin{proposition}\label{prop:galois-descent-potentials}
Let $\theta'=\pi^*\theta$, and assume a continuous semipositive solution
$\psi$ of
$\MA_{\theta'}(\psi)=\mu'$ is unique after the normalization
$\sup\psi=0$.  If $\mu'$ is $G$-invariant, then $\psi$ is
$G$-invariant.  If, in addition, the model descent property holds (or psh descent for
this potential has been established), then $\psi=\varphi\circ\pi$
for a continuous $\theta$-psh potential $\varphi$ satisfying
\[
                       \MA_\theta(\varphi)=\pi_*\mu'.
\]
\end{proposition}

\begin{proof}
For $g\in G$, functoriality gives
$\MA_{\theta'}(g^*\psi)=g^*\mu'=\mu'$, and the supremum normalization is
unchanged.  Uniqueness gives $g^*\psi=\psi$.  The stated positivity descent input gives a psh potential downstairs.
Formula \eqref{eq:base-change-ma} verifies
the descended equation after pushforward.

We expand the two steps.  For the first, $g$ is an automorphism of
$X_{K'}$ over $K$ fixing $\theta_{K'}$, so $g^*\psi$ is again
continuous and $\theta_{K'}$-psh, and the functoriality of the model
Monge--Amp\`ere measure under the isomorphism $g$ --- obtained because the semilinear map $g$ permutes models over
$R'$, their components and vertical intersections while acting on the
residue field by an automorphism --- gives
$\MA_{\theta_{K'}}(g^*\psi)=g^*\MA_{\theta_{K'}}(\psi)=g^*\mu'=\mu'$.
Since $\sup g^*\psi=\sup\psi=0$, uniqueness forces $g^*\psi=\psi$.

For the second, $\psi$ is now a $G$-invariant continuous function, so
$\psi=\varphi\circ\pi$ for a unique continuous $\varphi$ on $X^{\An}$
by \Cref{cor:ch30-invariant-functions}.  That $\varphi$ is
$\theta$-psh is exactly \Cref{thm:ch30-continuous-descent}, and is the
additional positivity descent input in the statement; descent of an
underlying line bundle alone is not a proof of this positivity.
Granting it, \eqref{eq:base-change-ma} gives
$\MA_\theta(\varphi)=\pi_*\MA_{\theta_{K'}}(\varphi\circ\pi)
=\pi_*\mu'$.
\end{proof}

\begin{corollary}
\label{cor:ch30-descent-cy}
Let $K'/K$ be finite Galois with the model descent property, let $\mu$
be a positive Radon measure on $X^{\An}$ of mass $V$, and let $\mu'$
be a $G$-invariant positive Radon measure on $X_{K'}^{\An}$ of mass
$V$ with $\pi_*\mu'=\mu$.  If the equation
$\MA_{\theta_{K'}}(\psi)=\mu'$ has a continuous solution which is unique
modulo constants, then $\MA_\theta(\varphi)=\mu$ has a continuous
solution, unique modulo constants.
\end{corollary}

\begin{proof}
Existence is \Cref{prop:galois-descent-potentials} together with
\Cref{thm:ch30-continuous-descent}.  For uniqueness, if $\varphi_i$ are two continuous solutions downstairs,
their pulled-back Monge--Amp\`ere measures $\nu_i$ are
$G$-invariant by the semilinear model argument above, and
$\pi_*\nu_i=\mu$.  \Cref{cor:ch30-invariant-functions} therefore gives
$\nu_1=\nu_2=\mu'=\mu^G$.  Uniqueness upstairs makes the two pullbacks
differ by a constant, and surjectivity of $\pi$ gives the same conclusion
downstairs.  One can also use \Cref{thm:ch29-uniqueness} directly.
The invariant-measure condition is essential to this argument:
unrestricted pushforward is not injective.
\end{proof}

\subsection{The limits of what is proved here}
\label{sec:ch30-descent-limits}

The last hypothesis in \Cref{prop:galois-descent-potentials} is
deliberately explicit.  Invariance yields a descended continuous function through the quotient;
its psh property still requires the specified positivity descent input.
For line bundles and models, numerical invariance alone is not a
faithfully flat descent datum.  Conversely,
if $\varphi$ is already a solution downstairs, then
$\varphi\circ\pi$ has Monge--Amp\`ere measure
$\MA_{\theta_{K'}}(\varphi\circ\pi)$, whose pushforward is the original
target by \eqref{eq:base-change-ma}.  In the finite Galois case this upstairs measure is precisely $\mu^G$,
by its invariance and \Cref{cor:ch30-invariant-functions}.  That
selected lift is distinct from inverting unrestricted pushforward.

Three boundaries should be recorded precisely.

\emph{Ramified model descent is not proved here.}  If $K'/K$ is
ramified, then $\Spec R'\to\Spec R$ is not étale, the isomorphism
$R'\otimes_RR'\simeq\prod_GR'$ fails, and a semilinear $G$-action on a
model over $R'$ is not a descent datum.  The proof of
\Cref{thm:ch30-model-descent} therefore breaks down at its first line.  We
do not replace this missing descent datum by skeletal compatibility:
descending a $G$-invariant model function requires an actual model and a
Cartier divisor downstairs.  Constructing them through the ramified finite
quotient is an additional step not supplied here.  Thus
\Cref{thm:ch30-continuous-descent} remains conditional on the model descent
property in the ramified case.  Once that property is granted, its
decreasing-net proof already covers singular as well as continuous
potentials.

\emph{Infinite and non-Galois extensions.}  For a Galois algebraic
extension $L/K$, the finite Galois subextensions give compatible
finite-level quotient statements.  Passing from those statements to a
quotient description after completion of $L$ requires an inverse-limit
comparison and is not asserted here; moreover the class $\PSH$ over
$\widehat L$ is not covered by \eqref{eq:ch20-standing}.  Thus the
statement of (Q2) first needs a new ambient theory.  For a non-Galois
finite extension one may pass to its Galois closure when formulating a
descent problem, which is harmless in characteristic zero.

\emph{Arbitrary complete extensions.}  Whether semipositivity of a
continuous metric descends along an arbitrary complete extension
$K'/K$ is a question about the local characterization of semipositivity
rather than about Galois theory.  The relevant comparison is
\Cref{thm:gm-metric-comparison} and the work of Gubler and Martin
\cite{GM19}; that route is not developed in this book, and no statement of
that generality is used anywhere in \Cref{part:bfj}.

\section{Pushforward and the general projection formula}
\label{sec:ch30-pushforward}

\subsection{One formula for all three columns}
\label{sec:ch30-general-projection}

\Cref{thm:finite-ma-projection},
\Cref{prop:birational-model-projection} and
\Cref{thm:base-change-ma} are three instances of a single statement, and it
is worth recording it once in the generality in which it is proved, if only
because the shape of the hypotheses is then visible.

\begin{theorem}
\label{thm:ch30-general-projection}
\index{projection formula!general form}
Let $g:Z\to W$ be one of the following:
\begin{enumerate}
\item a surjective, generically finite morphism of degree $d$ between
smooth connected projective $K$-varieties of the same dimension $n$
(so $d=1$ if $g$ is birational); or
\item the projection $\pi:X_{K'}^{\An}\to X^{\An}$ attached to a finite
extension $K'/K$, with $d=1$.
\end{enumerate}
Let $\theta$ be a closed model form on the target with ample class,
$\theta_Z=g^*\theta$ its pullback, and let
$\varphi_1,\ldots,\varphi_n$ be model or continuous $\theta$-psh
functions on the target.  Then
\begin{equation}
\label{eq:ch30-general-projection}
 g_*\MA_{\theta_Z}\bigl(g^*\varphi_1,\ldots,g^*\varphi_n\bigr)
 =d\,\MA_\theta(\varphi_1,\ldots,\varphi_n),
 \qquad
 E_{\theta_Z}(g^*\varphi)=d\,E_\theta(\varphi),
\end{equation}
and the total masses are $dV$ upstairs and $V$ downstairs.  If moreover
$\{\theta_Z\}$ is ample --- which holds in case (1) when $g$ is finite,
and in case (2) always, but fails for a nontrivial birational $g$ ---
then \eqref{eq:ch30-general-projection} extends to bounded and to
finite-energy potentials, with mixed measures.
\end{theorem}

\begin{proof}
The model and continuous cases of (1) are
\Cref{thm:finite-ma-projection} and
\Cref{cor:finite-pullback-energy}; the proofs there used only that
$\tilde f:\cY\to\cX$ is proper dominant with generically finite generic
fibre of degree $d$, which is exactly the hypothesis of
\Cref{thm:model-projection-formula}, so finiteness of $g$ was never
needed and the birational case $d=1$
(\Cref{thm:ch30-birational-identification}) is included.  Case (2) is
\Cref{thm:base-change-ma}: there the model-level input is not
\Cref{thm:model-projection-formula} but the ramification cancellation
\eqref{eq:ch30-coefficient-cancellation}, which plays the same role.  The
extension under the additional ampleness hypothesis is
\Cref{thm:ch30-finite-bounded-projection} and
\Cref{thm:ch30-finite-energy-transfer} in case (1), and
\Cref{thm:ch30-base-change-psh} in case (2).
\end{proof}

\begin{remark}[The two model-level inputs are genuinely different]
\label{rem:ch30-two-inputs}
It is tempting to regard \eqref{eq:ch30-coefficient-cancellation} as the
projection formula \eqref{eq:ch30-finite-projection-input} applied to the
finite morphism $\rho:\cX'\to\cX$ of degree $[K':K]$; that would give a
factor $[K':K]$ and be wrong.  The discrepancy is the convention on
degrees: \eqref{eq:ch30-finite-projection-input} is an identity between
intersection numbers computed over one and the same base field, whereas the
coefficients of the Chambert--Loir measure upstairs are degrees over
$k'$.  The passage from one convention to the other divides by
$f=[k':k]$, and \eqref{eq:ch30-multiplicity-relation} divides by
$e$; the product is $ef=[K':K]$, which is exactly the factor that must
disappear.  This is the same bookkeeping as in
\Cref{cor:trop-finite-field-extension}, where the calibration of a
polyhedral subspace is multiplied by $[F:K]$ precisely because the base
field, and hence the convention, has been changed.
\end{remark}

\subsection{Pushforward is not an inverse operation}
\label{sec:ch30-not-inverse}

The following elementary observation prevents several false functorial
claims.

\begin{proposition}
\label{prop:ch30-pushforward-not-inverse}
Let $f:Y\to X$ be finite of degree $d$.  If
$\nu=\MA_{f^*\theta}(f^*\varphi)$, then
$f_*\nu=d\MA_\theta(\varphi)$.  If a fibre contains at least two points, unrestricted pushforward on
positive Radon measures is not injective: in general
$f_*\nu=d\mu$ does not determine $\nu$ or identify it as the measure
of a pulled-back potential.  This assertion of nonuniqueness is not made
for an isomorphism or for measures constrained by an invariant-lift
condition.
\end{proposition}

\begin{proof}
The first assertion is \Cref{thm:finite-ma-projection}.  For the
second, any fibre containing two points admits distinct positive
distributions with the same total fibre mass and hence the same pushforward.
Monge--Amp\`ere structure imposes additional nonlinear conditions that are
not encoded by pushforward alone.
\end{proof}

\begin{example}[Two measures with the same pushforward]
\label{ex:ch30-two-measures}
Keep the square map $f(z)=z^2$ of \Cref{ex:ch30-power-map}, and choose
$b\in K^\times$.  Since the residue characteristic is zero, the two rigid
points
\[
                    y_+:=(z=b),\qquad y_-:=(z=-b)
\]
are distinct and have the same image $x:=(z=b^2)$.  Consequently
\[
 f_*\bigl(2\delta_{y_+}\bigr)
   =2\delta_x
   =f_*\bigl(\delta_{y_+}+\delta_{y_-}\bigr),
\]
although the two measures upstairs are different.  More generally every
measure
\[
              2t\,\delta_{y_+}+2(1-t)\,\delta_{y_-},
              \qquad 0\leq t\leq1,
\]
has the same pushforward.  This is only a measure-theoretic example: it
does not assert that these measures are Monge--Amp\`ere measures.  Its point
is precisely that the pushforward equality alone cannot decide that extra
nonlinear condition.
\end{example}

Similarly, a pushforward of currents is available under proper maps in the
CLD calculus, but a pullback of arbitrary currents is not.  Pullback is safe
for smooth forms, delta-forms under their transversality hypotheses, Cartier
divisors, and metrized line bundles.  It is not a formal dual of current
pushforward.

\begin{remark}[What would be needed for a pullback of measures]
\label{rem:ch30-measure-pullback}
For a finite flat morphism $g$ of degree $d$ between \emph{schemes}
there is a norm map on functions and a transfer on cycles, and these do
produce a pullback of \emph{cycles}.  The analytic analogue would require a
consistent assignment of local degrees to the points of a fibre of
$g^{\An}$, and such an assignment exists in the tropical setting of
\Cref{chap:analytic-tropicalization}, where it is exactly the local degree
$\deg_y(g)$ of \Cref{thm:trop-multiplicity-formula}.  The pushforward multiplicity formula alone does not prove a transfer
theorem for all Radon measures; such a theorem also needs the weighted
trace on continuous test functions.  Its possible existence is not
contradicted by noninjectivity of unrestricted pushforward.
In the finite Galois case
\eqref{eq:ch30-invariant-measure-lift} constructs a selected lift
directly, and a closed embedding such as $\sigma$ supplies a lift by
its own pushforward.  The operator projection identities in this
chapter remain statements about $g_*$.
\end{remark}

\section{The Gauss extension}
\label{sec:ch30-gauss}

The most useful bridge from a trivially valued field to a nontrivially
valued one is not an identification of analytic spaces, but a canonical
section built from Gauss valuations.

\subsection{The special base change and its additional section}
\label{sec:ch30-gauss-setting}

Throughout this section $k$ is a trivially valued field, algebraically
closed of characteristic zero unless the contrary is said, $X$ is a smooth integral projective $k$-variety of dimension $n$,
and $L$ is an ample
line bundle on $X$ with numerical class $\omega=c_1(L)$ and
\[
                       V:=(\omega^n)>0 .
\]
We put
\[
 K:=k(\!(t)\!),\qquad R:=k[\![t]\!],\qquad v_K(t)=1,
\]
so that, in this section only, the uniformizer is called $t$ rather than
$\varpi$, in agreement with \eqref{eq:gauss-extension} below.  We write
$X_K=X\times_kK$, $L_K$, $\omega_K$ for the curvature of the trivial model metric on $L_K$.
Its generic numerical class is $c_1(L_K)$; distinguishing the metric
reference from that class fixes the meaning of the potentials below.

Three observations fix the scope.

\emph{The nontrivially valued side satisfies \textup{(BFJ)}.}  The field
$K=k(\!(t)\!)$ is complete and discretely valued, with residue field
$k$ of characteristic zero, and $X_K$ is an integral projective
$K$-variety, smooth under the stated hypothesis on $X$.  The trivial
model metric is semipositive with ample class, so
\eqref{eq:ch20-standing} holds for
$(X_K,\omega_K)$ and the whole of
\Crefrange{chap:bfj-model-forms}{chap:bfj-continuity} applies verbatim on
that side.  This is the reason the trivially valued theory of
\Cref{part:bj} can borrow from \Cref{part:bfj} at all, and it is worth
stating as the single structural fact of the section.

\emph{The trivially valued side is not a specialization of the
nontrivially valued one.}  The field $k$ is complete for the trivial
absolute value, and $X^{\An}$ is defined; but \eqref{eq:ch20-standing}
fails at its first line, models over $k^\circ=k$ carry no information,
and $\mathcal D(X)$ as defined in \Cref{chap:bfj-model-forms} would be
zero.  The trivially valued theory replaces models over the ground ring by
test configurations over $\mathbf A^1_k$
(\Cref{def:bj-test-configuration}); it is a different construction, not a
degenerate case.

\emph{The base-change projection and the comparison section.}  The inclusion
$k\subset K$ is an extension of complete valued fields, so
\Cref{lem:ch30-base-change-map} provides a continuous proper surjection
\[
 r:X_K^{\An}\longrightarrow X^{\An} .
\]
This is the ordinary base-change map of case (B).  There is no map in the
other direction induced by a morphism of schemes.  What exists instead is a
canonical continuous \emph{section} $\sigma$ of $r$, whose existence is
special to the situation in which the smaller field is trivially valued.
Since $\sigma$ is a section, $r$ restricts to the inverse of $\sigma$ on
its image, and every statement below is naturally stated as a compatibility
between $\sigma^*$ on functions and $r_*$ on measures.

\subsection{Construction and first properties}
\label{sec:ch30-gauss-construction}

Let $k$ be trivially valued and put $K=k((t))$ with $v_K(t)=1$.  On an
affine chart $\Spec A\subset X$, a point $x\in X^{\An}$ is represented
by a semivaluation $v_x$ on $A$.  Define on Laurent polynomials
\begin{equation}
\label{eq:gauss-extension}
 \sigma(v_x)\left(\sum_i f_it^i\right)
       :=\min_i\{v_x(f_i)+i\}.
\end{equation}
This extends to the relevant completed residue fields and glues to a map
\[
              \sigma:X^{\An}\longrightarrow X_K^{\An}.
\]
Restriction of a valuation from $A\widehat\otimes_kK$ to $A$ defines a
map $r:X_K^{\An}\to X^{\An}$, and
\begin{equation}
\label{eq:gauss-section}
                              r\circ\sigma=\mathrm{id}.
\end{equation}

We now prove that this is well posed.  The only point requiring an argument
is multiplicativity, and the argument is the Gauss lemma in the form in
which it is used repeatedly in \Cref{part:bj}.

\begin{theorem}
\label{thm:ch30-gauss-well-defined}
\index{Gauss extension!construction}
Let $A$ be a finitely generated $k$-algebra and let
$v:A\to\mathbf R\cup\{+\infty\}$ be a semivaluation trivial on
$k^\times$.
\begin{enumerate}
\item\label{itm:ch30-gw:expansion}
Every element of $A\otimes_kK$ has a unique expansion
$F=\sum_{i\geq i_0}f_it^i$ with $f_i\in A$ lying in a common
finite-dimensional $k$-subspace of $A$; the infimum in
\eqref{eq:gauss-extension} is attained.
\item\label{itm:ch30-gw:semival}
Formula \eqref{eq:gauss-extension} defines a semivaluation
$\sigma(v)$ on $A\otimes_kK$ which restricts to $v$ on $A$ and
satisfies $\sigma(v)(t)=1$; in particular it restricts to $v_K$ on
$K$.
\item\label{itm:ch30-gw:glue}
The construction is compatible with localization, hence glues over $X$
to a map $\sigma:X^{\An}\to X_K^{\An}$ satisfying
\eqref{eq:gauss-section}.
\item\label{itm:ch30-gw:cont}
$\sigma$ is continuous and a homeomorphism onto its image.
For proper $X$, as in this section, that image is compact and closed
in $X_K^{\An}$.
\end{enumerate}
\end{theorem}

\begin{proof}
\cref{itm:ch30-gw:expansion}  Choose a $k$-basis $(e_\lambda)$ of $A$.
Then $A\otimes_kK=\bigoplus_\lambda K\,e_\lambda$, and an element is a
finite sum $\sum_\lambda c_\lambda e_\lambda$ with
$c_\lambda\in K$; collecting the coefficients of $t^i$ in the finitely
many Laurent series $c_\lambda$ gives the expansion, with all $f_i$ in
the span $W$ of the finitely many $e_\lambda$ occurring.  Uniqueness is
the directness of the sum.  If $e_1,\ldots,e_m$ is the resulting finite
list, the ultrametric inequality and triviality on $k^\times$ give
\[
 v(f_i)\geq \min_jv(e_j)
\]
whenever $f_i\neq0$ (with the assertion tautological if every $e_j$ lies
in $\Supp(v)$).  Thus $v(f_i)+i\to+\infty$, so only finitely many indices
can give a value below any fixed bound and the infimum is attained.

\cref{itm:ch30-gw:semival}  Normalization and the ultrametric inequality are
immediate from \eqref{eq:gauss-extension}: for $F=\sum f_it^i$ and
$G=\sum g_it^i$, the $i$-th coefficient of $F+G$ is $f_i+g_i$ and
$v(f_i+g_i)\geq\min\{v(f_i),v(g_i)\}$.  Also
$\sigma(v)(f)=v(f)$ for $f\in A$ and $\sigma(v)(t)=1$.

For multiplicativity, let $\lambda:=\sigma(v)(F)$ and
$\mu:=\sigma(v)(G)$.  If $\lambda=+\infty$, then $v(f_i)=+\infty$ for
every $i$, so every coefficient of $FG$ lies in the prime
$\Supp(v)$ and $\sigma(v)(FG)=+\infty$; similarly if
$\mu=+\infty$.  Assume therefore $\lambda,\mu<+\infty$.  Writing
$FG=\sum_lh_lt^l$ with $h_l=\sum_{i+j=l}f_ig_j$, the ultrametric
inequality gives
\[
 v(h_l)+l\geq\min_{i+j=l}\bigl(v(f_i)+i+v(g_j)+j\bigr)\geq\lambda+\mu,
\]
so $\sigma(v)(FG)\geq\lambda+\mu$.  For the reverse inequality let
$i_0$ be the \emph{smallest} index with $v(f_{i_0})+i_0=\lambda$ and
$j_0$ the smallest with $v(g_{j_0})+j_0=\mu$, and consider
$l_0=i_0+j_0$.  In the sum $h_{l_0}=\sum_{i+j=l_0}f_ig_j$ the term
$i=i_0,j=j_0$ has
\[
 v(f_{i_0}g_{j_0})=v(f_{i_0})+v(g_{j_0})=\lambda+\mu-l_0 ,
\]
while every other term is strictly larger: if $i<i_0$, then
$v(f_i)+i>\lambda$ by minimality of $i_0$ and $v(g_j)+j\geq\mu$; if
$i>i_0$, then $j<j_0$, so $v(g_j)+j>\mu$ and $v(f_i)+i\geq\lambda$.
A sum in which the minimum of the values is attained exactly once has value
equal to that minimum, so $v(h_{l_0})=\lambda+\mu-l_0$ and
\[
 \sigma(v)(FG)\leq v(h_{l_0})+l_0=\lambda+\mu .
\]
Hence $\sigma(v)(FG)=\sigma(v)(F)+\sigma(v)(G)$.

\cref{itm:ch30-gw:glue}  If $f\in A$ with $v(f)<+\infty$, then
$\sigma(v)$ extends to $A_f\otimes_kK$ by
$\sigma(v)(F/f^m)=\sigma(v)(F)-m\,v(f)$, and the formula
\eqref{eq:gauss-extension} is unchanged; the same computation on the
completed residue field handles points with nonzero support, since
\eqref{eq:gauss-extension} descends to $A/\Supp(v)$ and extends to its
fraction field and completion.  Thus the locally defined maps agree on
overlaps and glue.  Identity \eqref{eq:gauss-section} is
$\sigma(v)|_A=v$, which is
\cref{itm:ch30-gw:semival}.

\cref{itm:ch30-gw:cont}  The topology of $X_K^{\An}$ is the coarsest making
all functions $y\mapsto|F(y)|$ continuous, for $F$ a regular function on
an affine chart.  For $F=\sum_{i\geq i_0}f_it^i$, choose
$e_1,\ldots,e_m$ spanning the common finite-dimensional coefficient space
from \cref{itm:ch30-gw:expansion}.  Then
\[
 \bigl|F(\sigma(v))\bigr|=\max_i|f_i(v)|\,\rme^{-i},
\]
and, by the ultrametric inequality,
$|f_i(v)|\leq\max_j|e_j(v)|$.  Around any fixed $v_0$, continuity of the
finitely many functions $|e_j|$ gives a neighbourhood on which the right
side is bounded by one constant $C$.  On that neighbourhood the partial
maxima over $i\leq N$ converge to the full maximum with tail bounded by
$C\rme^{-N}$.  They are continuous, so their locally uniform limit
$|F\circ\sigma|$ is continuous.  Thus $\sigma$ is continuous.  The section identity and continuity of $r$ show directly that $\sigma$
is a homeomorphism onto its image, with inverse the restriction of $r$.
For proper $X$, compactness
(\Cref{thm:berk-proper-compact}) and the Hausdorff property of
$X_K^{\An}$ additionally make this image compact and closed.
\end{proof}

\subsection{Centers and supports}
\label{sec:ch30-gauss-centers}

The comparison programme of \Cref{part:comparisons} requires four data to be
tracked along $\sigma$: centers, supports, line-bundle classes, and total
Monge--Amp\`ere mass.  The first two are settled by a direct computation.

\begin{theorem}
\label{thm:ch30-gauss-centers-supports}
\index{Gauss extension!centers}\index{Gauss extension!supports}
Let $v\in X^{\An}$, with center $c_X(v)\in X$ and support
$\Supp(v)\in X$, and let $X_R:=X\times_kR$ be the trivial model of
$X_K$, whose special fibre is $X$ itself.
\begin{enumerate}
\item\label{itm:ch30-gcs:center}
$\sigma(v)$ has a center on $X_R$, and
\begin{equation}
\label{eq:ch30-gauss-center}
 c_{X_R}\bigl(\sigma(v)\bigr)=c_X(v)
 \qquad\text{under }(X_R)_s=X .
\end{equation}
\item\label{itm:ch30-gcs:support}
\begin{equation}
\label{eq:ch30-gauss-support}
 \Supp\bigl(\sigma(v)\bigr)=\Supp(v)\cdot(\cO_X\otimes_kK),
\end{equation}
that is, the support of $\sigma(v)$ is the generic point of the base
change to $K$ of the closure of $\Supp(v)$.  In particular $\sigma$
maps points of zero support to points of zero support, and maps the
proper-support point $v_{Z,\triv}$ attached to an integral closed
subvariety $Z\subset X$ to a point of $X_K^{\An}$ supported at the
generic point of $Z_K$; for $Z$ a closed $k$-point this is a rigid
point of $X_K^{\An}$.
\item\label{itm:ch30-gcs:divisorial}
$\sigma$ carries the divisorial point $c\,\ord_E$, $c\in\mathbf
Q_{>0}$ and $E$ a prime divisor over $X$, to a divisorial point of
$X_K^{\An}$, and quasimonomial points to quasimonomial points.
\end{enumerate}
\end{theorem}

\begin{proof}
\cref{itm:ch30-gcs:center}  Choose an affine neighbourhood
$U=\Spec A$ of $c_X(v)$.  By the defining property of the center,
$v\geq0$ on $A$.  An element of $\cO_{X_R}$ over this chart is
$F=\sum_{i\geq0}f_it^i$, so
$\sigma(v)(F)=\min_{i\geq0}\{v(f_i)+i\}\geq0$ because $v\geq0$.  The
center of a nonnegative semivaluation on $\Spec B$ is the prime
$\{F\in B:\sigma(v)(F)>0\}$, and here
\[
 \sigma(v)(F)>0
 \iff v(f_0)>0,
\]
since the terms with $i\geq1$ contribute at least $1$.  Thus the center
is the preimage under $B=A\otimes_kR\to A$, $F\mapsto f_0$, of the prime
$\{f\in A:v(f)>0\}$, which is $c_X(v)$.  This preimage is a prime of
$B$ containing $t$, that is, a point of the special fibre
$(X_R)_s=X$, namely $c_X(v)$ itself.

\cref{itm:ch30-gcs:support}  By \eqref{eq:gauss-extension},
$\sigma(v)(F)=+\infty$ if and only if $v(f_i)=+\infty$ for every $i$,
that is, if and only if every coefficient of $F$ lies in
$\mathfrak p:=\Supp(v)$.  The set of such $F$ is
$\mathfrak p\otimes_kK=\mathfrak p\cdot(A\otimes_kK)$, which is prime
because $(A/\mathfrak p)\otimes_kK$ is a domain --- $K/k$ being a
separable, indeed a regular, field extension and $A/\mathfrak p$ a
domain.  The remaining assertions are the special cases
$\mathfrak p=0$ and $\mathfrak p=I(Z)$.

\cref{itm:ch30-gcs:divisorial}  This is
\cite[\S1]{BJ22}, restated in \Cref{chap:trivial-analytification} after
\eqref{eq:triv-gauss-section}; the concrete mechanism is
\Cref{thm:ch30-gauss-model-dictionary} below, which exhibits $\sigma(v_E)$
as the divisorial point of a component of the special fibre of a model of
$X_K$, and \Cref{ex:ch08-gauss-ray} in the simplest case.
\end{proof}

\begin{remark}[The two endpoints of a ray]
\label{rem:ch30-gauss-ray-endpoints}
\Cref{thm:ch30-gauss-centers-supports} makes the picture of
\Cref{ex:ch08-gauss-ray} precise.  A ray $\{c\,\ord_E:c\geq0\}$ in
$X^{\An}$ has two endpoints: the trivial valuation $v_{X,\triv}$ at
$c=0$, of zero support and center the generic point of $X$, and a
proper-support point at $c=+\infty$.  Under $\sigma$, the first goes to
the Gauss point of $X_K^{\An}$ attached to the trivial model $X_R$ ---
its center on $X_R$ is the generic point of the special fibre, by
\eqref{eq:ch30-gauss-center} --- and the second goes to a point of positive
support, that is, to a point which no Monge--Amp\`ere measure of a bounded
potential can charge
(\Cref{prop:ch24-model-support}\,\cref{itm:ch24-ms:pluripolar}).  A positive rational weight on the ray gives a divisorial point upstairs;
an arbitrary positive real weight gives a quasimonomial point.
The image is a ray in the monomial locus, not a bounded segment.
The trivial valuation can also carry model mass and maps to the
trivial-model Gauss point.
\end{remark}

\subsection{Scaling, and the dependence on the normalization}
\label{sec:ch30-gauss-scaling}

The section depends on the choice $v_K(t)=1$.  Replacing $v_x$ by
$rv_x$ produces the scaling action central to the trivially valued theory.
The precise relation is the following, and it is the form in which the
scaling action of \Cref{sec:trivial-valuation-cone} is compared with a
change of normalization of the ground valuation.

\begin{proposition}
\label{prop:ch30-gauss-scaling}
Let $K_\lambda$ denote the field $k((t))$ with $v(t)=\lambda$.
For $\lambda>0$ define
$\sigma_\lambda:X^{\An}\to X_{K_\lambda}^{\An}$ by
\begin{equation}
\label{eq:ch30-gauss-lambda}
 \sigma_\lambda(v)\Bigl(\sum_if_it^i\Bigr):=\min_i\{v(f_i)+\lambda i\},
\end{equation}
so that $\sigma_1=\sigma$ and $\sigma_\lambda$ is the Gauss extension
formed with the normalization $v_K(t)=\lambda$.  Then for all $r>0$
\begin{equation}
\label{eq:ch30-gauss-scaling}
 \sigma(r\cdot v)=r\cdot\sigma_{1/r}(v),
 \qquad
 \sigma(r\cdot v)|_A=r\,v,
 \qquad
 \sigma(r\cdot v)(t)=1 .
\end{equation}
Consequently the $\mathbf R_{>0}$-action of
\eqref{eq:appJ-scaling-action} on $X^{\An}$ corresponds, under
$\sigma$, to changing the normalization of the ground valuation of $K$
and rescaling.  Multiplying an upstairs valuation by $r$ alone
does not stay in the analytification with its original fixed ground
absolute value.
\end{proposition}

\begin{proof}
By \eqref{eq:gauss-extension} and \eqref{eq:ch30-gauss-lambda},
\[
 \sigma(rv)\Bigl(\sum f_it^i\Bigr)=\min_i\{r\,v(f_i)+i\}
 =r\min_i\Bigl\{v(f_i)+\frac ir\Bigr\}
 =r\,\sigma_{1/r}(v)\Bigl(\sum f_it^i\Bigr).
\]
The two evaluations are the cases $F=f\in A$ and $F=t$.  For the final assertion, multiplication of a valuation on the full
scalar-extended function field by $r$ changes $v(t)=1$ to $v(t)=r$.
The formula compensates by starting with $K_{1/r}$ before multiplying.
It makes no claim excluding every other possible action on the analytic
space.
\end{proof}

\newpage
\begin{warning}{Two rescalings that look alike}
The renormalization of \Cref{sec:ch30-bc-normalization} and the scaling
action of \Cref{prop:ch30-gauss-scaling} both multiply valuations by a
positive constant, and they are constantly confused.  They are different
operations: the first rescales the ground valuation on $K'$ and
therefore its restriction to $K$, comparing different ground absolute
values by the bookkeeping
governed by \eqref{eq:ch30-ramified-renormalization}; the second moves a
point of the trivially valued space $X^{\An}$ along its ray, and is part
of the structure of that space.  On $X^{\An}$, every trivially valued point
$v_{Z,\triv}$ is fixed by scaling, including $v_{X,\triv}$.
On the zero-support valuation locus only the trivial valuation is fixed
by a scale factor different from one.  Multiplication of valuations on
the nontrivially valued side changes the ground absolute value and
therefore compares different analytic spaces.
\end{warning}

\subsection{The image of the Gauss extension}
\label{sec:ch30-gauss-image}

\begin{theorem}
\label{thm:ch30-gauss-image}
\index{Gauss extension!image}
Assume $k$ infinite --- which holds here, $k$ being algebraically
closed by the standing hypothesis of \Cref{sec:ch30-gauss-setting}.  Let
$\mathbf G_m(k)=k^\times$ act on $K=k(\!(t)\!)$ by $t\mapsto\lambda t$
and hence on $X_K=X\times_kK$ over $X$, and therefore on
$X_K^{\An}$.  Then
\begin{equation}
\label{eq:ch30-gauss-image}
 \sigma\bigl(X^{\An}\bigr)
 =\bigl\{y\in X_K^{\An}\ :\ y\ \text{is}\ \mathbf G_m(k)\text{-invariant
   and}\ y(t)=1\bigr\},
\end{equation}
and $\sigma\circ r$ is the identity on this set.  It need not be the
identity on the whole space, as the projective-line example below shows.
\end{theorem}

\begin{proof}
\emph{The image is contained in the right-hand side.}  For
$\lambda\in k^\times$ and $F=\sum f_it^i$, the translate is
$F_\lambda=\sum f_i\lambda^it^i$, and
$v(f_i\lambda^i)=v(f_i)$ because $v$ is trivial on $k^\times$; so
$\sigma(v)(F_\lambda)=\sigma(v)(F)$.  Also $\sigma(v)(t)=1$.

\emph{Conversely.}  Let $y\in X_K^{\An}$ be $\mathbf G_m(k)$-invariant
with $y(t)=1$, choose an affine chart $\Spec A$ containing its support,
and put $v:=r(y)=y|_A$.  If every coefficient of $F$ lies in $\Supp(v)$, a finite basis of
the coefficient space expresses $F$ in
$\Supp(v)\cdot(A\otimes_kK)$, so both $y(F)$ and the proposed minimum
are $+\infty$.  Otherwise the coefficient minimum is finite.
It suffices in that case to treat
$F=\sum_{i=i_0}^{i_1}f_it^i$ with finitely many terms.  Indeed, for a
general $F=\sum_{i\geq i_0}f_it^i$ let $F_N$ be its truncation at
$i=N$.  If $e_1,\ldots,e_m$ span the common coefficient space, then
$v(f_i)\geq c:=\min_jv(e_j)$, and hence
$y(F-F_N)\geq c+N+1$.  The finite Laurent-polynomial argument below gives
$y(F_N)=\min_{i\leq N}\{v(f_i)+i\}$.
Once $N$ includes an index attaining the finite coefficient minimum
and $c+N+1$ exceeds it, the tail has strictly larger value.
The ultrametric equality then gives $y(F)=y(F_N)$, and both
sides are unchanged by truncation.  The
inequality $\geq$ is the ultrametric inequality together with
$y(f_it^i)=v(f_i)+i$.  For $\leq$, choose $N:=i_1-i_0+1$ distinct
elements $\lambda_1,\ldots,\lambda_N\in k^\times$, which exist since
$k$ is infinite.  The Vandermonde matrix
$\bigl(\lambda_s^{\,i}\bigr)_{1\leq s\leq N,\,i_0\leq i\leq i_1}$ is
invertible over $k$, so each monomial $f_it^i$ is a $k$-linear
combination of the translates $F_{\lambda_1},\ldots,F_{\lambda_N}$.  Since
$y$ is trivial on $k^\times$ and ultrametric,
\[
 v(f_i)+i=y(f_it^i)\geq\min_sy(F_{\lambda_s})=y(F),
\]
the last equality by $\mathbf G_m(k)$-invariance.  Taking the minimum over
$i$ gives $\min_i\{v(f_i)+i\}\geq y(F)$, as required.  Hence
$y=\sigma(v)$.

The section identity gives the assertion on the invariant image.
\Cref{ex:ch30-gauss-image-strict} establishes that equality need not
extend to the whole space; no universal strictness assertion is made
in dimension zero.
\end{proof}

\begin{example}[The image meets each fibre of $r$ in one point, and the
fibres are big]
\label{ex:ch30-gauss-image-strict}
Take $X=\mathbf P^1_k$ with coordinate $z$ and $K=k(\!(t)\!)$.  Set
$w:=z-1-t\in K(X_K)$ and let $y$ be the monomial valuation in the
coordinates $(w,t)$ with weights
\[
                     y(w)=2,\qquad y(t)=1 .
\]
It is a divisorial point of $X_K^{\An}$, obtained by blowing up
$\mathbf P^1_R$ twice along the section $z=1+t$, and it satisfies
$y(t)=1$.  It is not $\mathbf G_m(k)$-invariant: for
$\lambda\neq1$,
\[
 y\bigl(z-1-\lambda t\bigr)
 =y\bigl(w+(1-\lambda)t\bigr)=\min\{2,1\}=1\neq2=y(w),
\]
so $y$ and its translate differ.  By
\Cref{thm:ch30-gauss-image}, $y\notin\sigma(X^{\An})$, and one sees this
directly: $r(y)=y|_{k(z)}$ satisfies $r(y)(z-1)=y(w+t)=1$, so
$r(y)=\ord_1$, while
$\sigma(\ord_1)(w)=\min\{\ord_1(z-1),\,\ord_1(1)+1\}=1\neq2=y(w)$.  Thus
the fibre $r^{-1}(\ord_1)$ contains at least the two distinct points
$\sigma(\ord_1)$ and $y$; in fact it is infinite, since the weight
$y(w)$ may be any real number $\geq1$.  The Gauss image is the
$\mathbf G_m(k)$-invariant slice of this fibration, and it meets each
fibre exactly once.
\end{example}

\subsection{Classes, models, and PL functions}
\label{sec:ch30-gauss-classes}

\begin{proposition}\label{prop:gauss-model-compatibility}
A test configuration over $\mathbf A^1_k$, completed along its central
fibre, gives a formal model after base change to $k[[t]]$.  On the Gauss
image $\sigma(X^{\An})$, its model function is the function defined by the
same vertical divisor or filtration.  The vertical intersection
coefficients used to define model energies and model Monge--Amp\`ere
measures are the same after the normalizations of Appendix~J.
\end{proposition}

\begin{proof}
Both model functions evaluate a local equation by the minimum formula
\eqref{eq:gauss-extension}.  The central fibre of the test configuration is
the special fibre of its $t$-adic completion.  Proper intersection numbers
are invariant under completion, which proves the asserted coefficient
comparison.  A global equality of Radon measures would additionally require
an explicitly specified pushforward along the Gauss embedding; none is
being silently imposed here.
\end{proof}

The next theorem is the quantitative form of
\Cref{prop:gauss-model-compatibility}, with the pushforward supplied and the
four tracked data displayed.  It is the statement on which
\Cref{part:bj} and \Cref{part:comparisons} rely.

\begin{theorem}
\label{thm:ch30-gauss-model-dictionary}
\index{Gauss extension!test configurations}
Let $(\cX,\cL)$ be a normal test configuration for $(X,L)$ in the sense
of \Cref{def:bj-test-configuration}.  Pass, if necessary, to a normal equivariant graph model dominating
both $\cX$ and $X\times\mathbf A^1$, and pull back $\cL$.
All formulas are invariant under this domination.
Denote the resulting morphism to the trivial configuration by $\rho$,
and put
\[
 \cX_R:=\cX\times_{\mathbf A^1_k}\Spec R,
 \qquad
 \cL_R:=\cL|_{\cX_R} .
\]
\begin{enumerate}
\item\label{itm:ch30-gmd:model}
$\cX_R$ is a normal model of $X_K$ in the sense of
\eqref{eq:ch20-standing}, $\cL_R$ is a model line bundle for $L_K$, and
the special fibre $(\cX_R)_s$ is canonically identified with the central
fibre $\cX_0$, with the same irreducible components and the same
multiplicities $b_E=\ord_E(t)$.
\item\label{itm:ch30-gmd:points}
For every component $E$ of $\cX_0$, the divisorial point
$x_E=b_E^{-1}\ord_E$ of $X_K^{\An}$ lies in the image of $\sigma$, and
\begin{equation}
\label{eq:ch30-gauss-divisorial}
 x_E=\sigma(v_E),
 \qquad v_E:=b_E^{-1}\ord_E|_{k(X)}\in X^{\An} .
\end{equation}
\item\label{itm:ch30-gmd:function}
Writing $\phi_{\cX,\cL}\in\PL(X)$ for the associated PL
function of \Cref{thm:tc-pl-correspondence} and $\varphi_{\cX_R,\cL_R}$ for the
model function of \Cref{chap:bfj-model-forms}, one has
\begin{equation}
\label{eq:ch30-gauss-pl}
 \phi_{\cX,\cL}=\sigma^*\varphi_{\cX_R,\cL_R}
 =\varphi_{\cX_R,\cL_R}\circ\sigma .
\end{equation}
\item\label{itm:ch30-gmd:intersections}
All intersection terms containing a vertical factor agree.
For test configurations $(\cX,\cL_0),\ldots,(\cX,\cL_n)$
on a common graph model, write $M=\rho^*L_{\mathbf A^1}$ and
$\cL_i=M+D_i$ with $D_i$ vertical.  Then
\begin{equation}
\label{eq:ch30-gauss-intersection}
 \bigl(\overline{\cL}_0\cdots\overline{\cL}_n\bigr)_{\overline\cX}
 =\sum_{\varnothing\ne I\subseteq\{0,\ldots,n\}}
   \left(\prod_{i\in I}D_{i,R}
         \cdot M_R^{\,n+1-|I|}\right)_{\cX_R},
\end{equation}
where the left side uses the standard trivialized compactification
\eqref{eq:tc-compactification}.  Every term on the right has at least
one vertical divisor, so is defined by
\Cref{lem:ch24-vertical-numbers}.
There is no intrinsic degree of an arbitrary product of $n+1$
ordinary line bundles on the model over the discrete valuation ring.
\item\label{itm:ch30-gmd:class}
$L$ is ample if and only if $L_K$ is, and
\begin{equation}
\label{eq:ch30-gauss-volume}
 (\omega_K^n)_{X_K}=(\omega^n)_X=V .
\end{equation}
\end{enumerate}
\end{theorem}

\begin{proof}
\cref{itm:ch30-gmd:model}  The morphism $\Spec R\to\mathbf A^1_k$ is
regular, being the base change of the completion of the excellent local ring
$\cO_{\mathbf A^1,0}$; hence normality and flatness ascend and $\cX_R$
is a normal flat projective $R$-scheme.  Its generic fibre is
$\cX\times_{\mathbf A^1}\Spec K$, which by
\eqref{eq:tc-punctured-trivialization} is $X\times_kK=X_K$, since the
generic point of $\Spec R$ maps into $\mathbf A^1\setminus\{0\}$.  Its
special fibre is the fibre of $\cX$ over $0$, namely $\cX_0$, with the
same scheme structure, hence the same components and multiplicities.

\cref{itm:ch30-gmd:points}  Each component $E$ of $\cX_0$ is preserved by
the connected group $\mathbf G_m$, hence $\ord_E$ is
$\mathbf G_m(k)$-invariant on $k(X)(t)$.
This subfield embeds in $K(X_K)$; it is not equal to it.
The induced component valuation on $K(X_K)$ remains invariant:
the field automorphisms preserve the model and its component, or one
can pass to coefficients in $k((t))$ by continuity.  Thus
$x_E(t)=b_E^{-1}\ord_E(t)=1$ by definition of $b_E$.  By
\Cref{thm:ch30-gauss-image}, $x_E=\sigma(r(x_E))$, and
$r(x_E)=x_E|_{k(X)}=v_E$.

\cref{itm:ch30-gmd:function}  Both sides are computed from the same local
equations.  On an affine chart of $\cX$, a local generator of
a Cartier multiple $qD$ of the vertical divisor
$D=\cL-\rho^*L_{\mathbf A^1}$ is a rational function $g$ on $\cX$.
By \eqref{eq:ch20-sign} its model weight upstairs is
$y\mapsto y(g)/q$, while the associated PL weight downstairs is
$v\mapsto\sigma(v)(g)/q$ by \Cref{thm:tc-pl-correspondence}.
These weights agree after composition with $\sigma$.
At matched divisorial points, the coefficient of $D$ divided by $b_E$
gives the same value on both sides by
\eqref{eq:ch30-gauss-divisorial}.  The local-equation calculation already proves the
identity at every point; it does not infer equality of two arbitrary
PL functions from the finite list of components of one configuration.

\cref{itm:ch30-gmd:intersections}  Write
$\cL_i=\rho^*L_{\mathbf A^1}+D_i$ with $D_i$ supported on $\cX_0$, as
in \eqref{eq:bj-common-model-pairs}, and expand
$(\cL_0\cdots\cL_n)_{\overline\cX}$ multilinearly.  The term with no
$D_i$ is $\bigl((\rho^*L_{\mathbf P^1})^{n+1}\bigr)=0$, because
$L_{\mathbf P^1}$ is pulled back from $X$ along
$\overline\cX\to X\times\mathbf P^1\to X$ and $L^{n+1}=0$ on the
$n$-dimensional $X$.  Every other term contains at least one factor
$D_i$, hence is an intersection number of cycles supported on $\cX_0$,
and such a number is unchanged by the flat base change
$\Spec R\to\mathbf A^1$, which has the same residue field and special fibre.  The proper
degree of each vertical intersection is computed on those fibre
cycles, and hence is unchanged under this flat completion.
The sum of these vertical terms on $\cX_R$ is exactly the right
side of \eqref{eq:ch30-gauss-intersection}, proving
\eqref{eq:ch30-gauss-intersection}.

\cref{itm:ch30-gmd:class}  Ampleness is preserved by the flat base change
$k\to K$ and descends along it, since $\Spec K\to\Spec k$ is faithfully
flat.  Equality \eqref{eq:ch30-gauss-volume} is the invariance of top
intersection numbers under extension of the base field.
\end{proof}

\subsection{Measures, mass, and energy}
\label{sec:ch30-gauss-mass}

We now record what the two theories compute, and in which normalization.
This is the point at which the plan's requirement to track total
Monge--Amp\`ere mass becomes a concrete statement about a factor of $V$.

\begin{center}
\fbox{\begin{minipage}{0.91\textwidth}
\textbf{Normalization checkpoint.}
The trivially valued theory of \Cref{part:bj} normalizes its
Monge--Amp\`ere measure to total mass $1$ and writes
$\MA^{\mathrm u}_\omega=V\MA_\omega$ for the unnormalized one
(\eqref{eq:bj-ma-normalizations}), and its energy satisfies
$E_\omega(\phi+c)=E_\omega(\phi)+c$
(\eqref{eq:bj-energy-basic-identities}).  The BFJ theory of this part uses
unnormalized measures of mass $V$ and energies with
$E(\varphi+c)=E(\varphi)+cV$.  Every identity below is therefore stated
with $\MA^{\mathrm u}$ on the trivially valued side, and the two energies
differ by exactly the factor $V$:
\begin{equation}
\label{eq:ch30-gauss-energy-normalization}
 E^{\mathrm{BFJ}}_{\omega_K}=V\cdot E^{\mathrm{BJ}}_{\omega}
 \quad\text{under the dictionary
 \eqref{eq:ch30-gauss-pl}}.
\end{equation}
When $V=1$, as on $\mathbf P^1$ with $\cO(1)$, the two coincide and
the discrepancy is invisible; \Cref{ex:ch30-gauss-two-rays} is included
precisely because there $V=2$.
\end{minipage}}
\end{center}

\begin{theorem}
\label{thm:ch30-gauss-measures}
\index{Gauss extension!Monge--Ampere mass}
Let $(\cX,\cL_1),\ldots,(\cX,\cL_n)$ be test configurations on a common
normal $\cX$, with associated PL functions $\phi_i\in\PL(X)$ and model
functions $\varphi_i$ on $X_K^{\An}$ related by
\eqref{eq:ch30-gauss-pl}, and assume each $\cL_i$ relatively
semipositive.  Then
\begin{equation}
\label{eq:ch30-gauss-measure-transfer}
 \begin{aligned}
 \sigma_*\Bigl(\bigwedge_{i=1}^n(\omega+\ddc\phi_i)^{\mathrm u}\Bigr)
 &=\bigwedge_{i=1}^n(\omega_K+\ddc\varphi_i),\\
 r_*\Bigl(\bigwedge_{i=1}^n(\omega_K+\ddc\varphi_i)\Bigr)
 &=\bigwedge_{i=1}^n(\omega+\ddc\phi_i)^{\mathrm u},
 \end{aligned}
\end{equation}
where the superscript $\mathrm u$ denotes the unnormalized trivially
valued product of \eqref{eq:bj-ma-normalizations}.  Both measures have total
mass
\begin{equation}
\label{eq:ch30-gauss-mass}
 (\omega^n)_X=(\omega_K^n)_{X_K}=V.
\end{equation}
The upstairs measure is carried by finitely many divisorial points
matched by \eqref{eq:ch30-gauss-divisorial}.  Their restrictions
downstairs are rational divisorial valuations or the trivial valuation,
which can also have positive mass.  Moreover
\begin{equation}
\label{eq:ch30-gauss-energy}
 E^{\mathrm{BFJ}}_{\omega_K}\bigl(\varphi\bigr)
 =V\,E^{\mathrm{BJ}}_{\omega}\bigl(\sigma^*\varphi\bigr)
\end{equation}
for every model function $\varphi$ on $X_K^{\An}$ arising from a test
configuration as in \Cref{thm:ch30-gauss-model-dictionary}.
For a nonsemipositive model metric, the two sides here mean the
polynomial intersection functional defined by the same displayed
formulas; for a psh model metric they are the usual energies.
\end{theorem}

\begin{proof}
By \eqref{eq:bj-model-ma-atomic} the unnormalized trivially valued product
is
\[
 \bigwedge_{i}(\omega+\ddc\phi_i)^{\mathrm u}
 =\sum_{E\subset\cX_0}b_E
   \bigl((\theta_{1,\cX}+D_1)|_E\cdots(\theta_{n,\cX}+D_n)|_E\bigr)
   \delta_{v_E},
\]
while by \eqref{eq:appJ-cl-measure} the BFJ product on $X_K^{\An}$,
computed on the model $\cX_R$, is
\[
 \bigwedge_i(\omega_K+\ddc\varphi_i)
 =\sum_{E\subset(\cX_R)_s}b_E
   \bigl(\cL_{1,R}|_E\cdots\cL_{n,R}|_E\bigr)\delta_{x_E}.
\]
By \Cref{thm:ch30-gauss-model-dictionary}\,\cref{itm:ch30-gmd:model} the two
index sets and the two multiplicities coincide, and the restricted line
bundles are the same; the intersection numbers are degrees over the residue
field, which is $k$ on both sides because $K=k(\!(t)\!)$ has residue
field $k$.  Hence the coefficients coincide.  By
\cref{itm:ch30-gmd:points} the atoms correspond under
$\sigma(v_E)=x_E$, and by
\Cref{lem:ch30-pushforward-basic}\,\cref{itm:ch30-pf:support}
$\sigma_*\delta_{v_E}=\delta_{x_E}$ and
$r_*\delta_{x_E}=\delta_{v_E}$, the latter by
\eqref{eq:gauss-section}.  This gives
\eqref{eq:ch30-gauss-measure-transfer}.  The total masses are the
generic-fibre intersection numbers on the two sides, equal by
\eqref{eq:ch30-gauss-volume}.

For \eqref{eq:ch30-gauss-energy}, the trivially valued energy is
$E^{\mathrm{BJ}}_\omega(\phi)=(\omega,\phi)^{n+1}/((n+1)V)$ by
\eqref{eq:bj-energy-definition}.  If $D_\varphi$ represents the BFJ model
function, then \eqref{eq:ch26-energy-intersection} gives
\[
 E^{\mathrm{BFJ}}_{\omega_K}(\varphi)
 =\frac1{n+1}\sum_{j=0}^n
   \bigl(D_\varphi\cdot\omega_{K,\varphi}^{\,j}
                         \cdot\omega_K^{\,n-j}\bigr).
\]
Every term has a vertical first factor and is matched by
\eqref{eq:ch30-gauss-intersection}; their sum is precisely the compactified
test-configuration number $(\omega,\phi)^{n+1}$.  Hence the two energies
differ exactly by the factor $V$.
\end{proof}

\begin{corollary}
\label{cor:ch30-gauss-four-data}
Along the Gauss extension:
\begin{center}
\small
\begin{tabular}{p{.20\textwidth}|p{.34\textwidth}|p{.36\textwidth}}
\toprule
\textbf{Datum}&\textbf{Trivially valued side}
 &\textbf{Statement}\\
\midrule
centers
 &$c_X(v)$ on $X$
 &$c_{X_R}(\sigma(v))=c_X(v)$,
  \eqref{eq:ch30-gauss-center}\\
supports
 &$\Supp(v)\subset X$
 &$\Supp(\sigma(v))=\Supp(v)_K$,
  \eqref{eq:ch30-gauss-support}\\
line-bundle classes
 &$\omega=c_1(L)\in\Amp(X)$
 &$\omega_K$ ample and
  $(\omega_K^n)=(\omega^n)$,
  \eqref{eq:ch30-gauss-volume}\\
total $\MA$ mass
 &$V$ unnormalized, $1$ normalized
 &$\sigma_*\MA^{\mathrm u}_\omega=\MA_{\omega_K}$, mass $V$,
  \eqref{eq:ch30-gauss-mass}\\
\bottomrule
\end{tabular}
\end{center}
\end{corollary}

\begin{proof}
This is
\Cref{thm:ch30-gauss-centers-supports},
\Cref{thm:ch30-gauss-model-dictionary}\,\cref{itm:ch30-gmd:class} and
\Cref{thm:ch30-gauss-measures}.
\end{proof}

\subsection{The Gauss extension of an explicit trivially valued situation}
\label{sec:ch30-gauss-examples}

\begin{example}[Deformation to the normal cone of a point on a line]
\label{ex:ch30-gauss-normal-cone}
\index{Gauss extension!worked example}
Let $k$ be algebraically closed of characteristic zero and trivially
valued, $X=\mathbf P^1_k$ with affine coordinate $z$,
$L=\cO_{\mathbf P^1}(1)$, so that
\[
                       \omega=c_1(L),\qquad V=(\omega)=1 .
\]
Fix $c\in\mathbf Z_{>0}$ and let $\mathfrak a:=\mathfrak m_0+(t^c)$ be
the flag ideal on $X\times\mathbf A^1_k$ of
\Cref{def:bj-flag-ideal}, where $\mathfrak m_0$ is the ideal of the point
$z=0$; its Fubini--Study function is, by
\eqref{eq:flag-ideal-function},
\begin{equation}
\label{eq:ch30-gauss-example-phi}
 \phi(v)=-\sigma(v)(\mathfrak a)
        =-\min\{v(\mathfrak m_0),\,c\}
        =\max\{-c,\,-v(\mathfrak m_0)\},
 \qquad v\in X^{\An}.
\end{equation}
Concretely, $\phi\equiv0$ off the ray $\{s\,\ord_0:s\geq0\}$, and
$\phi(s\,\ord_0)=-\min\{s,c\}$ on it.  For $c=1$ this is the
deformation to the normal cone of \Cref{ex:product-dnc-test-configurations}.

\emph{The test configuration and its base change.}  Let $\cY_c$ be
obtained from $X\times\mathbf A^1_k$ by blowing up the point $(0,0)$
and then, $c-1$ times, the point at which the strict transform of the
section $\{z=0\}$ meets the last exceptional divisor.  All the centers
are $\mathbf G_m$-invariant, so $\cY_c$ is a test configuration; its
central fibre is a chain $E_0+E_1+\cdots+E_c$ of projective lines, all of
multiplicity one, with
\[
 \ord_{E_i}(t)=1,\qquad \ord_{E_i}(z)=i .
\]
Polarize it by $\cL_c:=\rho^*L_{\mathbf A^1}+D$ with
$D=-\sum_{i=1}^ciE_i$.  Base changing to $R=k[\![t]\!]$ gives exactly
the model $\cX_c$ of $\mathbf P^1_K$ used in
\Cref{ex:ch30-power-map}, with the model function
\begin{equation}
\label{eq:ch30-gauss-example-Phi}
 \Phi=-\min\bigl\{c,\ \max\{0,v(z)\}\bigr\}
 =\varphi_D\in\mathcal D(X_K),
\end{equation}
and \Cref{thm:ch30-gauss-model-dictionary}\,\cref{itm:ch30-gmd:function}
gives $\sigma^*\Phi=\phi$.  One checks this by hand: for
$v\in X^{\An}$ one has $\sigma(v)(z)=v(z)$, which is $\geq0$ unless
$v$ is centered at $z=\infty$, where $v(\mathfrak m_0)=0$; so
$\Phi(\sigma(v))=-\min\{c,\max\{0,v(z)\}\}=\phi(v)$.

\emph{Centers.}  For $s>0$ the point $s\,\ord_0\in X^{\An}$ has center
the closed point $0\in\mathbf P^1_k$; by
\eqref{eq:ch30-gauss-center}, $\sigma(s\,\ord_0)$ has center the same
point $0$ of the special fibre of $X_R=\mathbf P^1_R$.  On the model
$\cX_c$, the center of $\sigma(i\,\ord_0)$ is the generic point of
$E_i$.  The vertex $v_{X,\triv}$ has center the generic point of $X$,
and $\sigma(v_{X,\triv})$ is the Gauss point $x_G$, whose center on
$\mathbf P^1_R$ is the generic point of the special fibre.

\emph{Supports.}  Every point $s\,\ord_0$ with $0\leq s<\infty$ has zero
support, and so does its image, by
\eqref{eq:ch30-gauss-support}.  The endpoint $v_{0,\triv}$ of the ray, of
support $\mathfrak m_0$, maps to the point of $X_K^{\An}$ of support
$\mathfrak m_0K[z]$, that is, to the rigid point $z=0$ of
$\mathbf P^{1,\An}_K$.  Neither measure below charges it.

\emph{Classes.}  $L$ is ample on $\mathbf P^1_k$, $L_K$ is ample on
$\mathbf P^1_K$, and $V=(\omega)=(\omega_K)=1$ by
\eqref{eq:ch30-gauss-volume}.

\emph{Total mass.}  On the nontrivially valued side the computation of
\Cref{ex:ch30-power-map}, with $a=c$, gives
\[
 \MA_{\omega_K}(\Phi)=\delta_{x_c},
 \qquad \text{mass }1=V,
\]
where $x_c$ is the divisorial point with $v(z)=c$, $v(t)=1$.  On the
trivially valued side, \eqref{eq:bj-model-ma-atomic} applied to
$\cY_c$ gives the same coefficients, namely
$\bigl(\cL_c\cdot E_i\bigr)
 =\bigl((\rho^*L_{\mathbf A^1}+D)\cdot E_i\bigr)=0$ for $i<c$ and
$=1$ for $i=c$ --- the twist $D$ is already part of $\cL_c$ and must
not be added twice --- so
\[
 \MA^{\mathrm u}_\omega(\phi)=\delta_{c\,\ord_0},
 \qquad\text{mass }1=V,
 \qquad
 \MA_\omega(\phi)=\delta_{c\,\ord_0}
\]
(the normalized and unnormalized measures agree here because $V=1$).  By
\eqref{eq:ch30-gauss-divisorial}, $\sigma(c\,\ord_0)=x_c$, so
\[
 \sigma_*\MA^{\mathrm u}_\omega(\phi)=\MA_{\omega_K}(\Phi),
 \qquad
 r_*\MA_{\omega_K}(\Phi)=\MA^{\mathrm u}_\omega(\phi),
\]
which is \eqref{eq:ch30-gauss-measure-transfer}.

\emph{Energy.}  From \Cref{ex:ch30-power-map} with $d=1$, $a=c$,
\[
 E^{\mathrm{BFJ}}_{\omega_K}(\Phi)=-\frac c2 ;
\]
and on the trivially valued side, using
$\MA_\omega(0)=\delta_{v_{X,\triv}}$ and $\phi(v_{X,\triv})=0$,
\[
 E^{\mathrm{BJ}}_\omega(\phi)
 =\tfrac12\Bigl(\int\phi\,\MA_\omega(\phi)+\int\phi\,\MA_\omega(0)\Bigr)
 =\tfrac12\bigl(-c+0\bigr)=-\frac c2
 =\frac1V\,E^{\mathrm{BFJ}}_{\omega_K}(\Phi),
\]
in agreement with \eqref{eq:ch30-gauss-energy}.  For $c=3$ the four
numbers to remember are
\[
 \MA_{\omega_K}(\Phi)=\delta_{x_3},\quad
 \MA^{\mathrm u}_\omega(\phi)=\delta_{3\ord_0},\quad
 V=1,\quad
 E^{\mathrm{BFJ}}=E^{\mathrm{BJ}}=-\tfrac32 .
\]
\end{example}

The next example is the same computation with $V=2$, and it is included
because it separates the two energy normalizations and the two mass
normalizations, which coincide numerically when $V=1$.

\begin{example}[Two rays, and the factor $V$]
\label{ex:ch30-gauss-two-rays}
\index{Gauss extension!mass bookkeeping}
Keep $X=\mathbf P^1_k$ but take $L=\cO_{\mathbf P^1}(2)$, so that
$V=(\omega)=2$.  Fix $a,b\in\mathbf Z_{>0}$ and let
\[
 \mathfrak a:=\bigl(\mathfrak m_0+(t^a)\bigr)
              \cdot\bigl(\mathfrak m_\infty+(t^b)\bigr),
\]
a flag ideal whose Fubini--Study function is, by multiplicativity of
$\sigma(v)$ on ideals,
\[
 \phi(v)=-\min\{v(\mathfrak m_0),a\}-\min\{v(\mathfrak m_\infty),b\} .
\]
Thus $\phi(s\,\ord_0)=-\min\{s,a\}$, $\phi(s\,\ord_\infty)=-\min\{s,b\}$,
and $\phi\equiv0$ on every other ray and at $v_{X,\triv}$.

\emph{The nontrivially valued side.}  Write $s=v(z)$ for the tropical
coordinate on the two branches of $\mathbf P^{1,\An}_K$ through the Gauss
point, so that $s>0$ towards $z=0$ and $s<0$ towards $z=\infty$.
The corresponding model function is
\[
 \Phi(s)=-\min\{a,s_+\}-\min\{b,(-s)_+\},
\]
and $\sigma^*\Phi=\phi$ exactly as before.  With
$\MA_{\omega_K}(0)=2\delta_{x_G}$ --- the special fibre of
$\mathbf P^1_R$ carries $\cO(2)$ of degree $2$ --- and with outgoing
slopes $-1$ in each of the two directions at $x_G$, the graph Laplacian
\eqref{eq:curve-ma-laplacian} gives
\begin{equation}
\label{eq:ch30-gauss-two-rays-ma}
 \MA_{\omega_K}(\Phi)
 =\bigl(2-1-1\bigr)\delta_{x_G}+\delta_{x_a}+\delta_{x_{-b}}
 =\delta_{x_a}+\delta_{x_{-b}},
 \qquad\text{mass }2=V .
\end{equation}

\emph{The trivially valued side.}  The corresponding test configuration has
central fibre a chain of $a+b+1$ components, and the same computation ---
or \eqref{eq:ch30-gauss-measure-transfer} --- gives
\begin{equation}
\label{eq:ch30-gauss-two-rays-triv}
 \MA^{\mathrm u}_\omega(\phi)=\delta_{a\,\ord_0}+\delta_{b\,\ord_\infty},
 \qquad
 \MA_\omega(\phi)
 =\tfrac12\delta_{a\,\ord_0}+\tfrac12\delta_{b\,\ord_\infty},
\end{equation}
of masses $2$ and $1$ respectively.  Here the two normalizations of
\eqref{eq:bj-ma-normalizations} genuinely differ, and only the first
matches the BFJ measure \eqref{eq:ch30-gauss-two-rays-ma} under
$\sigma_*$.

\emph{Energies.}  On the nontrivially valued side,
\[
 E^{\mathrm{BFJ}}_{\omega_K}(\Phi)
 =\tfrac12\Bigl(\Phi(x_a)+\Phi(x_{-b})+2\,\Phi(x_G)\Bigr)
 =\tfrac12\bigl(-a-b+0\bigr)=-\frac{a+b}2 ,
\]
while the normalized trivially valued energy is
\[
 E^{\mathrm{BJ}}_\omega(\phi)
 =\frac1V\,E^{\mathrm{BFJ}}_{\omega_K}(\Phi)=-\frac{a+b}4 .
\]
For $a=2$, $b=1$ the table of numbers is
\begin{center}
\small
\begin{tabular}{p{.30\textwidth}|p{.28\textwidth}|p{.28\textwidth}}
\toprule
\textbf{Quantity}&\textbf{over $k$ (BJ)}&\textbf{over $K$ (BFJ)}\\
\midrule
class&$\omega=c_1(\cO(2))$&$\omega_K=c_1(\cO(2))$\\
volume $V$&$2$&$2$\\
$\MA$, unnormalized
 &$\delta_{2\ord_0}+\delta_{\ord_\infty}$, mass $2$
 &$\delta_{x_2}+\delta_{x_{-1}}$, mass $2$\\
$\MA$, probability
 &$\tfrac12\delta_{2\ord_0}+\tfrac12\delta_{\ord_\infty}$
 &not used in \Cref{part:bfj}\\
energy&$-\tfrac34$&$-\tfrac32$\\
\bottomrule
\end{tabular}
\end{center}
The two energies differ by the factor $V=2$, as
\eqref{eq:ch30-gauss-energy-normalization} predicts; the two unnormalized
masses agree, as \eqref{eq:ch30-gauss-mass} predicts.  A comparison theorem
which reported these two energies as equal, or these two masses as
differing, would be using inconsistent normalizations.

\emph{Centers and supports.}  The atoms $2\ord_0$ and $\ord_\infty$ have
centers $0$ and $\infty$ on $\mathbf P^1_k$ and zero support; their
Gauss images $x_2$ and $x_{-1}$ have the corresponding centers on
$\mathbf P^1_R$ and zero support.  The two rigid points $z=0$ and
$z=\infty$ of $\mathbf P^{1,\An}_K$ are the images of the two
proper-support endpoints and carry no mass.
\end{example}

\subsection{What does not transfer}
\label{sec:ch30-gauss-limits}

What the results above compare is the finitely generated, or model, locus.
They do not identify the full BFJ psh class over $k(\!(t)\!)$ with the
full BJ psh class over $k$; the latter contains homogeneous and
filtration-theoretic information developed in \cref{part:bj}.  Four
specific failures should be kept in mind.

\emph{Measures can have mass outside the Gauss image.}  By \Cref{thm:ch30-gauss-image} the image is the
$\mathbf G_m(k)$-invariant slice, and
\Cref{ex:ch30-gauss-image-strict} shows the fibres of $r$ are infinite.
A model metric can have atoms outside $\sigma(X^{\An})$; lack of an
equivariant presentation alone does not force this, since two presentations
can give the same metric.  For a concrete example on $\mathbf P^1_K$,
the standard-reference potential
$-\min\{2,\max\{0,v(z-1-t)\}\}$ has unit curvature mass at the point
$y$ of \Cref{ex:ch30-gauss-image-strict}, outside the image.
For every positive Radon measure $\nu$, one can restrict it to the closed
Gauss image and transport that restriction by the inverse homeomorphism.
That operation discards the mass outside the image; it is not an inverse
to pushforward on all measures.  On measures carried by the image,
$r_*$ is the inverse of $\sigma_*$.  On general measures it still exists
but can lose information in the fibres.

\emph{The full psh comparison requires a separate theorem.}
Every $\omega_K$-psh function $\varphi$ on $X_K^{\An}$ restricts along
$\sigma$ to a function on $X^{\An}$, and restriction is order preserving.
But $\sigma$ is not surjective, so restriction is not injective on
arbitrary functions; no injectivity statement is proved here on the psh
subclass.  The model dictionary proved here does not by itself establish a
restriction theorem for arbitrary singular potentials.  Its approximants
come from test configurations with the required equivariance; general
models over $R$ do not have that presentation automatically.
This gap in the displayed argument is not a counterexample to psh
restriction.  The full-class comparison requires a separate theorem,
whose scope must be specified before using it.

\emph{Homogeneity requires a separate comparison.}  The scaling action of
\Cref{sec:trivial-valuation-cone} is intrinsic to the trivially valued
space and, by \Cref{prop:ch30-gauss-scaling}, corresponds under $\sigma$
to changing the normalization of $v_K$.  The ordinary multiplication of all valuations does not define that
scaling action on the fixed-valued full space $X_K^{\An}$.
Thus the homogeneous theory and $b$-divisorial calculus of
\Cref{chap:bj-homogeneity-bdivisors} do not follow from the model
dictionary alone.  No impossibility claim about other notions of
homogeneity is intended.

\emph{The residue characteristic is constrained.}  The trivially valued
theory of \Cref{part:bj} is developed over an arbitrary algebraically
closed $k$, whereas the transfer above lands in
\eqref{eq:ch20-standing} only when $\operatorname{char}k=0$.  For
$\operatorname{char}k=p>0$ the space $X_K^{\An}$ still exists and
\Cref{thm:ch30-gauss-well-defined,thm:ch30-gauss-image,%
thm:ch30-gauss-model-dictionary} remain valid --- their proofs use no
resolution of singularities and no multiplier ideals.  The model intersection and atomic-measure comparison is likewise
characteristic-free for the stated normal models.  Analytic extensions
beyond the model level need the corresponding psh, operator and
regularization inputs; they are not supplied by these algebraic
calculations alone.  This is the same boundary as in
\Cref{thm:ch23-gjkm-positive-characteristic} and
\Cref{thm:gjkm-positive-characteristic-cy}.

\section{Scope: where each hypothesis is used}
\label{sec:ch30-scope}

\subsection{Hypothesis by hypothesis}
\label{sec:ch30-scope-table}

\begin{center}
\small
\begin{tabular}{p{.24\textwidth}|p{.34\textwidth}|p{.34\textwidth}}
\toprule
\textbf{Hypothesis}&\textbf{Where it is used}
 &\textbf{What fails without it}\\
\midrule
$K$ discretely valued
 &everywhere models are quantified over:
  \Cref{lem:ch30-base-change-model},
  \Cref{thm:ch30-base-change-psh}
 &for a non-discrete value group use a valuation-ring or admissible
  formal model framework; \Cref{sec:ch30-bc-beyond}\\
$\operatorname{char}k=0$
 &the standing model/resolution and analytic framework,
  including \Cref{thm:bfj-decreasing-regularization}, hence in
  \Cref{thm:ch30-pullback-psh},
  \Cref{thm:ch30-finite-bounded-projection} and their consequences; and in
  \eqref{eq:ch30-efn}
 &the model statements
  (\Cref{lem:ch30-pullback-model-calculus},
  \Cref{thm:finite-ma-projection} for model potentials,
  \Cref{thm:base-change-ma} for model metrics,
  \Cref{thm:ch30-gauss-model-dictionary}) survive; the extensions to
  arbitrary singular, bounded and finite-energy potentials do not, by this
  proof\\
$X$ smooth
 &the standing BFJ analytic framework; the displayed model projection
  formulae use normal models and the supplied morphisms
 &those analytic extensions are not provided by the stated BFJ
  hypotheses; the model intersection identities still make sense\\
$\{\theta\}$ ample
 &the BFJ theory on $X$, hence
  \Cref{sec:ch30-finite,sec:ch30-base-change}
 &the singular and finite-energy calculus used on $X$ is unavailable\\
$f^*\{\theta\}$ ample
 &\Cref{thm:ch30-pullback-psh} and every finite-cover consequence
 &for a birational pullback it is generally false, which is why
  \Cref{sec:ch30-birational} retains model and continuous projection
  formulae only\\
$X$ geometrically integral
 &\Cref{lem:ch30-base-change-model}\,\cref{itm:ch30-bcmod:finite} and every
  degree computation after it
 &a sum over connected components of $X_{K'}$;
  \Cref{rem:ch30-not-geometrically-integral}\\
$f$ surjective
 &\Cref{thm:ch30-pullback-psh}\,\cref{itm:ch30-pbpsh:surj} and the
  finite-cover sections, through
  \eqref{eq:ch30-pullback-hypothesis} and equality of suprema
 &$f^*\varphi\equiv-\infty$ is possible;
  \Cref{ex:ch30-polar-pullback-fails}\\
$K'/K$ Galois
 &\Cref{sec:ch30-descent} only
 &pass to the Galois closure, harmless in characteristic zero\\
$K'/K$ unramified
 &\Cref{thm:ch30-model-descent} only
 &the semilinear action is not a descent datum;
  \Cref{sec:ch30-descent-limits}\\
$k$ infinite
 &\Cref{thm:ch30-gauss-image}, through the Vandermonde argument
 &the characterization of the Gauss image by
  $\mathbf G_m(k)$-invariance fails\\
\bottomrule
\end{tabular}
\end{center}

\subsection{What later work removes, and what it does not}
\label{sec:ch30-scope-later}

The dependence on $\operatorname{char}k=0$ is inherited from the
standing model and analytic inputs, including resolution/cofinality,
decreasing regularization and the bounded and finite-energy operator.
The displayed intersection projection identities are
characteristic-free once the required normal models and morphisms
are supplied.  Where
regularization is available by other means --- the test-ideal methods of
\cite{GJKM19} in the situations covered by
\Cref{thm:ch23-gjkm-positive-characteristic}, or the mixed-characteristic
results of \cite{FGK25} --- the corresponding extensions of
\Cref{thm:ch30-pullback-psh} and
\Cref{thm:ch30-finite-bounded-projection} follow when the corresponding psh stability, model pullback and
operator convergence package is available on both sides.  Replace the
regularization input and retain the ample-class signature wherever it
is used.  Regularization alone is not a substitute for these operator
and model-theoretic hypotheses.  The later envelope and regularization
conditions are recorded in
\Cref{chap:bfj-compactness,chap:bfj-envelopes,chap:bfj-continuity}.

The dependence on discreteness is different in kind.  It is not a
convenience: the definition of a model function, the normalization
$\varphi_{\cX_s}=1$ of \eqref{eq:ch20-sign}, and the multiplicity
$b_E=\ord_E(\varpi)$ that appears in every measure formula all presuppose
it.  Over a non-discretely valued field the correct framework replaces
$\mathcal D(X)$ by piecewise linear functions in the sense of
\Cref{thm:gm-metric-comparison} and models by admissible formal models; the metric and measure constructions in that framework require
their own hypotheses.  \cite{GM19} supplies the formal/model metric
comparison and \cite{BGM21} develops the volume theory; neither is
being quoted here as a replacement proof of every singular projection
formula.

\section{Dictionary and functoriality ledger}
\label{sec:ch30-ledger}

\subsection{The ledger of safe statements and boundaries}
\label{sec:ch30-ledger-list}

\begin{description}[style=nextline,leftmargin=0pt,labelindent=0pt,
  labelsep=0pt,itemsep=.8\baselineskip,font=\normalfont\bfseries]
\item[Finite pullback]
\emph{Safe statement.} Psh metrics pull back, and energy scales by $d$.
\par\emph{Boundary.} An unrestricted pushforward identity alone does
not select the measure upstairs; additional degree or invariance data can.

\item[Proper pushforward]
\emph{Safe statement.} Radon measures and CLD currents push forward.
\par\emph{Boundary.} Pushforward does not reconstruct a metric upstairs.

\item[Birational pullback]
\emph{Safe statement.} Model and continuous metrics satisfy the projection
formula.
\par\emph{Boundary.} The pulled-back ample class is usually only nef and big.

\item[Finite field extension]
\emph{Safe statement.} Model metrics pull back, with pushforward given by
\eqref{eq:base-change-ma}.
\par\emph{Boundary.} Psh functions and measures upstairs need not descend.

\item[Galois descent]
\emph{Safe statement.} Uniqueness forces invariance of a normalized solution.
\par\emph{Boundary.} Positivity descent is an additional input;
invariant measures have the unique lift of
\Cref{cor:ch30-invariant-functions}.

\item[Gauss extension]
\emph{Safe statement.} Model functions, test configurations, and
intersections match on the Gauss image.
\par\emph{Boundary.} The full singular theories are not thereby identified.

\item[Current pullback]
\emph{Safe statement.} Pullback exists in the stated smooth or formal
situations.
\par\emph{Boundary.} There is no general pullback of CLD currents.
\end{description}

\subsection{Objects and their two descriptions}
\label{sec:ch30-dict-objects}

\begin{center}
\small
\begin{tabular}{p{.28\textwidth}|p{.30\textwidth}|p{.34\textwidth}}
\toprule
\textbf{Geometric datum}&\textbf{Pluripotential-theoretic effect}
 &\textbf{Established at}\\
\midrule
degree $d$ of a generically finite morphism
 &every mass, mixed measure and energy multiplied by $d$
 &\Cref{thm:finite-ma-projection},
  \Cref{thm:ch30-finite-energy-transfer}\\
contraction of a curve by a birational $f$
 &$f^*\{\theta\}$ nef and big, not ample
 &\eqref{eq:birational-top-intersection},
  \Cref{sec:ch30-birational-setting}\\
identification $K(Y)=K(X)$
 &$f^{\An}$ bijective on zero-support points, hence on the atoms of
  model Monge--Amp\`ere measures
 &\Cref{lem:ch30-birational-bijection},
  \Cref{thm:ch30-birational-identification}\\
ramification index $e$
 &value group $\tfrac1e\mathbf Z$; multiplicities
  $b_{E'}=b_Em_{E'}/e$
 &\eqref{eq:ch30-uniformizer-value},
  \eqref{eq:ch30-multiplicity-relation}\\
residue degree $f$
 &degrees on the special fibre divided by $f$
 &\eqref{eq:ch30-degree-transfer}\\
$ef=[K':K]$ in the present defectless setting
 &no factor $[K':K]$ in \eqref{eq:base-change-ma}
 &\eqref{eq:ch30-coefficient-cancellation},
  \Cref{rem:ch30-cancellation}\\
renormalizing $v_{K'}(\varpi')=1$
 &weights, lengths, energies multiplied by $e$; measures unchanged
 &\eqref{eq:ch30-ramified-renormalization}\\
$K'/K$ Galois with group $G$
 &$X^{\An}=X_{K'}^{\An}/G$
 &\Cref{thm:ch30-galois-quotient}\\
test configuration over $\mathbf A^1_k$
 &model over $k[\![t]\!]$ with the same components and multiplicities
 &\Cref{thm:ch30-gauss-model-dictionary}\,\cref{itm:ch30-gmd:model}\\
$\mathbf G_m(k)$-invariance
 &membership in the Gauss image
 &\Cref{thm:ch30-gauss-image}\\
scaling $v\mapsto rv$ on $X^{\An}$
 &change of normalization of $v_K$
 &\Cref{prop:ch30-gauss-scaling}\\
probability versus mass-$V$ normalization
 &factor $V$ between the two energies
 &\eqref{eq:ch30-gauss-energy-normalization}\\
\bottomrule
\end{tabular}
\end{center}

\subsection{Reading the dictionary backwards}
\label{sec:ch30-dict-backward}

Five entries are one-way, and it is worth saying which.

\emph{From a pushed-forward measure to the measure upstairs} there is no
inverse without extra constraints; see
\Cref{prop:ch30-pushforward-not-inverse,ex:ch30-two-measures}.
There is an inverse on invariant measures for a finite Galois quotient
by \Cref{cor:ch30-invariant-functions}, and on measures carried by the
Gauss image via its inverse section.

\emph{From a psh function upstairs to one downstairs} there is no arrow in
general.  Along a morphism it fails for the trivial reason that a potential
need not be constant on fibres
(\Cref{rem:ch30-pullback-not-onto}); along a field extension it fails
because the class genuinely grows
(\Cref{rem:ch30-psh-grows}); what can be recovered, and under what
hypotheses, is \Cref{sec:ch30-descent}.

\emph{From the total mass to the geometry} there is no arrow: the two
examples \Cref{ex:ch30-power-map} and
\Cref{ex:ch30-conic-base-change} both begin with mass-$V$ measures and
end with mass-$dV$ and mass-$V$ measures respectively, and the numbers
alone do not say which operation was performed.  What distinguishes them is
\emph{not} the degree of the induced extension of function fields: that
degree is $d$ in the first case and $[K':K]$ in the second, both greater
than one, and only the first appears in the mass.  It is whether the ground
field --- and with it the residue field over which special-fibre degrees are
taken --- has changed; see \Cref{rem:ch30-two-inputs}.

\emph{From an equality of intersection numbers to an equality of measures}
there is no arrow.  \Cref{prop:gauss-model-compatibility} states the
coefficient comparison; the measure statement
\eqref{eq:ch30-gauss-measure-transfer} additionally requires the
identification of the atoms
\eqref{eq:ch30-gauss-divisorial}, which is where
\Cref{thm:ch30-gauss-image} is used.

\emph{From the Gauss image to the whole space} there is no arrow: $r$ is
defined everywhere, but $\sigma\circ r\neq\mathrm{id}$, and
\Cref{ex:ch30-gauss-image-strict} shows the fibres of $r$ are infinite.

\subsection{Numbers to remember}
\label{sec:ch30-dict-numbers}

\begin{center}
\small
\begin{tabular}{p{.40\textwidth}|p{.22\textwidth}|p{.30\textwidth}}
\toprule
\textbf{Quantity}&\textbf{Value}&\textbf{Where}\\
\midrule
$\MA$ of the tent on $\mathbf P^1$, $d=3$, $a=6$
 &$\delta_{x_6}$ below, $3\delta_{x_2}$ above
 &\eqref{eq:ch30-power-map-ma-down},
  \eqref{eq:ch30-power-map-ma-up}\\
its two energies
 &$-3$ and $-9$
 &\Cref{ex:ch30-power-map}\\
mass under a finite morphism
 &$V\mapsto dV$
 &\eqref{eq:finite-diagonal-projection}\\
$\MA$ split on the blown-up plane
 &$(1-c^2,\,c^2)$, $-1\leq c\leq0$
 &\eqref{eq:ch30-surface-ma-down}\\
its energy
 &$c^3/3$; $-\tfrac13$ at $c=-1$
 &\eqref{eq:ch30-surface-energy}\\
mass under a birational morphism
 &$V\mapsto V$
 &\Cref{thm:ch30-birational-identification}\\
conic: multiplicity at $v_{1/2}$
 &$b_E=2$ over $K$, $b'=1$ over $K'$
 &\Cref{ex:ch30-conic-base-change}\\
its ramification numbers
 &$e=2$, $f=1$, $m=1$, $d_{E''}=2$
 &\Cref{ex:ch30-conic-base-change}\\
its measure and energy
 &$2\delta_{v_{1/2}}$ and $-\tfrac12$, both sides
 &\eqref{eq:ch30-conic-ma-down},
  \eqref{eq:ch30-conic-energy-up}\\
the same after renormalizing $v(\varpi')=1$
 &energy $-1=e\cdot(-\tfrac12)$
 &\eqref{eq:ch30-ramified-renormalization}\\
Gauss extension of $\max\{-3,-v(\mathfrak m_0)\}$
 &$\MA=\delta_{3\ord_0}\mapsto\delta_{x_3}$, mass $1$
 &\Cref{ex:ch30-gauss-normal-cone}\\
its energy
 &$-\tfrac32$, the same on both sides since $V=1$
 &\Cref{ex:ch30-gauss-normal-cone}\\
two-ray example, $V=2$, $a=2$, $b=1$
 &masses $2$ and $2$; energies $-\tfrac34$ and $-\tfrac32$
 &\Cref{ex:ch30-gauss-two-rays}\\
\bottomrule
\end{tabular}
\end{center}

\section*{Source notes}

The model projection formulas are consequences of intersection theory and
the Chambert--Loir construction \cite{CL06}.  Approximation of semipositive
metrics is treated in \cite{GM19}; norm and scalar-extension methods are
developed further in \cite{BE21}.  The finite pullback statement for BFJ psh
functions follows from the skeletal characterization in \cite{BFJ16}.  The
Gauss construction is used systematically in the trivially valued theory of
\cite{BJ22}.  No pullback of arbitrary currents or preferred inverse to
pushforward under an arbitrary continuous map is asserted.
The finite Galois invariant measure lift and the Gauss-image
correspondence are explicit exceptions to unrestricted nonuniqueness.

\paragraph{Provenance, statement by statement.}
The framework in which all of this chapter takes place --- closed model
forms, model functions, the operator $\ddc$, and semipositivity as
relative nefness --- is \cite[\S\S2 and 4--5]{BFJ16}, and the projection
formula \Cref{thm:model-projection-formula} on which
\Cref{sec:ch30-finite,sec:ch30-birational} rest is classical intersection
theory in the form of \cite[\S1.4 and \S2.3]{Ful}.
\Cref{prop:finite-pullback-psh} and
\Cref{prop:birational-model-projection} are the model-level statements of
\cite[\S2]{BFJ16}; the extension to bounded and finite-energy potentials
(\Cref{thm:ch30-finite-bounded-projection},
\Cref{thm:ch30-finite-energy-transfer}) uses only the monotone convergence
property \cite[Theorem~3.1]{BFJ15} and the mixed products of
\cite[\S6]{BFJ15}, and the proofs given here are the standard
approximation arguments made explicit.
\Cref{thm:ch30-pullback-psh} is the statement that
\Cref{chap:bfj-psh} deferred: its proof is nothing but decreasing
regularization \cite[Theorem~8.3]{BFJ16} combined with stability under
decreasing nets, after requiring the pulled-back class to remain ample.
Within that framework, \eqref{eq:ch30-pullback-hypothesis} is necessary to
exclude the identically $-\infty$ limit, as
\Cref{ex:ch30-polar-pullback-fails} shows.

The base-change material of \Cref{sec:ch30-base-change} is standard, but the
bookkeeping is rarely displayed.  Berkovich's description of the fibres of
$\pi$ and the surjectivity of base change are \cite[\S1.4 and
\S3.4]{Berk12}; under the separability and residue-characteristic-zero
hypotheses above, the fundamental equality $ef=[K':K]$ is
\cite[Ch.~VI]{ZariskiSamuelII}.  The cancellation
\eqref{eq:ch30-coefficient-cancellation}, which is the whole of
\Cref{thm:base-change-ma}, is implicit in the invariance of Chambert--Loir
measures under ground-field extension used throughout \cite{CL06} and
\cite{GM19}; we have preferred to prove it from
\eqref{eq:ch30-ramification-sum} and
\eqref{eq:ch30-degree-transfer} because those two identities, and not the
conclusion, are what a reader must check when changing conventions.  The
analogous statement for tropical calibrations is
\Cref{cor:trop-finite-field-extension}, taken from
\cite[\S13--14]{CLD25}, and the factor $[F:K]$ appearing there rather than
here is explained in \Cref{rem:ch30-two-inputs}.  Behaviour of models under
a ramified base change, and the semistable reduction that
\Cref{ex:ch30-conic-base-change} illustrates, belong to
\cite{deJong1996}; the formal-model framework needed when the value group
ceases to be discrete is \cite{BoschLutkebohmert1993} and
\cite{RaynaudGruson1971}.  Non-Archimedean volumes over general complete
fields, where the same projection formulae are used systematically, are
\cite{BGM21} and \cite{BGGJKM20}.

\Cref{sec:ch30-descent} separates its topological, positivity and
solution inputs.
\Cref{thm:ch30-galois-quotient} is classical --- it is the analytic
counterpart of the transitivity of a finite group action on the primes above
a fixed prime, \cite[Tag 0BRI]{stacks-project}, recorded in
\Cref{app:stacks} --- and the proof given is self-contained.
\Cref{thm:ch30-model-descent} uses faithfully flat descent in the form of
\cite[Exp.~VIII, Theorem~1.1, Corollary~5.8 and
Proposition~7.8]{SGA1}, and is proved here for unramified extensions
because that is where a semilinear action is a descent datum;
\Cref{thm:ch30-continuous-descent} is therefore conditional, and
\Cref{sec:ch30-descent-limits} says exactly on what.  No claim for the
ramified model descent property is made.  The local characterization of
semipositivity in \cite{GM19} provides a different framework for studying
descent along arbitrary complete extensions; that route is not taken here.

The Gauss extension of \Cref{sec:ch30-gauss} is
\cite[\S1.3 and Remark~1.9]{BJ22}, where \eqref{eq:gauss-extension} and
\eqref{eq:gauss-section} appear, and it is used in the same way in
\cite{BHJ17} to compare test configurations with valuations.  The
characterization of the image by $\mathbf G_m(k)$-invariance,
\Cref{thm:ch30-gauss-image}, is stated in
\Cref{chap:trivial-analytification} after
\eqref{eq:triv-gauss-section} following \cite[\S1]{BJ22} and
\cite[\S4.1, Lemma~4.2]{BHJ17}; the Vandermonde proof given here is elementary and is
included because it is the step that makes
\eqref{eq:ch30-gauss-divisorial} --- and hence the measure transfer --- a
theorem rather than a definition.  The intersection-number comparison
\eqref{eq:ch30-gauss-intersection} is the standard compatibility between
the compactified intersection numbers of
\cite[\S6.6, Definition~6.11 and Remark~6.12]{BHJ17}
and \cite[\S7]{BJ22} with their vertical-term expansions on the
formal completion.  The normalization discrepancy
\eqref{eq:ch30-gauss-energy-normalization} is the one recorded in
\eqref{eq:bj-ma-normalizations} and \eqref{eq:appJ-energy}; it is not a
disagreement between sources but a deliberate difference of convention
between the two parts of this book, and
\Cref{ex:ch30-gauss-two-rays} exists to make it numerical.

\paragraph{What is organized differently here.}
Three choices depart from the usual arrangement.  First, the three
transport situations of \Cref{sec:ch30-three-kinds} are kept apart
throughout, and in particular the extra Gauss section is never conflated
with its underlying base-change projection; the ledger of
\Cref{sec:ch30-ledger-list} exists to make the difference between the
degree $d$, the pair $(e,f)$, and the absence of any such number in the
Gauss case impossible to overlook.  Second, every projection formula is
proved three times --- for model, for bounded, and for finite-energy
potentials --- rather than once with an unspecified passage to the limit,
because the three arguments use different convergence inputs.
The model intersection step itself is characteristic-free once its
models are supplied; the later steps retain their full analytic
hypotheses.  Third, descent is separated into the
three questions of \Cref{sec:ch30-descent-questions}; the literature tends
to make the quotient in (Q1) and the uniqueness step in (Q3) explicit while
leaving the positivity step (Q2) implicit.  The conditional form of
\Cref{thm:ch30-continuous-descent} records exactly that missing step.

\paragraph{Terminology and scope.}
Function and metric pullback means composition with the stated map.
The invariant measure lift of \eqref{eq:ch30-invariant-measure-lift}
and the measure correspondence on the Gauss image have their own
definitions; neither is an inverse to unrestricted pushforward under
an arbitrary continuous map.  Pullback of general currents remains
outside this scope.  The word \emph{descent} is used only in the sense of
\Cref{sec:ch30-descent-questions}.  \emph{Model function} always means an
element of $\mathcal D(X)_{\mathbf R}$ as defined in
\Cref{chap:bfj-model-forms}, over the ground field currently in force;
when two ground fields are present, the field is named.  Finally, no
statement of this chapter asserts an equality between the trivially valued
class $\PSH(X,\omega)$ of \Cref{part:bj} and the BFJ class
$\PSH(X_K,\omega_K)$; what is asserted is the model dictionary
\Cref{thm:ch30-gauss-model-dictionary} and its measure-theoretic corollary
\Cref{thm:ch30-gauss-measures}, and the comparison of the two full theories
is deferred to \Cref{part:comparisons}.
